
\Data\ID{1003084910}
\Updated{2020-07-19 11:07:09}
\Level{A}
\SubArea{Geometric Sequences}
\LANGen{
\Question{
We were given a geometric sequence \( \frac12\text{, }\ \frac14\text{, }\ \dots \). 
What is the formula for the \( n \)th element of the sequence?}
{\( a_n=\frac1{2^n}\text{, }\ n\in\mathbb{N} \)}{\( a_n=\frac1{2^{n+1}}\text{, }\ n\in\mathbb{N} \)}{\( a_n=\frac1{2^{n-1}}\text{, }\ n\in\mathbb{N} \)}{\( a_n=\frac1{2^{2n}}\text{, }\ n\in\mathbb{N} \)}{}{}}
\LANGpl{
\Question{
Dany jest ciąg geometryczny \( \frac12\text{, }\ \frac14\text{, }\ \dots \). 
Jaki jest wzór  \( n \)-tego wyrazu tego ciągu?}
{\( a_n=\frac1{2^n}\text{, }\ n\in\mathbb{N} \)}{\( a_n=\frac1{2^{n+1}}\text{, }\ n\in\mathbb{N} \)}{\( a_n=\frac1{2^{n-1}}\text{, }\ n\in\mathbb{N} \)}{\( a_n=\frac1{2^{2n}}\text{, }\ n\in\mathbb{N} \)}{}{}}
\LANGcs{
\Question{
Je dána geometrická posloupnost \( \frac12\text{, }\ \frac14\text{, }\ \dots \). 
Vzorec pro \( n \)-tý člen této posloupnosti je:}
{\( a_n=\frac1{2^n}\text{, }\ n\in\mathbb{N} \)}{\( a_n=\frac1{2^{n+1}}\text{, }\ n\in\mathbb{N} \)}{\( a_n=\frac1{2^{n-1}}\text{, }\ n\in\mathbb{N} \)}{\( a_n=\frac1{2^{2n}}\text{, }\ n\in\mathbb{N} \)}{}{}}
\LANGsk{
\Question{
Je daná geometrická postupnosť \( \frac12\text{, }\ \frac14\text{, }\ \dots \). 
Vzorec pre \( n \)-tý člen tejto postupnosti je:}
{\( a_n=\frac1{2^n}\text{, }\ n\in\mathbb{N} \)}{\( a_n=\frac1{2^{n+1}}\text{, }\ n\in\mathbb{N} \)}{\( a_n=\frac1{2^{n-1}}\text{, }\ n\in\mathbb{N} \)}{\( a_n=\frac1{2^{2n}}\text{, }\ n\in\mathbb{N} \)}{}{}}
\LANGes{
\Question{
Dada la progresión geométrica \( \frac12\text{, }\ \frac14\text{, }\ \dots \). 
Halla el término general de la sucesión.}
{\( a_n=\frac1{2^n}\text{, }\ n\in\mathbb{N} \)}{\( a_n=\frac1{2^{n+1}}\text{, }\ n\in\mathbb{N} \)}{\( a_n=\frac1{2^{n-1}}\text{, }\ n\in\mathbb{N} \)}{\( a_n=\frac1{2^{2n}}\text{, }\ n\in\mathbb{N} \)}{}{}}
\End

\Data\ID{1003107308}
\Updated{2023-03-27 08:03:58}
\Level{A}
\SubArea{Geometric Sequences}
\LANGen{
\Question{
The first five terms of a geometrical sequence are: \( -2,\ 1,\,-\frac12,\ \frac14,\,-\frac18 \).
Find the recursive formula of this  sequence.}
{\( a_1=-2\,;\ a_{n+1}=a_n\cdot\left(-\frac12\right),\ n\in\mathbb{N} \)}{\( a_1=-2\,;\ a_{n+1}=a_n\cdot\left(-\frac14\right),\ n\in\mathbb{N} \)}{\( a_1=-2\,;\ a_{n+1}=a_n\cdot\frac12,\ n\in\mathbb{N} \)}{\( a_1=-2\,;\ a_{n+1}=a_n\cdot\left(-\frac12\right)^n,\ n\in\mathbb{N} \)}{}{}}
\LANGpl{
\Question{
Pięć pierwszych wyrazów ciągu arytmetycznego to: \( -2,\ 1,\,-\frac12,\ \frac14,\,-\frac18 \).
Określ rekurencyjny wzór tego ciągu.}
{\( a_1=-2\,;\ a_{n+1}=a_n\cdot\left(-\frac12\right),\ n\in\mathbb{N} \)}{\( a_1=-2\,;\ a_{n+1}=a_n\cdot\left(-\frac14\right),\ n\in\mathbb{N} \)}{\( a_1=-2\,;\ a_{n+1}=a_n\cdot\frac12,\ n\in\mathbb{N} \)}{\( a_1=-2\,;\ a_{n+1}=a_n\cdot\left(-\frac12\right)^n,\ n\in\mathbb{N} \)}{}{}}
\LANGcs{
\Question{
Je dáno prvních pět členů geometrické posloupnosti: \( -2,\ 1,\,-\frac12,\ \frac14,\,-\frac18 \).
Rekurentní vyjádření této posloupnosti je:}
{\( a_1=-2\,;\ a_{n+1}=a_n\cdot\left(-\frac12\right),\ n\in\mathbb{N} \)}{\( a_1=-2\,;\ a_{n+1}=a_n\cdot\left(-\frac14\right),\ n\in\mathbb{N} \)}{\( a_1=-2\,;\ a_{n+1}=a_n\cdot\frac12,\ n\in\mathbb{N} \)}{\( a_1=-2\,;\ a_{n+1}=a_n\cdot\left(-\frac12\right)^n,\ n\in\mathbb{N} \)}{}{}}
\LANGsk{
\Question{
Je daných prvých päť členov geometrickej postupnosti: \( -2,\ 1,\,-\frac12,\ \frac14,\,-\frac18 \).
Rekurentné vyjadrenie tejto postupnosti je:}
{\( a_1=-2\,;\ a_{n+1}=a_n\cdot\left(-\frac12\right),\ n\in\mathbb{N} \)}{\( a_1=-2\,;\ a_{n+1}=a_n\cdot\left(-\frac14\right),\ n\in\mathbb{N} \)}{\( a_1=-2\,;\ a_{n+1}=a_n\cdot\frac12,\ n\in\mathbb{N} \)}{\( a_1=-2\,;\ a_{n+1}=a_n\cdot\left(-\frac12\right)^n,\ n\in\mathbb{N} \)}{}{}}
\LANGes{
\Question{
Dados los cinco primeros términos de una progresión geométrica: \( -2,\ 1,\,-\frac12,\ \frac14,\,-\frac18 \).
Halla la fórmula recursiva de la sucesión.}
{\( a_1=-2\,;\ a_{n+1}=a_n\cdot\left(-\frac12\right),\ n\in\mathbb{N} \)}{\( a_1=-2\,;\ a_{n+1}=a_n\cdot\left(-\frac14\right),\ n\in\mathbb{N} \)}{\( a_1=-2\,;\ a_{n+1}=a_n\cdot\frac12,\ n\in\mathbb{N} \)}{\( a_1=-2\,;\ a_{n+1}=a_n\cdot\left(-\frac12\right)^n,\ n\in\mathbb{N} \)}{}{}}
\End

\Data\ID{1003112801}
\Updated{2021-12-26 08:12:51}
\Level{A}
\SubArea{Geometric Sequences}
\LANGen{
\Question{
The third term of a geometric sequence is \( 100 \) and the common ratio is \( -5 \). Find the first term of this sequence.}
{\( 4 \)}{\( -4 \)}{\( 20 \)}{\( -20 \)}{\( 500 \)}{}}
\LANGpl{
\Question{
Trzeci wyraz ciągu geometrycznego wynosi \( 100 \), zaś wspólny współczynnik to \( -5 \). Wyznacz pierwszy wyraz tego ciągu.}
{\( 4 \)}{\( -4 \)}{\( 20 \)}{\( -20 \)}{\( 500 \)}{}}
\LANGcs{
\Question{
Třetí člen geometrické posloupnosti je roven \( 100 \) a kvocient je \( -5 \). Určete první člen této posloupnosti.}
{\( 4 \)}{\( -4 \)}{\( 20 \)}{\( -20 \)}{\( 500 \)}{}}
\LANGsk{
\Question{
Tretí člen geometrickej postupnosti je rovný \( 100 \) a kvocient je \( -5 \). Určte prvý člen tejto postupnosti.}
{\( 4 \)}{\( -4 \)}{\( 20 \)}{\( -20 \)}{\( 500 \)}{}}
\LANGes{
\Question{
El tercer término de una progresión geométrica es \( 100 \) y su razón es \( -5 \). Averigua el primer término de la progresión.}
{\( 4 \)}{\( -4 \)}{\( 20 \)}{\( -20 \)}{\( 500 \)}{}}
\End

\Data\ID{1003112802}
\Updated{2022-01-06 01:01:04}
\Level{A}
\SubArea{Geometric Sequences}
\LANGen{
\Question{
The first term of a geometric sequence is \( -36 \) and the fourth term is \( -\frac{32}3 \). Find the common ratio.}
{\( \frac23 \)}{\( -\frac23 \)}{\( \frac32 \)}{\( -\frac32 \)}{\( \frac13 \)}{}}
\LANGpl{
\Question{
Pierwszy wyraz ciągu geometrycznego wynosi \( -36 \), a czwarty wyraz ciągu to \( -\frac{32}3 \). Wyznacz wspólny współczynnik.}
{\( \frac23 \)}{\( -\frac23 \)}{\( \frac32 \)}{\( -\frac32 \)}{\( \frac13 \)}{}}
\LANGcs{
\Question{
První člen geometrické posloupnosti je roven \( -36 \) a čtvrtý člen je \( -\frac{32}3 \). Najděte kvocient této posloupnosti.}
{\( \frac23 \)}{\( -\frac23 \)}{\( \frac32 \)}{\( -\frac32 \)}{\( \frac13 \)}{}}
\LANGsk{
\Question{
Prvý člen geometrickej postupnosti je rovný \( -36 \) a štvrtý člen je \( -\frac{32}3 \). Nájdite kvocient tejto postupnosti.}
{\( \frac23 \)}{\( -\frac23 \)}{\( \frac32 \)}{\( -\frac32 \)}{\( \frac13 \)}{}}
\LANGes{
\Question{
El primer término de una sucesión geométrica es igual a \( -36 \) y su cuarto término es \( -\frac{32}3 \). Encuentra la razón de la sucesión.}
{\( \frac23 \)}{\( -\frac23 \)}{\( \frac32 \)}{\( -\frac32 \)}{\( \frac13 \)}{}}
\End

\Data\ID{1003112803}
\Updated{2022-04-23 05:04:14}
\Level{A}
\SubArea{Geometric Sequences}
\LANGen{
\Question{
The second term of a geometric sequence is \( 24 \) and the fifth term is \( 3 \). Choose the correct formula to find the third term of this sequence.}
{\( a_3=24\cdot\sqrt[3]{\frac3{24}} \)}{\( a_3=24\cdot\sqrt[3]{\frac{24}3} \)}{\( a_3=3\cdot\sqrt[3]{\frac3{24}} \)}{\( a_3=3\cdot\sqrt[3]{\frac{24}3} \)}{\( a_3=8\cdot\sqrt[3]{\frac3{24}} \)}{}}
\LANGpl{
\Question{
Drugi wyraz ciągu geometrycznego wynosi \( 24 \), piąty  \( 3 \). Wybierz poprawny wzór, aby wyznaczyć trzeci wyraz ciągu.}
{\( a_3=24\cdot\sqrt[3]{\frac3{24}} \)}{\( a_3=24\cdot\sqrt[3]{\frac{24}3} \)}{\( a_3=3\cdot\sqrt[3]{\frac3{24}} \)}{\( a_3=3\cdot\sqrt[3]{\frac{24}3} \)}{\( a_3=8\cdot\sqrt[3]{\frac3{24}} \)}{}}
\LANGcs{
\Question{
Druhý člen geometrické posloupnosti je roven \( 24 \) a pátý člen je \( 3 \). Zvolte správný postup pro výpočet třetího členu této posloupnosti.}
{\( a_3=24\cdot\sqrt[3]{\frac3{24}} \)}{\( a_3=24\cdot\sqrt[3]{\frac{24}3} \)}{\( a_3=3\cdot\sqrt[3]{\frac3{24}} \)}{\( a_3=3\cdot\sqrt[3]{\frac{24}3} \)}{\( a_3=8\cdot\sqrt[3]{\frac3{24}} \)}{}}
\LANGsk{
\Question{
Druhý člen geometrickej postupnosti je rovný \( 24 \) a piaty člen je \( 3 \). Zvoľte správny postup pre výpočet tretieho člena tejto postupnosti.}
{\( a_3=24\cdot\sqrt[3]{\frac3{24}} \)}{\( a_3=24\cdot\sqrt[3]{\frac{24}3} \)}{\( a_3=3\cdot\sqrt[3]{\frac3{24}} \)}{\( a_3=3\cdot\sqrt[3]{\frac{24}3} \)}{\( a_3=8\cdot\sqrt[3]{\frac3{24}} \)}{}}
\LANGes{
\Question{
El segundo término de una progresión geométrica es igual a \( 24 \) y su quinto término es \( 3 \). Elije el procedimiento correcto para calcular el tercer término de la progresión.}
{\( a_3=24\cdot\sqrt[3]{\frac3{24}} \)}{\( a_3=24\cdot\sqrt[3]{\frac{24}3} \)}{\( a_3=3\cdot\sqrt[3]{\frac3{24}} \)}{\( a_3=3\cdot\sqrt[3]{\frac{24}3} \)}{\( a_3=8\cdot\sqrt[3]{\frac3{24}} \)}{}}
\End

\Data\ID{1003112805}
\Updated{2022-04-23 05:04:14}
\Level{A}
\SubArea{Geometric Sequences}
\LANGen{
\Question{
The \(n\)th term of a geometric sequence is \(\frac52 \), the common ratio is \( \frac12 \) and the fourth term is \( 20 \). Find \( n \).}
{\( 7 \)}{\( 8 \)}{\( 6 \)}{\( 10 \)}{\( 5 \)}{}}
\LANGpl{
\Question{
\(n\)-ty wyraz ciągu geometrycznego  to \(\frac52 \), wspólny współczynnik wynosi \( \frac12 \), a czwarty wyraz ciągu to  \( 20 \). Wyznacz \( n \).}
{\( 7 \)}{\( 8 \)}{\( 6 \)}{\( 10 \)}{\( 5 \)}{}}
\LANGcs{
\Question{
\(n\)-tý člen geometrické posloupnosti je roven \(\frac52 \), kvocient je \( \frac12 \) a čtvrtý člen je \( 20 \). Určete \( n \).}
{\( 7 \)}{\( 8 \)}{\( 6 \)}{\( 10 \)}{\( 5 \)}{}}
\LANGsk{
\Question{
\(n\)-tý člen geometrickej postupnosti je rovný \(\frac52 \), kvocient je \( \frac12 \) a štvrtý člen je \( 20 \). Určte \( n \).}
{\( 7 \)}{\( 8 \)}{\( 6 \)}{\( 10 \)}{\( 5 \)}{}}
\LANGes{
\Question{
El término \(n\)-ésimo de una progresión geométrica es igual a \(\frac52 \), su razón es \( \frac12 \) y el cuarto término es \( 20 \). Averigua \( n \).}
{\( 7 \)}{\( 8 \)}{\( 6 \)}{\( 10 \)}{\( 5 \)}{}}
\End

\Data\ID{1003112809}
\Updated{2022-01-06 01:01:10}
\Level{A}
\SubArea{Geometric Sequences}
\LANGen{
\Question{
The first term of a geometric sequence is \( 16 \) and the fourth term equals \( 54 \). Which of the following inequalities is true for the second term of the sequence?}
{\( a_2 > a_1 \)}{\( a_2 < a_1 \)}{\( a_2 < 0 \)}{\( a_2 > a_4 \)}{\( a_2 < \frac{a_4}{a_1} \)}{}}
\LANGpl{
\Question{
Pierwszy wyraz ciągu geometrycznego wynosi \( 16 \), a czwarty wyraz równy jest \( 54 \). Która z poniższych nierówności jest prawdziwa dla drugiego wyrazu ciągu?}
{\( a_2 > a_1 \)}{\( a_2 < a_1 \)}{\( a_2 < 0 \)}{\( a_2 > a_4 \)}{\( a_2 < \frac{a_4}{a_1} \)}{}}
\LANGcs{
\Question{
První člen geometrické posloupnosti je roven \( 16 \) a čtvrtý člen je \( 54 \). Pro druhý člen této posloupnosti platí:}
{\( a_2 > a_1 \)}{\( a_2 < a_1 \)}{\( a_2 < 0 \)}{\( a_2 > a_4 \)}{\( a_2 < \frac{a_4}{a_1} \)}{}}
\LANGsk{
\Question{
Prvý člen geometrickej postupnosti je rovný \( 16 \) a štvrtý člen je \( 54 \). Pre druhý člen tejto postupnosti platí:}
{\( a_2 > a_1 \)}{\( a_2 < a_1 \)}{\( a_2 < 0 \)}{\( a_2 > a_4 \)}{\( a_2 < \frac{a_4}{a_1} \)}{}}
\LANGes{
\Question{
El primer término de una progresión geométrica es \( 16 \) y su cuarto término es \( 54 \). Para el segundo término se cumple que:}
{\( a_2 > a_1 \)}{\( a_2 < a_1 \)}{\( a_2 < 0 \)}{\( a_2 > a_4 \)}{\( a_2 < \frac{a_4}{a_1} \)}{}}
\End

\Data\ID{1003112811}
\Updated{2022-01-06 01:01:54}
\Level{A}
\SubArea{Geometric Sequences}
\LANGen{
\Question{
The third term of a geometric sequence is \( 27 \) and the sixth term equals \( 1 \). Find the common ratio of this sequence:}
{\( \frac13 \)}{\( 3 \)}{\( \sqrt{27} \)}{\( 27 \)}{\( \frac1{27} \)}{}}
\LANGpl{
\Question{
Trzeci wyraz ciągu geometrycznego wynosi \( 27 \), a szósty wyraz jest równy \( 1 \). Wyznacz wspólny współczynnik tego ciągu:}
{\( \frac13 \)}{\( 3 \)}{\( \sqrt{27} \)}{\( 27 \)}{\( \frac1{27} \)}{}}
\LANGcs{
\Question{
Třetí člen geometrické posloupnosti je roven \( 27 \) a šestý člen je \( 1 \). Určete kvocient této posloupnosti.}
{\( \frac13 \)}{\( 3 \)}{\( \sqrt{27} \)}{\( 27 \)}{\( \frac1{27} \)}{}}
\LANGsk{
\Question{
Tretí člen geometrickej postupnosti je rovný \( 27 \) a šiesty člen je \( 1 \). Určte kvocient tejto postupnosti.}
{\( \frac13 \)}{\( 3 \)}{\( \sqrt{27} \)}{\( 27 \)}{\( \frac1{27} \)}{}}
\LANGes{
\Question{
El tercer término de una progresión geométrica es igual a \( 27 \) y su sexto término es \( 1 \). Calcula la razón de la progresión.}
{\( \frac13 \)}{\( 3 \)}{\( \sqrt{27} \)}{\( 27 \)}{\( \frac1{27} \)}{}}
\End

\Data\ID{1003124701}
\Updated{2022-01-06 01:01:02}
\Level{A}
\SubArea{Geometric Sequences}
\LANGen{
\Question{
The third term of a geometric sequence is \( 9 \) and the common ratio is \( 3 \). Find the recursive formula of the sequence.}
{\( a_1=1 \), \( a_{n+1}=3a_n \)}{\( a_1=3 \), \( a_{n+1}=a_n+3 \)}{\( a_1=9 \), \( a_{n+1}=3a_n \)}{\( a_1=3 \), \( a_{n+1}=a_n^2 \)}{\( a_1=1\), \(a_{n+1}=\frac13a_n \)}{}}
\LANGpl{
\Question{
Trzeci wyraz ciągu geometrycznego wynosi  \( 9 \), a wspólny współczynnik \( 3 \). Wyznacz wzór rekurencyjny tego ciągu.}
{\( a_1=1 \), \( a_{n+1}=3a_n \)}{\( a_1=3 \), \( a_{n+1}=a_n+3 \)}{\( a_1=9 \), \( a_{n+1}=3a_n \)}{\( a_1=3 \), \( a_{n+1}=a_n^2 \)}{\( a_1=1\), \(a_{n+1}=\frac13a_n \)}{}}
\LANGcs{
\Question{
Najděte rekurentní vyjádření geometrické posloupnosti, je-li její třetí člen roven \( 9 \) a kvocient je \( 3 \).}
{\( a_1=1 \), \( a_{n+1}=3a_n \)}{\( a_1=3 \), \( a_{n+1}=a_n+3 \)}{\( a_1=9 \), \( a_{n+1}=3a_n \)}{\( a_1=3 \), \( a_{n+1}=a_n^2 \)}{\( a_1=1\), \(a_{n+1}=\frac13a_n \)}{}}
\LANGsk{
\Question{
Nájdite rekurentné vyjadrenie geometrickej postupnosti, ak je jej tretí člen rovný \( 9 \) a kvocient je \( 3 \).}
{\( a_1=1 \), \( a_{n+1}=3a_n \)}{\( a_1=3 \), \( a_{n+1}=a_n+3 \)}{\( a_1=9 \), \( a_{n+1}=3a_n \)}{\( a_1=3 \), \( a_{n+1}=a_n^2 \)}{\( a_1=1\), \(a_{n+1}=\frac13a_n \)}{}}
\LANGes{
\Question{
Halla la fórmula recursiva de una progresión geométrica, cuyo tercer término es igual a \( 9 \) y cuya razón es \( 3 \).}
{\( a_1=1 \), \( a_{n+1}=3a_n \)}{\( a_1=3 \), \( a_{n+1}=a_n+3 \)}{\( a_1=9 \), \( a_{n+1}=3a_n \)}{\( a_1=3 \), \( a_{n+1}=a_n^2 \)}{\( a_1=1\), \(a_{n+1}=\frac13a_n \)}{}}
\End

\Data\ID{1003124702}
\Updated{2022-01-06 01:01:21}
\Level{A}
\SubArea{Geometric Sequences}
\LANGen{
\Question{
Find the recursive formula of the geometric sequence \( a_n=2\cdot 3^n \), \( n\in\mathbb{N} \).}
{\( a_1=6 \), \( a_{n+1} = 3a_n \)}{\( a_1=2 \), \( a_{n+1} = 3a_n \)}{\( a_1=3 \), \( a_{n+1} = 6a_n \)}{\( a_1=6 \), \( a_{n+1} = \frac13a_n \)}{\( a_1=2 \), \( a_{n+1} = a_n+3 \)}{}}
\LANGpl{
\Question{
Wyznacz wzór rekurencyjny ciągu geometrycznego \( a_n=2\cdot 3^n \), \( n\in\mathbb{N} \).}
{\( a_1=6 \), \( a_{n+1} = 3a_n \)}{\( a_1=2 \), \( a_{n+1} = 3a_n \)}{\( a_1=3 \), \( a_{n+1} = 6a_n \)}{\( a_1=6 \), \( a_{n+1} = \frac13a_n \)}{\( a_1=2 \), \( a_{n+1} = a_n+3 \)}{}}
\LANGcs{
\Question{
Najděte rekurentní vyjádření geometrické posloupnosti, je-li \( a_n=2\cdot 3^n \), \( n\in\mathbb{N} \).}
{\( a_1=6 \), \( a_{n+1} = 3a_n \)}{\( a_1=2 \), \( a_{n+1} = 3a_n \)}{\( a_1=3 \), \( a_{n+1} = 6a_n \)}{\( a_1=6 \), \( a_{n+1} = \frac13a_n \)}{\( a_1=2 \), \( a_{n+1} = a_n+3 \)}{}}
\LANGsk{
\Question{
Nájdite rekurentné vyjadrenie geometrickej postupnosti, ak je \( a_n=2\cdot 3^n \), \( n\in\mathbb{N} \).}
{\( a_1=6 \), \( a_{n+1} = 3a_n \)}{\( a_1=2 \), \( a_{n+1} = 3a_n \)}{\( a_1=3 \), \( a_{n+1} = 6a_n \)}{\( a_1=6 \), \( a_{n+1} = \frac13a_n \)}{\( a_1=2 \), \( a_{n+1} = a_n+3 \)}{}}
\LANGes{
\Question{
Halla la fórmula recursiva para una progresión geométrica si \( a_n=2\cdot 3^n \), \( n\in\mathbb{N} \).}
{\( a_1=6 \), \( a_{n+1} = 3a_n \)}{\( a_1=2 \), \( a_{n+1} = 3a_n \)}{\( a_1=3 \), \( a_{n+1} = 6a_n \)}{\( a_1=6 \), \( a_{n+1} = \frac13a_n \)}{\( a_1=2 \), \( a_{n+1} = a_n+3 \)}{}}
\End

\Data\ID{1003124703}
\Updated{2022-01-06 01:01:42}
\Level{A}
\SubArea{Geometric Sequences}
\LANGen{
\Question{
The second term of a geometric sequence is \( 15 \) and the third term is \( 3 \). Find the recursive formula of the sequence.}
{\( a_1=75 \), \( a_{n+1} = \frac15a_n \)}{\( a_1=3 \), \( a_{n+1} = 5a_n \)}{\( a_1=\frac35 \), \( a_{n+1} = \frac15a_n \)}{\( a_1=\frac35 \), \( a_{n+1} = 5a_n \)}{\( a_1=27 \), \( a_{n+1} = a_n-12 \)}{}}
\LANGpl{
\Question{
Drugi wyraz ciągu geometrycznego wynosi \( 15 \), a trzeci wyraz ciągu jest równy \( 3 \). Wyznacz wzór rekurencyjny tego ciągu.}
{\( a_1=75 \), \( a_{n+1} = \frac15a_n \)}{\( a_1=3 \), \( a_{n+1} = 5a_n \)}{\( a_1=\frac35 \), \( a_{n+1} = \frac15a_n \)}{\( a_1=\frac35 \), \( a_{n+1} = 5a_n \)}{\( a_1=27 \), \( a_{n+1} = a_n-12 \)}{}}
\LANGcs{
\Question{
Najděte rekurentní vyjádření geometrické posloupnosti, je-li její druhý člen roven \( 15 \) a třetí člen je \( 3 \).}
{\( a_1=75 \), \( a_{n+1} = \frac15a_n \)}{\( a_1=3 \), \( a_{n+1} = 5a_n \)}{\( a_1=\frac35 \), \( a_{n+1} = \frac15a_n \)}{\( a_1=\frac35 \), \( a_{n+1} = 5a_n \)}{\( a_1=27 \), \( a_{n+1} = a_n-12 \)}{}}
\LANGsk{
\Question{
Nájdite rekurentné vyjadrenie geometrickej postupnosti, ak je jej druhý člen rovný \( 15 \) a tretí člen je \( 3 \).}
{\( a_1=75 \), \( a_{n+1} = \frac15a_n \)}{\( a_1=3 \), \( a_{n+1} = 5a_n \)}{\( a_1=\frac35 \), \( a_{n+1} = \frac15a_n \)}{\( a_1=\frac35 \), \( a_{n+1} = 5a_n \)}{\( a_1=27 \), \( a_{n+1} = a_n-12 \)}{}}
\LANGes{
\Question{
Halla la fórmula recursiva de una progresión geométrica si su segundo término es \( 15 \) y el tercero es \( 3 \).}
{\( a_1=75 \), \( a_{n+1} = \frac15a_n \)}{\( a_1=3 \), \( a_{n+1} = 5a_n \)}{\( a_1=\frac35 \), \( a_{n+1} = \frac15a_n \)}{\( a_1=\frac35 \), \( a_{n+1} = 5a_n \)}{\( a_1=27 \), \( a_{n+1} = a_n-12 \)}{}}
\End

\Data\ID{1003124704}
\Updated{2022-01-06 01:01:01}
\Level{A}
\SubArea{Geometric Sequences}
\LANGen{
\Question{
The \( 10 \)th term of a geometric sequence is \( 1 \) and the \( 15 \)th term is \( -1 \). Find the recursive formula of the sequence.}
{\( a_1=-1 \), \( a_{n+1}=-a_n \)}{\( a_1=1 \), \( a_{n+1}=-a_n \)}{\( a_1=-1 \), \( a_{n+1}=a_n \)}{\( a_1=1 \), \( a_{n+1}=a_n \)}{\( a_1=-1 \), \( a_{n+1}=a_n-1 \)}{}}
\LANGpl{
\Question{
\( 10 \)-ty  wyraz ciągu geometrycznego wynosi  \( 1 \), a \( 15 \)-ty wyraz  \( -1 \). Wyznacz wzór rekurencyjny tego ciągu.}
{\( a_1=-1 \), \( a_{n+1}=-a_n \)}{\( a_1=1 \), \( a_{n+1}=-a_n \)}{\( a_1=-1 \), \( a_{n+1}=a_n \)}{\( a_1=1 \), \( a_{n+1}=a_n \)}{\( a_1=-1 \), \( a_{n+1}=a_n-1 \)}{}}
\LANGcs{
\Question{
Desátý člen geometrické posloupnosti je roven \( 1 \) a patnáctý člen je \( -1 \). Najděte její rekurentní vyjádření.}
{\( a_1=-1 \), \( a_{n+1}=-a_n \)}{\( a_1=1 \), \( a_{n+1}=-a_n \)}{\( a_1=-1 \), \( a_{n+1}=a_n \)}{\( a_1=1 \), \( a_{n+1}=a_n \)}{\( a_1=-1 \), \( a_{n+1}=a_n-1 \)}{}}
\LANGsk{
\Question{
Desiaty člen geometrickej postupnosti je rovný \( 1 \) a pätnásty člen je \( -1 \). Nájdite jej rekurentné vyjadrenie.}
{\( a_1=-1 \), \( a_{n+1}=-a_n \)}{\( a_1=1 \), \( a_{n+1}=-a_n \)}{\( a_1=-1 \), \( a_{n+1}=a_n \)}{\( a_1=1 \), \( a_{n+1}=a_n \)}{\( a_1=-1 \), \( a_{n+1}=a_n-1 \)}{}}
\LANGes{
\Question{
El término décimo de una progresión geométrica es igual a \( 1 \) y su término decimoquinto es \( -1 \). Halla su fórmula recursiva.}
{\( a_1=-1 \), \( a_{n+1}=-a_n \)}{\( a_1=1 \), \( a_{n+1}=-a_n \)}{\( a_1=-1 \), \( a_{n+1}=a_n \)}{\( a_1=1 \), \( a_{n+1}=a_n \)}{\( a_1=-1 \), \( a_{n+1}=a_n-1 \)}{}}
\End

\Data\ID{1003124705}
\Updated{2022-01-06 01:01:19}
\Level{A}
\SubArea{Geometric Sequences}
\LANGen{
\Question{
The third term of a geometric sequence is \( 3 \) and the common ratio is \( 3 \). Find the \( n \)th term.}
{\( a_n=3^{n-2} \), \( n\in\mathbb{N} \)}{\( a_n=3^{n-1} \), \( n\in\mathbb{N} \)}{\( a_n=3^{n} \), \( n\in\mathbb{N} \)}{\( a_n=\frac3n \), \( n\in\mathbb{N} \)}{\( a_n=3n \), \( n\in\mathbb{N} \)}{}}
\LANGpl{
\Question{
Trzeci wyraz ciągu geometrycznego wynosi \( 3 \), a wspólny współczynnik \( 3 \). Wyznacz \( n \)-ty wyraz tego ciągu.}
{\( a_n=3^{n-2} \), \( n\in\mathbb{N} \)}{\( a_n=3^{n-1} \), \( n\in\mathbb{N} \)}{\( a_n=3^{n} \), \( n\in\mathbb{N} \)}{\( a_n=\frac3n \), \( n\in\mathbb{N} \)}{\( a_n=3n \), \( n\in\mathbb{N} \)}{}}
\LANGcs{
\Question{
Třetí člen geometrické posloupnosti je roven \( 3 \) a kvocient je \( 3 \). Najděte vzorec pro \( n \)-tý člen.}
{\( a_n=3^{n-2} \), \( n\in\mathbb{N} \)}{\( a_n=3^{n-1} \), \( n\in\mathbb{N} \)}{\( a_n=3^{n} \), \( n\in\mathbb{N} \)}{\( a_n=\frac3n \), \( n\in\mathbb{N} \)}{\( a_n=3n \), \( n\in\mathbb{N} \)}{}}
\LANGsk{
\Question{
Tretí člen geometrickej postupnosti je rovný \( 3 \) a kvocient je \( 3 \). Nájdite vzorec pre \( n \)-tý člen.}
{\( a_n=3^{n-2} \), \( n\in\mathbb{N} \)}{\( a_n=3^{n-1} \), \( n\in\mathbb{N} \)}{\( a_n=3^{n} \), \( n\in\mathbb{N} \)}{\( a_n=\frac3n \), \( n\in\mathbb{N} \)}{\( a_n=3n \), \( n\in\mathbb{N} \)}{}}
\LANGes{
\Question{
El tercer término de una progresión geométrica es igual a \( 3 \) y su razón es \( 3 \). Halla la fórmula para el término \( n \)-ésimo.}
{\( a_n=3^{n-2} \), \( n\in\mathbb{N} \)}{\( a_n=3^{n-1} \), \( n\in\mathbb{N} \)}{\( a_n=3^{n} \), \( n\in\mathbb{N} \)}{\( a_n=\frac3n \), \( n\in\mathbb{N} \)}{\( a_n=3n \), \( n\in\mathbb{N} \)}{}}
\End

\Data\ID{1003124706}
\Updated{2022-01-06 01:01:41}
\Level{A}
\SubArea{Geometric Sequences}
\LANGen{
\Question{
The first term of a geometric sequence is \( 5 \) and the fourth term is \( 40 \). Find the \( n \)th term.}
{\( a_n=5\cdot2^{n-1} \), \( n\in\mathbb{N} \)}{\( a_n=\frac{5n}2 \), \( n\in\mathbb{N} \)}{\( a_n=5\cdot2^n \), \( n\in\mathbb{N} \)}{\( a_n=5n \), \( n\in\mathbb{N} \)}{\( a_n=5\cdot\left(2^{n}-1\right) \), \( n\in\mathbb{N} \)}{}}
\LANGpl{
\Question{
Pierwszy wyraz ciągu geometrycznego wynosi  \( 5 \), a czwarty jest równy \( 40 \). Wyznacz \( n \)-ty wyraz tego ciągu.}
{\( a_n=5\cdot2^{n-1} \), \( n\in\mathbb{N} \)}{\( a_n=\frac{5n}2 \), \( n\in\mathbb{N} \)}{\( a_n=5\cdot2^n \), \( n\in\mathbb{N} \)}{\( a_n=5n \), \( n\in\mathbb{N} \)}{\( a_n=5\cdot\left(2^{n}-1\right) \), \( n\in\mathbb{N} \)}{}}
\LANGcs{
\Question{
První člen geometrické posloupnosti je roven \( 5 \) a čtvrtý člen je \( 40 \). Najděte vzorec pro \( n \)-tý člen.}
{\( a_n=5\cdot2^{n-1} \), \( n\in\mathbb{N} \)}{\( a_n=\frac{5n}2 \), \( n\in\mathbb{N} \)}{\( a_n=5\cdot2^n \), \( n\in\mathbb{N} \)}{\( a_n=5n \), \( n\in\mathbb{N} \)}{\( a_n=5\cdot\left(2^{n}-1\right) \), \( n\in\mathbb{N} \)}{}}
\LANGsk{
\Question{
Prvý člen geometrickej postupnosti je rovný \( 5 \) a štvrtý člen je \( 40 \). Nájdite vzorec pre \( n \)-tý člen.}
{\( a_n=5\cdot2^{n-1} \), \( n\in\mathbb{N} \)}{\( a_n=\frac{5n}2 \), \( n\in\mathbb{N} \)}{\( a_n=5\cdot2^n \), \( n\in\mathbb{N} \)}{\( a_n=5n \), \( n\in\mathbb{N} \)}{\( a_n=5\cdot\left(2^{n}-1\right) \), \( n\in\mathbb{N} \)}{}}
\LANGes{
\Question{
El primer término de una progresión geométrica es igual a \( 5 \) y su cuarto término es \( 40 \). Halla la fórmula para el término \( n \)-ésimo.}
{\( a_n=5\cdot2^{n-1} \), \( n\in\mathbb{N} \)}{\( a_n=\frac{5n}2 \), \( n\in\mathbb{N} \)}{\( a_n=5\cdot2^n \), \( n\in\mathbb{N} \)}{\( a_n=5n \), \( n\in\mathbb{N} \)}{\( a_n=5\cdot\left(2^{n}-1\right) \), \( n\in\mathbb{N} \)}{}}
\End

\Data\ID{1003124707}
\Updated{2022-01-10 08:01:09}
\Level{A}
\SubArea{Geometric Sequences}
\LANGen{
\Question{
The \( n \)th term of a geometric sequence is \( 2\cdot3^{n-2} \). Find the third term and the common ratio.}
{\( a_3=6 \), \( q=3 \)}{\( a_3=3 \), \( q=-2 \)}{\( a_3=6 \), \( q=-3 \)}{\( a_3=6 \), \( q=2 \)}{\( a_3=3 \), \( q=6 \)}{}}
\LANGpl{
\Question{
\( n \)-ty wyraz ciągu geometrycznego wynosi  \( 2\cdot3^{n-2} \). Wyznacz trzeci wyraz ciągu i wspólny współczynnik.}
{\( a_3=6 \), \( q=3 \)}{\( a_3=3 \), \( q=-2 \)}{\( a_3=6 \), \( q=-3 \)}{\( a_3=6 \), \( q=2 \)}{\( a_3=3 \), \( q=6 \)}{}}
\LANGcs{
\Question{
\( n \)-tý člen geometrické posloupnosti je roven \( 2\cdot3^{n-2} \). Určete třetí člen a kvocient této posloupnosti.}
{\( a_3=6 \), \( q=3 \)}{\( a_3=3 \), \( q=-2 \)}{\( a_3=6 \), \( q=-3 \)}{\( a_3=6 \), \( q=2 \)}{\( a_3=3 \), \( q=6 \)}{}}
\LANGsk{
\Question{
\( n \)-tý člen geometrickej postupnosti je rovný \( 2\cdot3^{n-2} \). Určte tretí člen a kvocient tejto postupnosti.}
{\( a_3=6 \), \( q=3 \)}{\( a_3=3 \), \( q=-2 \)}{\( a_3=6 \), \( q=-3 \)}{\( a_3=6 \), \( q=2 \)}{\( a_3=3 \), \( q=6 \)}{}}
\LANGes{
\Question{
El término \( n \)-ésimo de una progresión geométrica es igual a \( 2\cdot3^{n-2} \). Averigua el tercer término y la razón de la progresión.}
{\( a_3=6 \), \( q=3 \)}{\( a_3=3 \), \( q=-2 \)}{\( a_3=6 \), \( q=-3 \)}{\( a_3=6 \), \( q=2 \)}{\( a_3=3 \), \( q=6 \)}{}}
\End

\Data\ID{1003124708}
\Updated{2022-04-23 05:04:14}
\Level{A}
\SubArea{Geometric Sequences}
\LANGen{
\Question{
The \( n \)th term of a geometric sequence is \( 3^{n-1}\cdot2^{n+1} \). Find the second term and the common ratio.}
{\( a_2=24 \), \( q=6 \)}{\( a_2=6 \), \( q=6 \)}{\( a_2=2 \), \( q=3 \)}{\( a_2=24 \), \( q=2 \)}{\( a_2=3 \), \( q=3 \)}{}}
\LANGpl{
\Question{
\( n \)- ty wyraz ciągu geometrycznego wynosi  \( 3^{n-1}\cdot2^{n+1} \). Wyznacz drugi wyraz ciągu i wspólny współczynnik.}
{\( a_2=24 \), \( q=6 \)}{\( a_2=6 \), \( q=6 \)}{\( a_2=2 \), \( q=3 \)}{\( a_2=24 \), \( q=2 \)}{\( a_2=3 \), \( q=3 \)}{}}
\LANGcs{
\Question{
\( n \)-tý člen geometrické posloupnosti je roven \( 3^{n-1}\cdot2^{n+1} \). Určete druhý člen a kvocient této posloupnosti.}
{\( a_2=24 \), \( q=6 \)}{\( a_2=6 \), \( q=6 \)}{\( a_2=2 \), \( q=3 \)}{\( a_2=24 \), \( q=2 \)}{\( a_2=3 \), \( q=3 \)}{}}
\LANGsk{
\Question{
\( n \)-tý člen geometrickej postupnosti je rovný \( 3^{n-1}\cdot2^{n+1} \). Určte druhý člen a kvocient tejto postupnosti.}
{\( a_2=24 \), \( q=6 \)}{\( a_2=6 \), \( q=6 \)}{\( a_2=2 \), \( q=3 \)}{\( a_2=24 \), \( q=2 \)}{\( a_2=3 \), \( q=3 \)}{}}
\LANGes{
\Question{
El término \( n \)-ésimo de una progresión geométrica es igual a \( 3^{n-1}\cdot2^{n+1} \). Halla el segundo término de la progresión y su razón.}
{\( a_2=24 \), \( q=6 \)}{\( a_2=6 \), \( q=6 \)}{\( a_2=2 \), \( q=3 \)}{\( a_2=24 \), \( q=2 \)}{\( a_2=3 \), \( q=3 \)}{}}
\End

\Data\ID{1003134601}
\Updated{2022-01-10 08:01:36}
\Level{A}
\SubArea{Geometric Sequences}
\LANGen{
\Question{
The fifth term of a geometric sequence is $4\sqrt2$ and the second term is $2$. Find the fourth term.}
{$4$}{$2\sqrt2$}{$3\sqrt2$}{$\sqrt2$}{$2$}{}}
\LANGpl{
\Question{
Piąty wyraz ciągu geometrycznego wynosi $4\sqrt2$, a drugi wyraz jest równy  $2$. Wyznacz czwarty wyraz ciągu.}
{$4$}{$2\sqrt2$}{$3\sqrt2$}{$\sqrt2$}{$2$}{}}
\LANGcs{
\Question{
Určete čtvrtý člen geometrické posloupnosti, je-li pátý člen $4\sqrt2$ a druhý člen je $2$.}
{$4$}{$2\sqrt2$}{$3\sqrt2$}{$\sqrt2$}{$2$}{}}
\LANGsk{
\Question{
Určte štvrtý člen geometrickej postupnosti, ak je piaty člen $4\sqrt2$ a druhý člen je $2$.}
{$4$}{$2\sqrt2$}{$3\sqrt2$}{$\sqrt2$}{$2$}{}}
\LANGes{
\Question{
Averigua el cuarto término de una progresión geométrica si su quinto término es $4\sqrt2$ y el segundo término es $2$.}
{$4$}{$2\sqrt2$}{$3\sqrt2$}{$\sqrt2$}{$2$}{}}
\End

\Data\ID{2000014011}
\Updated{2022-11-02 09:11:42}
\Level{A}
\SubArea{Geometric Sequences}
\IMG{2010014001-figure16.pdf}
\LANGen{
\Question{
The picture shows a part of the graph of a geometric sequence. What is the common ratio of the geometric sequence?}
{\(-\frac12\)}{\(-2\)}{\(\frac12\)}{\(\left(-\frac12\right)^n\)}{}{}}
\LANGpl{
\Question{
Rysunek przedstawia fragment wykresu ciągu geometrycznego. Jaki jest iloraz ciągu geometrycznego?}
{\(-\frac12\)}{\(-2\)}{\(\frac12\)}{\(\left(-\frac12\right)^n\)}{}{}}
\LANGcs{
\Question{
Na obrázku je část grafu geometrické posloupnosti. Jaký je její kvocient?}
{\(-\frac12\)}{\(-2\)}{\(\frac12\)}{\(\left(-\frac12\right)^n\)}{}{}}
\LANGsk{
\Question{
Na obrázku je časť grafu geometrickej postupnosti. Aký je kvocient tejto geometrickej postupnosti?}
{\(-\frac12\)}{\(-2\)}{\(\frac12\)}{\(\left(-\frac12\right)^n\)}{}{}}
\LANGes{
\Question{
La imagen muestra una parte de la gráfica de una sucesión geométrica. ¿Cuál es la razón de la sucesión geométrica?}
{\(-\frac12\)}{\(-2\)}{\(\frac12\)}{\(\left(-\frac12\right)^n\)}{}{}}
\End

\Data\ID{2010004903}
\Updated{2022-08-25 10:08:52}
\Level{A}
\SubArea{Geometric Sequences}
\LANGen{
\Question{
The seventh term of a geometric sequence is \( 32 \) and the tenth term is \( 4 \). Choose the correct formula to find the eighth term of this sequence.}
{\( a_8=32\cdot\sqrt[3]{\frac4{32}} \)}{\( a_8=32\cdot\sqrt[3]{\frac{32}4} \)}{\( a_8=4\cdot\sqrt[3]{\frac4{32}} \)}{\( a_8=4\cdot\sqrt[3]{\frac{32}4} \)}{\( a_8=8\cdot\sqrt[3]{\frac3{24}} \)}{}}
\LANGpl{
\Question{
Siódmy wyraz ciągu geometrycznego to \( 32 \), a dziesiąty to \( 4 \). Wybierz poprawny wzór, aby znaleźć ósmy wyraz tego ciągu.}
{\( a_8=32\cdot\sqrt[3]{\frac4{32}} \)}{\( a_8=32\cdot\sqrt[3]{\frac{32}4} \)}{\( a_8=4\cdot\sqrt[3]{\frac4{32}} \)}{\( a_8=4\cdot\sqrt[3]{\frac{32}4} \)}{\( a_8=8\cdot\sqrt[3]{\frac3{24}} \)}{}}
\LANGcs{
\Question{
Sedmý člen geometrické posloupnosti je \( 32 \) a desátý člen \( 4 \). Určete správný postup pro výpočet osmého členu této posloupnosti.}
{\( a_8=32\cdot\sqrt[3]{\frac4{32}} \)}{\( a_8=32\cdot\sqrt[3]{\frac{32}4} \)}{\( a_8=4\cdot\sqrt[3]{\frac4{32}} \)}{\( a_8=4\cdot\sqrt[3]{\frac{32}4} \)}{\( a_8=8\cdot\sqrt[3]{\frac3{24}} \)}{}}
\LANGsk{
\Question{
Siedmy člen geometrickej postupnosti je rovný \( 32 \) a desiaty člen je \( 4 \). Zvoľte správny postup pre výpočet ôsmeho člena tejto postupnosti.}
{\( a_8=32\cdot\sqrt[3]{\frac4{32}} \)}{\( a_8=32\cdot\sqrt[3]{\frac{32}4} \)}{\( a_8=4\cdot\sqrt[3]{\frac4{32}} \)}{\( a_8=4\cdot\sqrt[3]{\frac{32}4} \)}{\( a_8=8\cdot\sqrt[3]{\frac3{24}} \)}{}}
\LANGes{
\Question{
El séptimo término de una progresión geométrica es \( 32 \) y el décimo término es \( 4 \). Averigua la fórmula correcta para el cálculo del octavo término.}
{\( a_8=32\cdot\sqrt[3]{\frac4{32}} \)}{\( a_8=32\cdot\sqrt[3]{\frac{32}4} \)}{\( a_8=4\cdot\sqrt[3]{\frac4{32}} \)}{\( a_8=4\cdot\sqrt[3]{\frac{32}4} \)}{\( a_8=8\cdot\sqrt[3]{\frac3{24}} \)}{}}
\End

\Data\ID{2010004904}
\Updated{2022-08-25 10:08:52}
\Level{A}
\SubArea{Geometric Sequences}
\LANGen{
\Question{
The \(n\)th term of a geometric sequence is \(\frac29 \), the common ratio is \( \frac13 \) and the fourth term is \( 6 \). Find \( n \).}
{\( 7 \)}{\( 8 \)}{\( 6 \)}{\( 10 \)}{\( 5 \)}{}}
\LANGpl{
\Question{
Ciąg geometryczny ma \(n\)-ty wyraz \(\frac29 \), iloraz to \( \frac13 \), a czwarty wyraz to \( 6 \). Znajdź \( n \).}
{\( 7 \)}{\( 8 \)}{\( 6 \)}{\( 10 \)}{\( 5 \)}{}}
\LANGcs{
\Question{
\(n\)-tý člen geometrické posloupnosti je \(\frac29 \), kvocient je \( \frac13 \) a čtvrtý člen je \( 6 \). Určete \( n \).}
{\( 7 \)}{\( 8 \)}{\( 6 \)}{\( 10 \)}{\( 5 \)}{}}
\LANGsk{
\Question{
\(n\)-tý člen geometrickej postupnosti je rovný \(\frac29 \), kvocient je \( \frac13 \) a štvrtý člen je \( 6 \). Určte \( n \).}
{\( 7 \)}{\( 8 \)}{\( 6 \)}{\( 10 \)}{\( 5 \)}{}}
\LANGes{
\Question{
\(n\)-ésimo término de una progresión geométrica es \(\frac29 \), su razón es \( \frac13 \) y el cuarto término es \( 6 \). Averigua \( n \).}
{\( 7 \)}{\( 8 \)}{\( 6 \)}{\( 10 \)}{\( 5 \)}{}}
\End

\Data\ID{2010004905}
\Updated{2022-08-25 10:08:52}
\Level{A}
\SubArea{Geometric Sequences}
\LANGen{
\Question{
The \( n \)th term of a geometric sequence is \( 4^{n-1}\cdot5^{2-n} \). Find the second term and the common ratio.}
{\( a_2=4 \), \( q=\frac45 \)}{\( a_2=4 \), \( q=\frac54 \)}{\( a_2=5 \), \( q=\frac45 \)}{\( a_2=4 \), \( q=20 \)}{\( a_2=5 \), \( q=20 \)}{}}
\LANGpl{
\Question{
Ciąg geometryczny ma \( n \)-ty wyraz \( 4^{n-1}\cdot5^{2-n} \). Znajdź drugi wyraz i iloraz tego ciągu.}
{\( a_2=4 \), \( q=\frac45 \)}{\( a_2=4 \), \( q=\frac54 \)}{\( a_2=5 \), \( q=\frac45 \)}{\( a_2=4 \), \( q=20 \)}{\( a_2=5 \), \( q=20 \)}{}}
\LANGcs{
\Question{
\( n \)-tý člen geometrické posloupnosti je roven \( 4^{n-1}\cdot5^{2-n} \). Určete druhý člen a kvocient této posloupnosti.}
{\( a_2=4 \), \( q=\frac45 \)}{\( a_2=4 \), \( q=\frac54 \)}{\( a_2=5 \), \( q=\frac45 \)}{\( a_2=4 \), \( q=20 \)}{\( a_2=5 \), \( q=20 \)}{}}
\LANGsk{
\Question{
\( n \)-tý člen geometrickej postupnosti je rovný \( 4^{n-1}\cdot5^{2-n} \). Určte druhý člen a kvocient tejto postupnosti.}
{\( a_2=4 \), \( q=\frac45 \)}{\( a_2=4 \), \( q=\frac54 \)}{\( a_2=5 \), \( q=\frac45 \)}{\( a_2=4 \), \( q=20 \)}{\( a_2=5 \), \( q=20 \)}{}}
\LANGes{
\Question{
\( n \)-ésimo término de una progresión geométrica \( 4^{n-1}\cdot5^{2-n} \). Averigua el segundo término y la razón.}
{\( a_2=4 \), \( q=\frac45 \)}{\( a_2=4 \), \( q=\frac54 \)}{\( a_2=5 \), \( q=\frac45 \)}{\( a_2=4 \), \( q=20 \)}{\( a_2=5 \), \( q=20 \)}{}}
\End

\Data\ID{2010004906}
\Updated{2022-08-25 10:08:52}
\Level{A}
\SubArea{Geometric Sequences}
\IMG{2010004901_6-figure0.pdf}
\LANGen{
\Question{
The picture shows a part of the graph of the geometric sequence. What is the explicit formula for the \(n\)th term of the following geometric sequence?}
{\( a_n =2\cdot5^{n-1}\)}{\( a_n =2\cdot5^{n}\)}{\( a_n =5^{n-1}\)}{\( a_n =2\cdot5^{n+1}\)}{}{}}
\LANGpl{
\Question{
Rysunek przedstawia fragment wykresu ciągu geometrycznego. Jaki jest wzór na \(n\)-ty wyraz tego ciągu geometrycznego?}
{\( a_n =2\cdot5^{n-1}\)}{\( a_n =2\cdot5^{n}\)}{\( a_n =5^{n-1}\)}{\( a_n =2\cdot5^{n+1}\)}{}{}}
\LANGcs{
\Question{
Na obrázku je část grafu geometrické posloupnosti. Jaký je vzorec pro \(n\)-tý člen této posloupnosti?}
{\( a_n =2\cdot5^{n-1}\)}{\( a_n =2\cdot5^{n}\)}{\( a_n =5^{n-1}\)}{\( a_n =2\cdot5^{n+1}\)}{}{}}
\LANGsk{
\Question{
Na obrázku je časť grafu geometrickej postupnosti. Vzorec pre \(n\)-tý člen tejto postupnosti je:}
{\( a_n =2\cdot5^{n-1}\)}{\( a_n =2\cdot5^{n}\)}{\( a_n =5^{n-1}\)}{\( a_n =2\cdot5^{n+1}\)}{}{}}
\LANGes{
\Question{
En el dibujo tenemos una parte de la gráfica de una progresión geométrica. ¿Cuál es la fórmula para calcular el término \(n\)-ésimo?}
{\( a_n =2\cdot5^{n-1}\)}{\( a_n =2\cdot5^{n}\)}{\( a_n =5^{n-1}\)}{\( a_n =2\cdot5^{n+1}\)}{}{}}
\End

\Data\ID{2010004907}
\Updated{2022-08-25 10:08:52}
\Level{A}
\SubArea{Geometric Sequences}
\IMG{2010004901_6-figure1.pdf}
\LANGen{
\Question{
The picture shows a part of the graph of the geometric sequence. What is the explicit formula for the \(n\)th term of the following geometric sequence?}
{\( a_n =4\cdot3^{n+1}\)}{\( a_n =4\cdot3^{n}\)}{\( a_n =4\cdot3^{n-1}\)}{\( a_n =3^{n-1}\)}{}{}}
\LANGpl{
\Question{
Rysunek przedstawia fragment wykresu ciągu geometrycznego. Jaki jest dokładny wzór na \(n\)-ty wyraz tego ciągu geometrycznego?}
{\( a_n =4\cdot3^{n+1}\)}{\( a_n =4\cdot3^{n}\)}{\( a_n =4\cdot3^{n-1}\)}{\( a_n =3^{n-1}\)}{}{}}
\LANGcs{
\Question{
Na obrázku je část grafu geometrické posloupnosti. Jaký je vzorec pro \(n\)-tý člen této posloupnosti?}
{\( a_n =4\cdot3^{n+1}\)}{\( a_n =4\cdot3^{n}\)}{\( a_n =4\cdot3^{n-1}\)}{\( a_n =3^{n-1}\)}{}{}}
\LANGsk{
\Question{
Na obrázku je časť grafu geometrickej postupnosti. Vzorec pre \(n\)-tý člen tejto postupnosti je:}
{\( a_n =4\cdot3^{n+1}\)}{\( a_n =4\cdot3^{n}\)}{\( a_n =4\cdot3^{n-1}\)}{\( a_n =3^{n-1}\)}{}{}}
\LANGes{
\Question{
En el dibujo tenemos una parte de la gráfica de una progresión geométrica. ¿Cuál es la fórmula para calcular el término \(n\)-ésimo?}
{\( a_n =4\cdot3^{n+1}\)}{\( a_n =4\cdot3^{n}\)}{\( a_n =4\cdot3^{n-1}\)}{\( a_n =3^{n-1}\)}{}{}}
\End

\Data\ID{2010014003}
\Updated{2022-11-02 09:11:26}
\Level{A}
\SubArea{Geometric Sequences}
\LANGen{
\Question{
First five terms of four sequences are given. One of these sequences is not a geometric sequence. Choose this sequence.}
{\( \frac12 , \frac14, \frac16, \frac18, \frac1{10}\)}{\(1,-1,1,-1,1\)}{\(2,2,2,2,2\)}{\( 6,12,24,48,96\)}{}{}}
\LANGpl{
\Question{
Podano pięć pierwszych wyrazów czterech ciągów. Jeden z tych ciągów nie jest ciągiem geometrycznym. Wybierz ten ciąg.}
{\( \frac12 , \frac14, \frac16, \frac18, \frac1{10}\)}{\(1,-1,1,-1,1\)}{\(2,2,2,2,2\)}{\( 6,12,24,48,96\)}{}{}}
\LANGcs{
\Question{
Je dáno prvních pět členů čtyř posloupností. Určete, která z posloupností není geometrická.}
{\( \frac12 , \frac14, \frac16, \frac18, \frac1{10}\)}{\(1,-1,1,-1,1\)}{\(2,2,2,2,2\)}{\( 6,12,24,48,96\)}{}{}}
\LANGsk{
\Question{
Nižšie je uvedených prvých päť členov daných postupností. Jedna z nich nie je geometrická. Určte túto postupnosť.}
{\( \frac12 , \frac14, \frac16, \frac18, \frac1{10}\)}{\(1,-1,1,-1,1\)}{\(2,2,2,2,2\)}{\( 6,12,24,48,96\)}{}{}}
\LANGes{
\Question{
Se muestran los cinco primeros términos de varias sucesiones. Elige cuál de ellas no es una progresión geométrica.}
{\( \frac12 , \frac14, \frac16, \frac18, \frac1{10}\)}{\(1,-1,1,-1,1\)}{\(2,2,2,2,2\)}{\( 6,12,24,48,96\)}{}{}}
\End

\Data\ID{2010014004}
\Updated{2022-11-02 09:11:26}
\Level{A}
\SubArea{Geometric Sequences}
\LANGen{
\Question{
First five terms of four sequences are given. One of these sequences is not a geometric sequence. Choose this sequence.}
{\( \frac13 , \frac16, \frac19, \frac1{12}, \frac1{15}\)}{\(2,-2,2,-2,2\)}{\(7,7,7,7,7\)}{\( 1,4,4^2,4^3,4^4\)}{}{}}
\LANGpl{
\Question{
Podano pięć pierwszych wyrazów czterech ciągów. Jeden z tych ciągów nie jest ciągiem geometrycznym. Wybierz ten ciąg.}
{\( \frac13 , \frac16, \frac19, \frac1{12}, \frac1{15}\)}{\(2,-2,2,-2,2\)}{\(7,7,7,7,7\)}{\( 1,4,4^2,4^3,4^4\)}{}{}}
\LANGcs{
\Question{
Je dáno prvních pět členů čtyř posloupností. Určete, která z posloupností není geometrická.}
{\( \frac13 , \frac16, \frac19, \frac1{12}, \frac1{15}\)}{\(2,-2,2,-2,2\)}{\(7,7,7,7,7\)}{\( 1,4,4^2,4^3,4^4\)}{}{}}
\LANGsk{
\Question{
Nižšie je uvedených prvých päť členov daných postupností. Jedna z nich nie je geometrická. Určte túto postupnosť.}
{\( \frac13 , \frac16, \frac19, \frac1{12}, \frac1{15}\)}{\(2,-2,2,-2,2\)}{\(7,7,7,7,7\)}{\( 1,4,4^2,4^3,4^4\)}{}{}}
\LANGes{
\Question{
Se muestran los cinco primeros términos de varias sucesiones. Elige cuál de ellas no es una progresión geométrica.}
{\( \frac13 , \frac16, \frac19, \frac1{12}, \frac1{15}\)}{\(2,-2,2,-2,2\)}{\(7,7,7,7,7\)}{\( 1,4,4^2,4^3,4^4\)}{}{}}
\End

\Data\ID{2010014007}
\Updated{2022-11-02 09:11:26}
\Level{A}
\SubArea{Geometric Sequences}
\IMG{2010014001_1_0.pdf}
\LANGen{
\Question{
The picture shows first three patterns formed by green dots. The numbers of dots in patterns correspond to terms of a geometric sequence. To which term corresponds a pattern that will use \(192\) dots?}
{Seventh}{Eighth}{Tenth}{Nineth}{}{}}
\LANGpl{
\Question{
Zdjęcie przedstawia pierwsze trzy wzory utworzone przez zielone kropki. Liczby kropek we wzorach odpowiadają wyrazom ciągu geometrycznego. Któremu wyrazowi odpowiada wzór, który będzie używał kropek \(192\)?}
{siódmy}{ósmy}{dziesiąty}{dziewiąty}{}{}}
\LANGcs{
\Question{
Počet teček tvarů na obrázku reprezentuje členy geometrické posloupnosti, na obrázku jsou první tři členy. Kolikátý člen posloupnosti bude obsahovat \(192\) teček?}
{sedmý}{osmý}{desátý}{devátý}{}{}}
\LANGsk{
\Question{
Geometrická postupnosť je reprezentovaná počtom bodiek v daných útvaroch. Na obrázku sú prvé tri útvary. Ktorý útvar bude tvorený \(192\) bodkami?}
{siedmy}{ôsmy}{desiaty}{deviaty}{}{}}
\LANGes{
\Question{
La imagen muestra los tres primeros patrones formados por puntos verdes. El número de puntos en los patrones corresponde a los términos de una progresión geométrica. ¿A qué término corresponderá el patrón con \(192\) puntos?}
{Séptimo}{Octavo}{Décimo}{Noveno}{}{}}
\End

\Data\ID{2010014008}
\Updated{2022-11-02 09:11:26}
\Level{A}
\SubArea{Geometric Sequences}
\IMG{2010014001_2.pdf}
\LANGen{
\Question{
The picture shows first three patterns formed by green dots. The numbers of dots in patterns correspond to terms of a geometric sequence. To which term corresponds a pattern that will use \(640\) dots?}
{Eighth}{Seventh}{Tenth}{Nineth}{}{}}
\LANGpl{
\Question{
Zdjęcie przedstawia pierwsze trzy wzory utworzone przez zielone kropki. Liczby kropek we wzorach odpowiadają wyrazom ciągu geometrycznego. Któremu wyrazowi odpowiada wzór, który będzie używał kropek \(640\)?}
{ósmemu}{siódmemu}{dziesiątemu}{dziewiątemu}{}{}}
\LANGcs{
\Question{
Počet teček tvarů na obrázku reprezentuje členy geometrické posloupnosti, na obrázku jsou první tři členy. Kolikátý člen posloupnosti bude obsahovat \(640\) teček?}
{osmý}{sedmý}{desátý}{devátý}{}{}}
\LANGsk{
\Question{
Geometrická postupnosť je reprezentovaná počtom bodiek v daných útvaroch. Na obrázku sú prvé tri útvary. Ktorý útvar bude tvorený \(640\) bodkami?}
{ôsmy}{siedmy}{desiaty}{deviaty}{}{}}
\LANGes{
\Question{
La imagen muestra los tres primeros patrones formados por puntos verdes. El números de puntos en los patrones corresponde a los términos de una progresión geométrica. ¿A qué término corresponderá el patrón con \(640\) puntos?}
{Octavo}{Séptimo}{Décimo}{Noveno}{}{}}
\End

\Data\ID{2010014009}
\Updated{2022-11-02 09:11:26}
\Level{A}
\SubArea{Geometric Sequences}
\IMG{2010014001_3.pdf}
\LANGen{
\Question{
A sequence of patterns uses grey and purple squares. Identify the false statement about the numbers of squares in the patterns.}
{The numbers of purple squares in patterns form a geometric sequence with an even common ratio.}{The numbers of purple squares in patterns form a geometric sequence with a positive common ratio}{The numbers of grey squares in patterns form a geometric sequence with a positive common ratio.}{The numbers of grey squares in patterns form a geometric sequence with an even common ratio.}{}{}}
\LANGpl{
\Question{
Sekwencja wzorów wykorzystuje szare i fioletowe kwadraty. Zidentyfikuj fałsz stwierdzenie dotyczące liczby kwadratów we wzorach.}
{Liczby fioletowych kwadratów we wzorach tworzą ciąg geometryczny o równym, wspólnym ilorazie.}{Liczby fioletowych kwadratów we wzorach tworzą ciąg geometryczny o dodatnim wspólnym ilorazie.}{Liczby szarych kwadratów we wzorach tworzą ciąg geometryczny o dodatnim wspólnym ilorazie.}{Liczby szarych kwadratów we wzorach tworzą ciąg geometryczny o równym, wspólnym ilorazie.}{}{}}
\LANGcs{
\Question{
Mějme posloupnost útvarů tvořených šedými a fialovými čtverci. Vyberte nepravdivý výrok o počtu čtverců v jednotlivých útvarech.}
{Počet fialových čtverců tvoří geometrickou posloupnost se sudým kvocientem.}{Počet fialových čtverců tvoří geometrickou posloupnost s kladným kvocientem.}{Počet šedých čtverců tvoří geometrickou posloupnost s kladným kvocientem.}{Počet šedých čtverců tvoří geometrickou posloupnost se sudým kvocientem.}{}{}}
\LANGsk{
\Question{
Majme postupnosť útvarov, ktoré sú tvorené šedými a fialovými štvorcami. Vyberte nepravdivé tvrdenie o postupnosti počtu štvorcov v jednotlivých útvaroch.}
{Počet fialových štvorcov tvorí geometrickú postupnosť s párnym kvocientom.}{Počet fialových štvorcov tvorí geometrickú postupnosť s kladným kvocientom.}{Počet šedých štvorcov tvorí geometrickú postupnosť s kladným kvocientom.}{Počet šedých štvorcov tvorí geometrickú postupnosť s párnym kvocientom.}{}{}}
\LANGes{
\Question{
La siguiente sucesión de patrones utiliza cuadrados grises y morados. Identifica la afirmación falsa sobre el número de cuadrados en los patrones.}
{Los números de cuadrados morados en los patrones forman una progresión geométrica cuya razón es par.}{Los números de cuadrados morados en los patrones forman una progresión geométrica cuya razón es positiva.}{Los números de cuadrados grises en los patrones forman una progresión geométrica cuya razón es positiva.}{Los números de cuadrados grises en los patrones forman una progresión geométrica cuya razón es par.}{}{}}
\End

\Data\ID{2010014010}
\Updated{2022-11-02 09:11:26}
\Level{A}
\SubArea{Geometric Sequences}
\IMG{2010014001_4.pdf}
\LANGen{
\Question{
A sequence of patterns uses grey and purple squares. Identify the false statement about the numbers of squares in the patterns.}
{The numbers of grey squares in patterns do not form a geometric sequence.}{The numbers of purple squares in patterns form a geometric sequence with an even common ratio.}{The numbers of purple squares in patterns form a geometric sequence with a positive common ratio.}{The numbers of grey squares in patterns form a geometric sequence with an even common ratio.}{}{}}
\LANGpl{
\Question{
Sekwencja wzorów wykorzystuje szare i fioletowe kwadraty. Zidentyfikuj fałsz stwierdzenie dotyczące liczby kwadratów we wzorach.}
{Liczby szarych kwadratów we wzorach nie tworzą ciągu geometrycznego.}{Liczby fioletowych kwadratów we wzorach tworzą ciąg geometryczny o równym, wspólnym ilorazie.}{Liczby fioletowych kwadratów we wzorach tworzą ciąg geometryczny o dodatnim wspólnym ilorazie.}{Liczby szarych kwadratów we wzorach tworzą ciąg geometryczny o równym, wspólnym ilorazie.}{}{}}
\LANGcs{
\Question{
Mějme posloupnost útvarů tvořených šedými a fialovými čtverci. Vyberte nepravdivý výrok o počtu čtverců v jednotlivých útvarech.}
{Počet šedých čtverců netvoří geometrickou posloupnost.}{Počet fialových čtverců tvoří geometrickou posloupnost se sudým kvocientem.}{Počet fialových čtverců tvoří geometrickou posloupnost s kladným kvocientem.}{Počet šedých čtverců tvoří geometrickou posloupnost se sudým kvocientem.}{}{}}
\LANGsk{
\Question{
Majme postupnosť útvarov, ktoré sú tvorené šedými a fialovými štvorcami. Vyberte nepravdivé tvrdenie o postupnosti počtu štvorcov v jednotlivých útvaroch.}
{Počet šedých štvorcov netvorí geometrickú postupnosť.}{Počet fialových štvorcov tvorí geometrickú postupnosť s párnym kvocientom.}{Počet fialových štvorcov tvorí geometrickú postupnosť s kladným kvocientom.}{Počet šedých štvorcov tvorí geometrickú postupnosť s párnym kvocientom.}{}{}}
\LANGes{
\Question{
La siguiente sucesión de patrones utiliza cuadrados grises y morados. Identifica la afirmación falsa sobre el número de cuadrados en los patrones.}
{Los números de cuadrados grises en los patrones no forman una progresión geométrica.}{Los números de cuadrados morados en los patrones forman una progresión geométrica cuya razón es par.}{Los números de cuadrados morados en los patrones forman una progresión geométrica cuya razón es positiva.}{Los números de cuadrados grises en los patrones forman una progresión geométrica cuya razón es par.}{}{}}
\End

\Data\ID{2110014001}
\Updated{2022-11-02 09:11:56}
\Level{A}
\SubArea{Geometric Sequences}
\LANGen{
\Question{
In one of the pictures given below is \textbf{not} a graph of a part of a geometric sequence. Choose this graph.}
{}{}{}{}{}{}}
\LANGpl{
\Question{
Jednen z poniższych obrazków \textbf{nie} jest wykresem części ciągu geometrycznego. Wybierz ten wykres.}
{}{}{}{}{}{}}
\LANGcs{
\Question{
Jeden z grafů \textbf{neznázorňuje} část geometrické posloupnosti Vyberte tento graf.}
{}{}{}{}{}{}}
\LANGsk{
\Question{
Na jednom z obrázkov \textbf{nie} je graf časti geometrickej postupnosti. Vyberte, ktorý to je.}
{}{}{}{}{}{}}
\LANGes{
\Question{
En una de las siguientes imágenes \textbf{hay} una gráfica que \textbf{no} es parte de una progresión geométrica. Elige dicha gráfica.}
{}{}{}{}{}{}}

\IMGanswer{2010014001-figure0.pdf}
\IMGanswer{2010014001-figure1_0.pdf}
\IMGanswer{2010014001-figure2_0.pdf}
\IMGanswer{2010014001-figure3_0.pdf}\End

\Data\ID{2110014002}
\Updated{2022-11-02 09:11:00}
\Level{A}
\SubArea{Geometric Sequences}
\LANGen{
\Question{
In one of the pictures given below is \textbf{not} a graph of a part of a geometric sequence. Choose this graph.}
{}{}{}{}{}{}}
\LANGpl{
\Question{
Jeden z poniższych obrazków \textbf{nie} jest wykresem części ciągu geometrycznego. Wybierz ten wykres.}
{}{}{}{}{}{}}
\LANGcs{
\Question{
Jeden z grafů \textbf{neznázorňuje} část geometrické posloupnosti Vyberte tento graf.}
{}{}{}{}{}{}}
\LANGsk{
\Question{
Na jednom z obrázkov \textbf{nie} je graf časti geometrickej postupnosti. Vyberte, ktorý to je.}
{}{}{}{}{}{}}
\LANGes{
\Question{
En una de las siguientes imágenes \textbf{hay} una gráfica que \textbf{no} es parte de una progresión geométrica. Elige dicha gráfica.}
{}{}{}{}{}{}}

\IMGanswer{2010014001-figure4_0.pdf}
\IMGanswer{2010014001-figure5.pdf}
\IMGanswer{2010014001-figure6_0.pdf}
\IMGanswer{2010014001-figure7_0.pdf}\End

\Data\ID{1003112804}
\Updated{2022-01-06 01:01:55}
\Level{B}
\SubArea{Geometric Sequences}
\LANGen{
\Question{
The third term of a geometric sequence is \( -5 \) and the eighth term is \( -5 \). \( s_5 \) is the sum of the first five terms and \( q \) is the common ratio of this sequence. Choose the formula which is not valid for this sequence:}
{\( s_5=-5\cdot\frac{q^5-1}{q-1} \)}{\( s_5=-25 \)}{\( s_5=5\cdot a_1 \)}{\( s_5=5\cdot a_3 \)}{\( s_5=5\cdot(-5) \)}{}}
\LANGpl{
\Question{
Trzeci wyraz ciągu wynosi \( -5 \), a ósmy wyraz to \( -5 \). \( s_5 \) jest sumą pięciu pierwszych wyrazów, a \( q \) jest wspólnym współczynnikiem ciągu. Wybierz wzór, który nie jest poprawny dla tego ciągu:}
{\( s_5=-5\cdot\frac{q^5-1}{q-1} \)}{\( s_5=-25 \)}{\( s_5=5\cdot a_1 \)}{\( s_5=5\cdot a_3 \)}{\( s_5=5\cdot(-5) \)}{}}
\LANGcs{
\Question{
Třetí člen geometrické posloupnosti je roven \( -5 \) a osmý člen je \( -5 \). \( s_5 \) je součet prvních pěti členů a \( q \) je kvocient této posloupnosti. Vyberte tvrzení, které neplatí v této posloupnosti.}
{\( s_5=-5\cdot\frac{q^5-1}{q-1} \)}{\( s_5=-25 \)}{\( s_5=5\cdot a_1 \)}{\( s_5=5\cdot a_3 \)}{\( s_5=5\cdot(-5) \)}{}}
\LANGsk{
\Question{
Tretí člen geometrické postupnosti je rovný \( -5 \) a ôsmy člen je \( -5 \). \( s_5 \) je súčet prvých piatich členov a \( q \) je kvocient tejto postupnosti. Vyberte tvrdenie, ktoré neplatí v tejto postupnosti.}
{\( s_5=-5\cdot\frac{q^5-1}{q-1} \)}{\( s_5=-25 \)}{\( s_5=5\cdot a_1 \)}{\( s_5=5\cdot a_3 \)}{\( s_5=5\cdot(-5) \)}{}}
\LANGes{
\Question{
El tercer término de una progresión geométrica es \( -5 \) y su término octavo es \( -5 \). \( s_5 \) es la suma de los cinco primeros términos y  \( q \) es la razón. Elige la declaración falsa.}
{\( s_5=-5\cdot\frac{q^5-1}{q-1} \)}{\( s_5=-25 \)}{\( s_5=5\cdot a_1 \)}{\( s_5=5\cdot a_3 \)}{\( s_5=5\cdot(-5) \)}{}}
\End

\Data\ID{1003112806}
\Updated{2022-01-06 01:01:02}
\Level{B}
\SubArea{Geometric Sequences}
\LANGen{
\Question{
The sum of the first four terms of a geometric sequence is \( 0 \) and the first term equals \( 2 \). Choose the correct formula to find the eighth term.}
{\( a_8 = 2\cdot (-1)^7 \)}{\( a_8 = 2\cdot (1)^7 \)}{\( a_8 = 2\cdot 2 \)}{\( a_8 = \frac02 \)}{\( a_8 = 2\cdot (-2) \)}{}}
\LANGpl{
\Question{
Suma czterech pierwszych wyrazów ciągu geometrycznego wynosi \( 0 \), a pierwszy wyraz równy jest \( 2 \). Wybierz poprawny wzór, aby wyznaczyć ósmy wyraz ciągu.}
{\( a_8 = 2\cdot (-1)^7 \)}{\( a_8 = 2\cdot (1)^7 \)}{\( a_8 = 2\cdot 2 \)}{\( a_8 = \frac02 \)}{\( a_8 = 2\cdot (-2) \)}{}}
\LANGcs{
\Question{
Součet prvních čtyř členů geometrické posloupnosti je \( 0 \) a první člen je roven \( 2 \). Pro osmý člen této posloupnosti platí:}
{\( a_8 = 2\cdot (-1)^7 \)}{\( a_8 = 2\cdot (1)^7 \)}{\( a_8 = 2\cdot 2 \)}{\( a_8 = \frac02 \)}{\( a_8 = 2\cdot (-2) \)}{}}
\LANGsk{
\Question{
Súčet prvých štyroch členov geometrickej postupnosti je \( 0 \) a prvý člen je rovný \( 2 \). Pre ôsmy člen tejto postupnosti platí:}
{\( a_8 = 2\cdot (-1)^7 \)}{\( a_8 = 2\cdot (1)^7 \)}{\( a_8 = 2\cdot 2 \)}{\( a_8 = \frac02 \)}{\( a_8 = 2\cdot (-2) \)}{}}
\LANGes{
\Question{
La suma de los cuatro primeros términos de una progresión geométrica es \( 0 \) y su primer término es igual a \( 2 \). Para el octavo término de la progresión se cumple:}
{\( a_8 = 2\cdot (-1)^7 \)}{\( a_8 = 2\cdot (1)^7 \)}{\( a_8 = 2\cdot 2 \)}{\( a_8 = \frac02 \)}{\( a_8 = 2\cdot (-2) \)}{}}
\End

\Data\ID{1003112807}
\Updated{2022-04-23 05:04:20}
\Level{B}
\SubArea{Geometric Sequences}
\LANGen{
\Question{
The sum of the first two terms of a geometric sequence is \( 28 \) and its first term equals \( 2 \). Which of the following statements about its common ratio \( q \) is not valid?}
{\( q \) is an even number.}{\( q > 10 \)}{\( q < 28 \)}{\( q \) is prime number.}{\( q \) is a divisor of \( 26 \).}{}}
\LANGpl{
\Question{
Suma dwóch pierwszych wyrazów ciągu geometrycznego wynosi \( 28 \), a jego pierwszy wyraz jest równy  \( 2 \). Które z poniższych stwierdzeń dotyczących wspólnego współczynnika  \( q \) nie jest prawdziwe?}
{\( q \) jest liczbą parzystą.}{\( q > 10 \)}{\( q < 28 \)}{\( q \) jest liczbą pierwszą.}{\( q \) jest dzielnikiem \( 26 \).}{}}
\LANGcs{
\Question{
Součet prvních dvou členů geometrické posloupnosti je \( 28 \) a první člen je roven \( 2 \). Pro kvocient této posloupnosti neplatí:}
{\( q \) je sudé číslo.}{\( q > 10 \)}{\( q < 28 \)}{\( q \) je prvočíslo.}{\( q \) je dělitel \( 26 \).}{}}
\LANGsk{
\Question{
Súčet prvých dvoch členov geometrickej postupnosti je \( 28 \) a prvý člen je rovný \( 2 \). Pre kvocient tejto postupnosti neplatí:}
{\( q \) je párne číslo.}{\( q > 10 \)}{\( q < 28 \)}{\( q \) je prvočíslo.}{\( q \) je deliteľ \( 26 \).}{}}
\LANGes{
\Question{
La suma de loa primeros dos términos de una sucesión geométrica es \( 28 \) y su primer término es \( 2 \). Para la razón de esta sucesión no se cumple que:}
{\( q \) es par}{\( q > 10 \)}{\( q < 28 \)}{\( q \) es primo}{\( q \) es divisor de \( 26 \).}{}}
\End

\Data\ID{1003112808}
\Updated{2022-01-06 01:01:34}
\Level{B}
\SubArea{Geometric Sequences}
\LANGen{
\Question{
The first term of a geometric sequence is \( 15 \) and the common ratio is \( -1 \). Find the sum of the first \( 10 \) terms of this sequence.}
{\( 0 \)}{\( 15 \)}{\( -15 \)}{\( 150 \)}{\( -150 \)}{}}
\LANGpl{
\Question{
Pierwszy wyraz ciągu geometrycznego wynosi \( 15 \), a  \( -1 \). Wyznacz sumę pierwszych \( 10 \) wyrazów tego ciągu.}
{\( 0 \)}{\( 15 \)}{\( -15 \)}{\( 150 \)}{\( -150 \)}{}}
\LANGcs{
\Question{
První člen geometrické posloupnosti je roven \( 15 \) a kvocient je \( -1 \). Určete součet prvních deseti členů této posloupnosti.}
{\( 0 \)}{\( 15 \)}{\( -15 \)}{\( 150 \)}{\( -150 \)}{}}
\LANGsk{
\Question{
Prvý člen geometrickej postupnosti je rovný \( 15 \) a kvocient je \( -1 \). Určte súčet prvých desiatich členov tejto postupnosti.}
{\( 0 \)}{\( 15 \)}{\( -15 \)}{\( 150 \)}{\( -150 \)}{}}
\LANGes{
\Question{
El primer término de una progresión geométrica es igual a \( 15 \) y su razón es \( -1 \). Calcula la suma de los primeros diez términos de la progresión.}
{\( 0 \)}{\( 15 \)}{\( -15 \)}{\( 150 \)}{\( -150 \)}{}}
\End

\Data\ID{1003112810}
\Updated{2022-04-23 05:04:20}
\Level{B}
\SubArea{Geometric Sequences}
\LANGen{
\Question{
The sum of the first \( n \) terms of a geometric sequence is \( -5 \), the common ratio is \( -2 \) and the first term equals \( 1 \). Find \( n \).}
{\( 4 \)}{\( 3 \)}{\( 5 \)}{\( 6 \)}{\( 2 \)}{}}
\LANGpl{
\Question{
Suma pierwszych \( n \) wyrazów ciągu geometrycznego wynosi \( -5 \), a wspólny współczynnik \( -2 \), zaś pierwszy wyraz ciągu jest równy \( 1 \). Wyznacz \( n \).}
{\( 4 \)}{\( 3 \)}{\( 5 \)}{\( 6 \)}{\( 2 \)}{}}
\LANGcs{
\Question{
Součet prvních \( n \) členů geometrické posloupnosti je roven \( -5 \), kvocient je \( -2 \) a první člen je \( 1 \). Určete \( n \).}
{\( 4 \)}{\( 3 \)}{\( 5 \)}{\( 6 \)}{\( 2 \)}{}}
\LANGsk{
\Question{
Súčet prvých \( n \) členov geometrickej postupnosti je rovný \( -5 \), kvocient je \( -2 \) a prvý člen je \( 1 \). Určte \( n \).}
{\( 4 \)}{\( 3 \)}{\( 5 \)}{\( 6 \)}{\( 2 \)}{}}
\LANGes{
\Question{
La suma de los primeros \( n \) términos de una progresión geométrica es igual a \( -5 \), la razón es \( -2 \) y su primer término es \( 1 \). Averigua \( n \).}
{\( 4 \)}{\( 3 \)}{\( 5 \)}{\( 6 \)}{\( 2 \)}{}}
\End

\Data\ID{1003112812}
\Updated{2022-01-06 01:01:40}
\Level{B}
\SubArea{Geometric Sequences}
\LANGen{
\Question{
The sum \( s_3 \) of the first three terms of a geometric sequence is \( 21 \) and its common ratio is \( -2 \). Which of the following (in)equalities is true for the first term of the sequence?}
{\( a_1 > 0 \)}{\( a_1 < 0 \)}{\( a_1 = 0 \)}{\( a_1 < a_2 \)}{\( a_1 > s_3 \)}{}}
\LANGpl{
\Question{
Suma \( s_3 \) trzech pierwszych wyrazów ciągu geometrycznego wynosi  \( 21 \), a jego wspólny współczynnik jest równy \( -2 \). Która z poniższych (nie)równości jest prawdziwa dla pierwszego wyrazu ciągu?}
{\( a_1 > 0 \)}{\( a_1 < 0 \)}{\( a_1 = 0 \)}{\( a_1 < a_2 \)}{\( a_1 > s_3 \)}{}}
\LANGcs{
\Question{
Součet \( s_3 \) prvních tří členů geometrické posloupnosti je roven \( 21 \) a kvocient je \( -2 \). Pro první člen této posloupnosti platí:}
{\( a_1 > 0 \)}{\( a_1 < 0 \)}{\( a_1 = 0 \)}{\( a_1 < a_2 \)}{\( a_1 > s_3 \)}{}}
\LANGsk{
\Question{
Súčet \( s_3 \) prvých troch členov geometrickej postupnosti je rovný \( 21 \) a kvocient je \( -2 \). Pre prvý člen tejto postupnosti platí:}
{\( a_1 > 0 \)}{\( a_1 < 0 \)}{\( a_1 = 0 \)}{\( a_1 < a_2 \)}{\( a_1 > s_3 \)}{}}
\LANGes{
\Question{
La suma \( s_3 \) de los tres primeros términos de una progresión geométrica es igual a \( 21 \) y su razón es \( -2 \). Para el primer término se cumple que:}
{\( a_1 > 0 \)}{\( a_1 < 0 \)}{\( a_1 = 0 \)}{\( a_1 < a_2 \)}{\( a_1 > s_3 \)}{}}
\End

\Data\ID{1003134602}
\Updated{2022-01-10 08:01:53}
\Level{B}
\SubArea{Geometric Sequences}
\LANGen{
\Question{
The sum of the second and the third term of a geometric sequence is $6\log3$ and the common ratio is $2$. Find the sum of the first three terms.}
{$7\log3$}{$2\log3^6$}{$2\log6$}{$\log9$}{$3\log6$}{}}
\LANGpl{
\Question{
Suma drugiego i trzeciego wyrazu ciągu jest równa  $6\log3$, a wspólny współczynnik wynosi $2$. Wyznacz sumę trzech pierwszych wyrazów ciągu.}
{$7\log3$}{$2\log3^6$}{$2\log6$}{$\log9$}{$3\log6$}{}}
\LANGcs{
\Question{
Určete součet prvních tří členů geometrické posloupnosti, je-li součet druhého a třetího členu $6\log3$ a kvocient je $2$.}
{$7\log3$}{$2\log3^6$}{$2\log6$}{$\log9$}{$3\log6$}{}}
\LANGsk{
\Question{
Určte súčet prvých troch členov geometrickej postupnosti, ak je súčet druhého a tretieho člena $6\log3$ a kvocient je $2$.}
{$7\log3$}{$2\log3^6$}{$2\log6$}{$\log9$}{$3\log6$}{}}
\LANGes{
\Question{
Averigua la suma de primeros tres términos de una progresión geométrica, si la suma del segundo y tercer término es  $6\log3$ y su razón es $2$.}
{$7\log3$}{$2\log3^6$}{$2\log6$}{$\log9$}{$3\log6$}{}}
\End

\Data\ID{1003134603}
\Updated{2022-01-13 02:01:11}
\Level{B}
\SubArea{Geometric Sequences}
\LANGen{
\Question{
The sum of the first five terms of a geometric sequence is less than $1$ and the common ratio is $10$. Determine all possible values of the first term.}
{$ a_1 < \frac1{11111}$, $a_1\in\mathbb{R}$}{$ -\frac1{10^5} < a_1 < \frac1{10^5}$, $a_1\in\mathbb{R}$}{$a_1 < \frac1{99999}$, $a_1\in\mathbb{R}$}{$a_1 < 10^{-4}$, $a_1\in\mathbb{R}$}{$a_1 < 10^{-5}$, $a_1\in\mathbb{R}$}{}}
\LANGpl{
\Question{
Suma pięciu pierwszych wyrazów ciągu geometrycznego jest mniejsza niż $1$, a wspólny współczynnik wynosi  $10$. Wyznacz wszystkie możliwe wartości pierwszego wyrazu ciągu.}
{$ a_1 < \frac1{11111}$, $a_1\in\mathbb{R}$}{$ -\frac1{10^5} < a_1 < \frac1{10^5}$, $a_1\in\mathbb{R}$}{$a_1 < \frac1{99999}$, $a_1\in\mathbb{R}$}{$a_1 < 10^{-4}$, $a_1\in\mathbb{R}$}{$a_1 < 10^{-5}$, $a_1\in\mathbb{R}$}{}}
\LANGcs{
\Question{
Součet prvních pěti členů geometrické posloupnosti je menší než $1$ a kvocient je $10$. Najděte všechny možné hodnoty pro první člen.}
{$ a_1 < \frac1{11111}$, $a_1\in\mathbb{R}$}{$ -\frac1{10^5} < a_1 < \frac1{10^5}$, $a_1\in\mathbb{R}$}{$a_1 < \frac1{99999}$, $a_1\in\mathbb{R}$}{$a_1 < 10^{-4}$, $a_1\in\mathbb{R}$}{$a_1 < 10^{-5}$, $a_1\in\mathbb{R}$}{}}
\LANGsk{
\Question{
Súčet prvých piatich členov geometrickej postupnosti je menší než $1$ a kvocient je $10$. Nájdite všetky možné hodnoty pre prvý člen.}
{$ a_1 < \frac1{11111}$, $a_1\in\mathbb{R}$}{$ -\frac1{10^5} < a_1 < \frac1{10^5}$, $a_1\in\mathbb{R}$}{$a_1 < \frac1{99999}$, $a_1\in\mathbb{R}$}{$a_1 < 10^{-4}$, $a_1\in\mathbb{R}$}{$a_1 < 10^{-5}$, $a_1\in\mathbb{R}$}{}}
\LANGes{
\Question{
La suma de los primeros cinco términos de una progresión geométrica es menor que $1$ y su razón es $10$. Encuentra todos los valores posibles para el primer término.}
{$ a_1 < \frac1{11111}$, $a_1\in\mathbb{R}$}{$ -\frac1{10^5} < a_1 < \frac1{10^5}$, $a_1\in\mathbb{R}$}{$a_1 < \frac1{99999}$, $a_1\in\mathbb{R}$}{$a_1 < 10^{-4}$, $a_1\in\mathbb{R}$}{$a_1 < 10^{-5}$, $a_1\in\mathbb{R}$}{}}
\End

\Data\ID{1003134604}
\Updated{2022-01-13 02:01:24}
\Level{B}
\SubArea{Geometric Sequences}
\LANGen{
\Question{
Find the $5$th term of a geometric sequence $\{a_n\}_{n=1}^{\infty}$, where $a_1\cdot a_3\cdot a_5=x^3$, and $\frac{a_4}{a_2} =y^4$.}
{$xy^4$}{$\frac x{y^4}$}{$\frac x{y^2}$}{$x$}{$xy^2$}{}}
\LANGpl{
\Question{
Wyznacz $5$-ty wyraz ciągu geometrycznego $\{a_n\}_{n=1}^{\infty}$, gdzie $a_1\cdot a_3\cdot a_5=x^3$, i $\frac{a_4}{a_2} =y^4$.}
{$xy^4$}{$\frac x{y^4}$}{$\frac x{y^2}$}{$x$}{$xy^2$}{}}
\LANGcs{
\Question{
Pátý člen geometrické posloupnosti, ve které platí $a_1\cdot a_3\cdot a_5=x^3$ a $\frac{a_4}{a_2} =y^4$, je:}
{$xy^4$}{$\frac x{y^4}$}{$\frac x{y^2}$}{$x$}{$xy^2$}{}}
\LANGsk{
\Question{
Piaty člen geometrickej postupnosti, v ktorej platí $a_1\cdot a_3\cdot a_5=x^3$ a $\frac{a_4}{a_2} =y^4$, je:}
{$xy^4$}{$\frac x{y^4}$}{$\frac x{y^2}$}{$x$}{$xy^2$}{}}
\LANGes{
\Question{
El quinto término de una progresión geométrica de la cuál sabemos que $a_1\cdot a_3\cdot a_5=x^3$ y $\frac{a_4}{a_2} =y^4$, es:}
{$xy^4$}{$\frac x{y^4}$}{$\frac x{y^2}$}{$x$}{$xy^2$}{}}
\End

\Data\ID{1003134701}
\Updated{2022-04-23 05:04:20}
\Level{B}
\SubArea{Geometric Sequences}
\LANGen{
\Question{
The sum of the first and the second term of a geometric sequence is \( 1 \), the sum of the third and the fourth term is \( 4 \) and the common ratio is negative. Find the first term.}
{\( -1 \)}{\( 2 \)}{\( -2 \)}{\( 3 \)}{\( \frac13 \)}{}}
\LANGpl{
\Question{
Suma pierwszego i drugiego wyrazu ciągu geometrycznego wynosi \( 1 \), suma trzeciego i czwartego wyrazu jest równa \( 4 \) a wspólny współczynnik jest ujemny. Wyznacz pierwszy wyraz tego ciągu.}
{\( -1 \)}{\( 2 \)}{\( -2 \)}{\( 3 \)}{\( \frac13 \)}{}}
\LANGcs{
\Question{
Určete první člen geometrické posloupnosti, jestliže součet prvního a druhého členu je \( 1 \), součet třetího a čtvrtého je \( 4 \) a kvocient je záporný.}
{\( -1 \)}{\( 2 \)}{\( -2 \)}{\( 3 \)}{\( \frac13 \)}{}}
\LANGsk{
\Question{
Určte prvý člen geometrickej postupnosti, ak súčet prvého a druhého členu je \( 1 \), súčet tretieho a štvrtého je \( 4 \) a kvocient je záporný.}
{\( -1 \)}{\( 2 \)}{\( -2 \)}{\( 3 \)}{\( \frac13 \)}{}}
\LANGes{
\Question{
Halla el primer término de una progresión geométrica si la suma de su primer y segundo término es \( 1 \), la suma de su tercer y cuarto término es \( 4 \) y la razón es negativa.}
{\( -1 \)}{\( 2 \)}{\( -2 \)}{\( 3 \)}{\( \frac13 \)}{}}
\End

\Data\ID{1003134702}
\Updated{2022-04-23 05:04:20}
\Level{B}
\SubArea{Geometric Sequences}
\LANGen{
\Question{
The sum of the first three terms of a geometric sequence is \( \frac78 \) and the common ratio is \( \frac12 \). Find the sum of all the terms from the \( 3 \)rd to the \( 5 \)th of the sequence.}
{\( \frac7{32} \)}{\( \frac{31}{32} \)}{\( \frac3{32} \)}{\( \frac78 \)}{\( \frac58 \)}{}}
\LANGpl{
\Question{
Suma trzech pierwszych wyrazów ciągu geometrycznego  jest równa \( \frac78 \), a wspólny współczynnik wynosi \( \frac12 \). Wyznacz sumę wszystkich wyrazów tego ciągu  od \( 3 \)-go do  \( 5 \)-go.}
{\( \frac7{32} \)}{\( \frac{31}{32} \)}{\( \frac3{32} \)}{\( \frac78 \)}{\( \frac58 \)}{}}
\LANGcs{
\Question{
Určete součet třetího až pátého členu geometrické posloupnosti, jestliže součet prvních tří členů je \( \frac78 \) a kvocient je \( \frac12 \).}
{\( \frac7{32} \)}{\( \frac{31}{32} \)}{\( \frac3{32} \)}{\( \frac78 \)}{\( \frac58 \)}{}}
\LANGsk{
\Question{
Určte súčet tretieho až piateho člena geometrickej postupnosti, ak súčet prvých troch členov je \( \frac78 \) a kvocient je \( \frac12 \).}
{\( \frac7{32} \)}{\( \frac{31}{32} \)}{\( \frac3{32} \)}{\( \frac78 \)}{\( \frac58 \)}{}}
\LANGes{
\Question{
Averigua la suma desde el tercer al quinto término de una progresión geométrica si la suma de los primeros tres términos es \( \frac78 \) y su razón es \( \frac12 \).}
{\( \frac7{32} \)}{\( \frac{31}{32} \)}{\( \frac3{32} \)}{\( \frac78 \)}{\( \frac58 \)}{}}
\End

\Data\ID{1003134703}
\Updated{2022-01-13 02:01:21}
\Level{B}
\SubArea{Geometric Sequences}
\LANGen{
\Question{
The first term of a geometric sequence is different from zero and \( s_3=-3s_2 \). (\( s_n \) is the sum of the first n terms.) Find the common ratio.}
{\( -2 \)}{\( 2 \)}{\( \frac12 \)}{\( 0 \)}{\( -3 \)}{}}
\LANGpl{
\Question{
Pierwszy wyraz ciągu geometrycznego jest różny od zera i   \( s_3=-3s_2 \). (\( s_n \) jest sumą pierwszych już n wyrazów.)  Wyznacz wspólny współczynnik.}
{\( -2 \)}{\( 2 \)}{\( \frac12 \)}{\( 0 \)}{\( -3 \)}{}}
\LANGcs{
\Question{
Určete kvocient geometrické posloupnosti, jestliže platí: \( s_3=-3s_2 \) a první člen je nenulový.}
{\( -2 \)}{\( 2 \)}{\( \frac12 \)}{\( 0 \)}{\( -3 \)}{}}
\LANGsk{
\Question{
Určte kvocient geometrickej postupnosti, ak platí: \( s_3=-3s_2 \) a prvý člen je nenulový.}
{\( -2 \)}{\( 2 \)}{\( \frac12 \)}{\( 0 \)}{\( -3 \)}{}}
\LANGes{
\Question{
Averigua la razón de una progresión geométrica si sabemos que:  \( s_3 =-3s_2 \) y el primer término no es cero.}
{\( -2 \)}{\( 2 \)}{\( \frac12 \)}{\( 0 \)}{\( -3 \)}{}}
\End

\Data\ID{1003134704}
\Updated{2022-01-13 02:01:01}
\Level{B}
\SubArea{Geometric Sequences}
\LANGen{
\Question{
The first term of a geometric sequence is \( 3 \). Find the common ratio so that the sum of the first \( 3 \) terms is minimal.}
{\( -\frac12 \)}{\( -5 \)}{\( 0 \)}{\( 2 \)}{\( -1 \)}{}}
\LANGpl{
\Question{
Pierwszy wyraz ciągu geometrycznego wynosi \( 3 \). Wyznacz wspólny współczynnik, taki dla którego suma trzech pierwszych wyrazów ciągu jest minimalna.}
{\( -\frac12 \)}{\( -5 \)}{\( 0 \)}{\( 2 \)}{\( -1 \)}{}}
\LANGcs{
\Question{
Určete kvocient geometrické posloupnosti, jejíž první člen je \( 3 \), tak, aby součet prvních tří členů byl minimální.}
{\( -\frac12 \)}{\( -5 \)}{\( 0 \)}{\( 2 \)}{\( -1 \)}{}}
\LANGsk{
\Question{
Určte kvocient geometrickej postupnosti, ktorej prvý člen je \( 3 \), tak, aby súčet prvých troch členov bol minimálny.}
{\( -\frac12 \)}{\( -5 \)}{\( 0 \)}{\( 2 \)}{\( -1 \)}{}}
\LANGes{
\Question{
Elije la razón que hace que la suma de los tres primeros términos de una progresión geométrica sea mínima si su primer término es \( 3 \).}
{\( -\frac12 \)}{\( -5 \)}{\( 0 \)}{\( 2 \)}{\( -1 \)}{}}
\End

\Data\ID{1003134705}
\Updated{2022-04-23 05:04:20}
\Level{B}
\SubArea{Geometric Sequences}
\LANGen{
\Question{
Find the first term of the geometric sequence \( \{a_n\}_{n=1}^{\infty} \), if: \[\begin{aligned} a_3\cdot a_4&=-25, \\ a_4-a_3&=10. \end{aligned} \]}
{\( -5 \)}{\( 5 \)}{\( -1 \)}{\( 1 \)}{\( -15 \)}{}}
\LANGpl{
\Question{
Wyznacz pierwszy wyraz ciągu geometrycznego \( \{a_n\}_{n=1}^{\infty} \), jeśli: \[\begin{aligned} a_3\cdot a_4&=-25, \\ a_4-a_3&=10. \end{aligned} \]}
{\( -5 \)}{\( 5 \)}{\( -1 \)}{\( 1 \)}{\( -15 \)}{}}
\LANGcs{
\Question{
Určete první člen geometrické posloupnosti \( \{a_n\}_{n=1}^{\infty} \), platí-li: \[\begin{aligned} a_3\cdot a_4&=-25, \\ a_4-a_3&=10. \end{aligned} \]}
{\( -5 \)}{\( 5 \)}{\( -1 \)}{\( 1 \)}{\( -15 \)}{}}
\LANGsk{
\Question{
Určte prvý člen geometrickej postupnosti \( \{a_n\}_{n=1}^{\infty} \), ak platí: \[\begin{aligned} a_3\cdot a_4&=-25, \\ a_4-a_3&=10. \end{aligned} \]}
{\( -5 \)}{\( 5 \)}{\( -1 \)}{\( 1 \)}{\( -15 \)}{}}
\LANGes{
\Question{
Averigua el primer término de una progresión geométrica \( \{a_n\}_{n=1}^{\infty} \), si sabemos que: \[\begin{aligned} a_3\cdot a_4&=-25, \\ a_4-a_3&=10. \end{aligned} \]}
{\( -5 \)}{\( 5 \)}{\( -1 \)}{\( 1 \)}{\( -15 \)}{}}
\End

\Data\ID{1003134706}
\Updated{2022-01-13 02:01:36}
\Level{B}
\SubArea{Geometric Sequences}
\LANGen{
\Question{
Find the second term and the common ratio of the geometric sequence \( \{a_n\}_{n=1}^{\infty} \), if: \[ \begin{aligned} a_2-a_1&=22, \\ a_3-a_2&=66. \end{aligned} \]}
{\( a_2=33 \), \( q=3 \)}{\( a_2=11 \), \( q=3 \)}{\( a_2=22 \), \( q=3 \)}{\( a_2=33 \), \( q=2 \)}{\( a_2=11 \), \( q=2 \)}{}}
\LANGpl{
\Question{
Wyznacz drugi wyraz ciągu i wspólny współczynnik ciągu geometrycznego \( \{a_n\}_{n=1}^{\infty} \), jeśli: \[ \begin{aligned} a_2-a_1&=22, \\ a_3-a_2&=66. \end{aligned} \]}
{\( a_2=33 \), \( q=3 \)}{\( a_2=11 \), \( q=3 \)}{\( a_2=22 \), \( q=3 \)}{\( a_2=33 \), \( q=2 \)}{\( a_2=11 \), \( q=2 \)}{}}
\LANGcs{
\Question{
Určete druhý člen a kvocient geometrické posloupnosti \( \{a_n\}_{n=1}^{\infty} \), platí-li: \[ \begin{aligned} a_2-a_1&=22, \\ a_3-a_2&=66. \end{aligned} \]}
{\( a_2=33 \), \( q=3 \)}{\( a_2=11 \), \( q=3 \)}{\( a_2=22 \), \( q=3 \)}{\( a_2=33 \), \( q=2 \)}{\( a_2=11 \), \( q=2 \)}{}}
\LANGsk{
\Question{
Určte druhý člen a kvocient geometrickej postupnosti \( \{a_n\}_{n=1}^{\infty} \), platí-li: \[ \begin{aligned} a_2-a_1&=22, \\ a_3-a_2&=66. \end{aligned} \]}
{\( a_2=33 \), \( q=3 \)}{\( a_2=11 \), \( q=3 \)}{\( a_2=22 \), \( q=3 \)}{\( a_2=33 \), \( q=2 \)}{\( a_2=11 \), \( q=2 \)}{}}
\LANGes{
\Question{
Averigua el segundo término y la razón de una progresión geométrica\( \{a_n\}_{n=1}^{\infty} \), si sabemos que: \[ \begin{aligned} a_2-a_1&=22, \\ a_3-a_2&=66. \end{aligned} \]}
{\( a_2=33 \), \( q=3 \)}{\( a_2=11 \), \( q=3 \)}{\( a_2=22 \), \( q=3 \)}{\( a_2=33 \), \( q=2 \)}{\( a_2=11 \), \( q=2 \)}{}}
\End

\Data\ID{1003134707}
\Updated{2022-04-23 05:04:20}
\Level{B}
\SubArea{Geometric Sequences}
\LANGen{
\Question{
Find the sum of the first five terms of the geometric sequence \( \{a_n\}_{n=1}^{\infty} \), if: \[ \begin{aligned} a_1+a_3&=8, \\ a_1+a_2&=0. \end{aligned} \]}
{\( 4 \)}{\( -4 \)}{\( 0 \)}{\( 8 \)}{\( -8 \)}{}}
\LANGpl{
\Question{
Wyznacz sumę pięciu pierwszych wyrazów ciągu geometrycznego \( \{a_n\}_{n=1}^{\infty} \), jeśli: \[ \begin{aligned} a_1+a_3&=8, \\ a_1+a_2&=0. \end{aligned} \]}
{\( 4 \)}{\( -4 \)}{\( 0 \)}{\( 8 \)}{\( -8 \)}{}}
\LANGcs{
\Question{
Určete součet prvních pěti členů geometrické posloupnosti \( \{a_n\}_{n=1}^{\infty} \), platí-li: \[ \begin{aligned} a_1+a_3&=8, \\ a_1+a_2&=0. \end{aligned} \]}
{\( 4 \)}{\( -4 \)}{\( 0 \)}{\( 8 \)}{\( -8 \)}{}}
\LANGsk{
\Question{
Určte súčet prvých piatich členov geometrickej postupnosti \( \{a_n\}_{n=1}^{\infty} \), ak platí: \[ \begin{aligned} a_1+a_3&=8, \\ a_1+a_2&=0. \end{aligned} \]}
{\( 4 \)}{\( -4 \)}{\( 0 \)}{\( 8 \)}{\( -8 \)}{}}
\LANGes{
\Question{
Halla la suma de primeros cinco términos de una progresión geométrica \( \{a_n\}_{n=1}^{\infty} \), si sabemos que: \[ \begin{aligned} a_1+a_3&=8, \\ a_1+a_2&=0. \end{aligned} \]}
{\( 4 \)}{\( -4 \)}{\( 0 \)}{\( 8 \)}{\( -8 \)}{}}
\End

\Data\ID{1003134708}
\Updated{2022-01-13 02:01:27}
\Level{B}
\SubArea{Geometric Sequences}
\LANGen{
\Question{
Find the sum of the second and the third term of the geometric sequence \( \{a_n\}_{n=1}^{\infty} \), if: \[ \begin{aligned} \frac{a_4}{a_1}&=-8, \\ a_4-a_2&=-9. \end{aligned} \]}
{\( 3 \)}{\( -3 \)}{\( \frac32 \)}{\( -2 \)}{\( \frac92 \)}{}}
\LANGpl{
\Question{
Wyznacz sumę drugiego i trzeciego wyrazu ciągu geometrycznego \( \{a_n\}_{n=1}^{\infty} \), jeśli: \[ \begin{aligned} \frac{a_4}{a_1}&=-8, \\ a_4-a_2&=-9. \end{aligned} \]}
{\( 3 \)}{\( -3 \)}{\( \frac32 \)}{\( -2 \)}{\( \frac92 \)}{}}
\LANGcs{
\Question{
Určete součet druhého a třetího členu geometrické posloupnosti \( \{a_n\}_{n=1}^{\infty} \), platí-li: \[ \begin{aligned} \frac{a_4}{a_1}&=-8, \\ a_4-a_2&=-9. \end{aligned} \]}
{\( 3 \)}{\( -3 \)}{\( \frac32 \)}{\( -2 \)}{\( \frac92 \)}{}}
\LANGsk{
\Question{
Určte súčet druhého a tretieho člena geometrickej postupnosti \( \{a_n\}_{n=1}^{\infty} \), ak platí: \[ \begin{aligned} \frac{a_4}{a_1}&=-8, \\ a_4-a_2&=-9. \end{aligned} \]}
{\( 3 \)}{\( -3 \)}{\( \frac32 \)}{\( -2 \)}{\( \frac92 \)}{}}
\LANGes{
\Question{
Averigua la suma del segundo y tercero término de una progresión geométrica \( \{a_n\}_{n=1}^{\infty} \), si sabemos que: \[ \begin{aligned} \frac{a_4}{a_1}&=-8, \\ a_4-a_2&=-9. \end{aligned} \]}
{\( 3 \)}{\( -3 \)}{\( \frac32 \)}{\( -2 \)}{\( \frac92 \)}{}}
\End

\Data\ID{1003134709}
\Updated{2022-01-13 02:01:56}
\Level{B}
\SubArea{Geometric Sequences}
\LANGen{
\Question{
Find the sum of the second and the third term of the geometric sequence \( \{a_n\}_{n=1}^{\infty} \), if: \[ \begin{aligned} a_1-a_2&=b,  \\ a_1+a_2&=-3b, b\in\mathbb{R}. \end{aligned}\]}
{\( -6b \)}{\( -5b \)}{\( -3b \)}{\( -2b \)}{\( 2b \)}{}}
\LANGpl{
\Question{
Wyznacz sumę drugiego i trzeciego wyrazu ciągu geometrycznego \( \{a_n\}_{n=1}^{\infty} \), jeśli: \[ \begin{aligned} a_1-a_2&=b,  \\ a_1+a_2&=-3b, b\in\mathbb{R}. \end{aligned}\]}
{\( -6b \)}{\( -5b \)}{\( -3b \)}{\( -2b \)}{\( 2b \)}{}}
\LANGcs{
\Question{
Určete součet druhého a třetího členu geometrické posloupnosti \( \{a_n\}_{n=1}^{\infty} \), platí-li: \[ \begin{aligned} a_1-a_2&=b,  \\ a_1+a_2&=-3b, b\in\mathbb{R}. \end{aligned}\]}
{\( -6b \)}{\( -5b \)}{\( -3b \)}{\( -2b \)}{\( 2b \)}{}}
\LANGsk{
\Question{
Určte súčet druhého a tretieho člena geometrickej postupnosti \( \{a_n\}_{n=1}^{\infty} \), platí-li: \[ \begin{aligned} a_1-a_2&=b,  \\ a_1+a_2&=-3b, b\in\mathbb{R}. \end{aligned}\]}
{\( -6b \)}{\( -5b \)}{\( -3b \)}{\( -2b \)}{\( 2b \)}{}}
\LANGes{
\Question{
Averigua la suma del segundo y tercer término de progresión geométrica \( \{a_n\}_{n=1}^{\infty} \), si sabemos que: \[ \begin{aligned} a_1-a_2&=b,  \\ a_1+a_2&=-3b, b\in\mathbb{R}. \end{aligned}\]}
{\( -6b \)}{\( -5b \)}{\( -3b \)}{\( -2b \)}{\( 2b \)}{}}
\End

\Data\ID{2010004901}
\Updated{2022-08-25 10:08:52}
\Level{B}
\SubArea{Geometric Sequences}
\LANGen{
\Question{
The following numbers form a geometric sequence. Find
\(x\).
\[ 
 25\, ,\ 5\, ,\ x\, ,\ a
\]}
{\(1\)}{\(\frac{1}
{5}\)}{\(-\frac{1} 
{5}\)}{\(- 1\)}{}{}}
\LANGpl{
\Question{
Poniższe liczby tworzą ciąg geometryczny. Wyznacz
\(x\).
\[ 
 25\, ,\ 5\, ,\ x\, ,\ a
\]}
{\(1\)}{\(\frac{1}
{5}\)}{\(-\frac{1} 
{5}\)}{\(- 1\)}{}{}}
\LANGcs{
\Question{
Je dáno několik po sobě jdoucích členů geometrické posloupnosti. Urči
\(x\).
\[ 
 25\, ,\ 5\, ,\ x\, ,\ a
\]}
{\(1\)}{\(\frac{1}
{5}\)}{\(-\frac{1} 
{5}\)}{\(- 1\)}{}{}}
\LANGsk{
\Question{
Je daných niekoľko po sebe idúcich členov geometrickej postupnosti. Určte \(x\).
\[ 
 25\, ,\ 5\, ,\ x\, ,\ a
\]}
{\(1\)}{\(\frac{1}
{5}\)}{\(-\frac{1} 
{5}\)}{\(- 1\)}{}{}}
\LANGes{
\Question{
Los siguientes son términos consecutivos de una progresión geométrica. Averigua
\(x\).
\[ 
 25\, ,\ 5\, ,\ x\, ,\ a
\]}
{\(1\)}{\(\frac{1}
{5}\)}{\(-\frac{1} 
{5}\)}{\(- 1\)}{}{}}
\End

\Data\ID{2010004902}
\Updated{2022-08-25 10:08:52}
\Level{B}
\SubArea{Geometric Sequences}
\LANGen{
\Question{
The following numbers form a geometric sequence. The fourth term satisfies
\(a < 0\). Find
\(x\).
\[ 
 7\, ,\ x\, ,\ \frac{1} 
{7}\, ,\ a
\]}
{\(- 1\)}{\(1\)}{\(-7\)}{\(-\frac{1} 
{7}\)}{}{}}
\LANGpl{
\Question{
Poniższe liczby tworzą ciąg geometryczny. Czwarty wyraz spełnia
\(a < 0\). Wyznacz
\(x\).
\[ 
 7\, ,\ x\, ,\ \frac{1} 
{7}\, ,\ a
\]}
{\(- 1\)}{\(1\)}{\(-7\)}{\(-\frac{1} 
{7}\)}{}{}}
\LANGcs{
\Question{
Je dáno několik po sobě jdoucích členů geometrické posloupnosti. Určete \(x\), když platí, že \(a < 0\).
\[ 
 7\, ,\ x\, ,\ \frac{1} 
{7}\, ,\ a
\]}
{\(- 1\)}{\(1\)}{\(-7\)}{\(-\frac{1} 
{7}\)}{}{}}
\LANGsk{
\Question{
Je daných niekoľko po sebe idúcich členov geometrickej postupnosti. Určte \(x\), ak platí \(a < 0\).
\[ 
 7\, ,\ x\, ,\ \frac{1} 
{7}\, ,\ a
\]}
{\(- 1\)}{\(1\)}{\(-7\)}{\(-\frac{1} 
{7}\)}{}{}}
\LANGes{
\Question{
Los siguientes son términos consecutivos de una progresión geométrica. Averigua \(x\), si vale \(a < 0\).
\[ 
 7\, ,\ x\, ,\ \frac{1} 
{7}\, ,\ a
\]}
{\(- 1\)}{\(1\)}{\(-7\)}{\(-\frac{1} 
{7}\)}{}{}}
\End

\Data\ID{2010004908}
\Updated{2022-08-25 10:08:52}
\Level{B}
\SubArea{Geometric Sequences}
\LANGen{
\Question{
At which term does the geometric sequence \( \left( \frac{1}{2187}\, ,\ \frac{1}{729}\, ,\ \frac{1}{243}\, ,\ \frac{1}{81}\, ,\ \dots \right)\) begin to have natural values?}
{\( 8\)}{\(9\)}{\(7\)}{\(10\)}{}{}}
\LANGpl{
\Question{
Na którym miejscu wyrazy ciągu geometrycznego \( \left( \frac{1}{2187}\, ,\ \frac{1}{729}\, ,\ \frac{1}{243}\, ,\ \frac {1}{81}\, ,\ \dots \right)\) zaczynają mieć naturalną wartość?}
{\( 8\)}{\(9\)}{\(7\)}{\(10\)}{}{}}
\LANGcs{
\Question{
Určete, od kolikátého členu geometrické posloupnosti  \( \left( \frac{1}{2187}\, ,\ \frac{1}{729}\, ,\ \frac{1}{243}\, ,\ \frac{1}{81}\, ,\ \dots \right)\) začnou být členy přirozená čísla.}
{\( 8\)}{\(9\)}{\(7\)}{\(10\)}{}{}}
\LANGsk{
\Question{
Určte, od ktorého člena geometrickej postupnosti \( \left( \frac{1}{2187}\, ,\ \frac{1}{729}\, ,\ \frac{1}{243}\, ,\ \frac{1}{81}\, ,\ \dots \right)\) začnú byť členmi prirodzené čísla.}
{\( 8\)}{\(9\)}{\(7\)}{\(10\)}{}{}}
\LANGes{
\Question{
Averigua desde cuál término de la progresión geométrica \( \left( \frac{1}{2187}\, ,\ \frac{1}{729}\, ,\ \frac{1}{243}\, ,\ \frac{1}{81}\, ,\ \dots \right)\) los términos empiezan a ser números naturales.}
{\( 8\)}{\(9\)}{\(7\)}{\(10\)}{}{}}
\End

\Data\ID{2010004909}
\Updated{2022-08-25 10:08:52}
\Level{B}
\SubArea{Geometric Sequences}
\LANGen{
\Question{
At which term does the geometric sequence \( \left( \frac{1}{4096}\, ,\ \frac{1}{1024}\, ,\ \frac{1}{256}\, ,\ \frac{1}{64}\, ,\ \dots \right)\) begin to have natural values?}
{\( 7\)}{\(8\)}{\(6\)}{\(9\)}{}{}}
\LANGpl{
\Question{
Na którym miejscu wyrazy ciągu geometrycznego \( \left( \frac{1}{4096}\, ,\ \frac{1}{1024}\, ,\ \frac{1}{256}\, ,\ \frac {1}{64}\, ,\ \dots \right)\) zaczynają mieć naturalną wartośći?}
{\( 7\)}{\(8\)}{\(6\)}{\(9\)}{}{}}
\LANGcs{
\Question{
Určete, od kolikátého členu geometrické posloupnosti \( \left( \frac{1}{4096}\, ,\ \frac{1}{1024}\, ,\ \frac{1}{256}\, ,\ \frac{1}{64}\, ,\ \dots \right)\) začnou být členy přirozená čísla.}
{\( 7\)}{\(8\)}{\(6\)}{\(9\)}{}{}}
\LANGsk{
\Question{
Určte, od ktorého člena geometrickej postupnosti \( \left( \frac{1}{4096}\, ,\ \frac{1}{1024}\, ,\ \frac{1}{256}\, ,\ \frac{1}{64}\, ,\ \dots \right)\) začnú byť členmi prirodzené čísla.}
{\( 7\)}{\(8\)}{\(6\)}{\(9\)}{}{}}
\LANGes{
\Question{
Averigua desde cuál término de la progresión geométrica  \( \left( \frac{1}{4096}\, ,\ \frac{1}{1024}\, ,\ \frac{1}{256}\, ,\ \frac{1}{64}\, ,\ \dots \right)\) los términos empiezan a ser números naturales.}
{\( 7\)}{\(8\)}{\(6\)}{\(9\)}{}{}}
\End

\Data\ID{2010004910}
\Updated{2022-08-25 10:08:56}
\Level{B}
\SubArea{Geometric Sequences}
\LANGen{
\Question{
At which term does the geometric sequence \( (2187\, ,\ 729 \, ,\ 243 \, , \ 81\, , \ \dots)\) begin to have values less than \(1\)?}
{\(9\)}{\(8\)}{\(7\)}{\(10\)}{}{}}
\LANGpl{
\Question{
Od którego wyrazu ciąg geometryczny \( (2187\, ,\ 729 \, ,\ 243 \, , \ 81\, , \ \dots)\) zaczyna mieć wartości mniejsze niż \(1\)?}
{\(9\)}{\(8\)}{\(7\)}{\(10\)}{}{}}
\LANGcs{
\Question{
Určete, od kolikátého členu geometrické posloupnosti \( (2187\, ,\ 729 \, ,\ 243 \, , \ 81\, , \ \dots)\) začnou být členy menší než \(1\).}
{\(9\)}{\(8\)}{\(7\)}{\(10\)}{}{}}
\LANGsk{
\Question{
Určte, od ktorého člena geometrickej postupnosti \( (2187\, ,\ 729 \, ,\ 243 \, , \ 81\, , \ \dots)\) začnú byť členy menšie ako \(1\).}
{\(9\)}{\(8\)}{\(7\)}{\(10\)}{}{}}
\LANGes{
\Question{
Averigua de cuál término de la progresión geométrica \( (2187\, ,\ 729 \, ,\ 243 \, , \ 81\, , \ \dots)\) los términos empiezan a ser menores que \(1\).}
{\(9\)}{\(8\)}{\(7\)}{\(10\)}{}{}}
\End

\Data\ID{2010004911}
\Updated{2022-08-25 10:08:56}
\Level{B}
\SubArea{Geometric Sequences}
\LANGen{
\Question{
At which term does the geometric sequence \( (4096\, ,\ 1024 \, ,\ 256 \, , \ 64\, , \ \dots)\) begin to have values less than \(1\)?}
{\(8\)}{\(7\)}{\(6\)}{\(9\)}{}{}}
\LANGpl{
\Question{
Od którego wyrazu ciąg geometryczny \( (4096\, ,\ 1024 \, ,\ 256 \, , \ 64\, , \ \dots)\) zaczyna mieć wartości mniejsze niż \(1\)?}
{\(8\)}{\(7\)}{\(6\)}{\(9\)}{}{}}
\LANGcs{
\Question{
Určete, od kolikátého členu geometrické posloupnosti \( (4096\, ,\ 1024 \, ,\ 256 \, , \ 64\, , \ \dots)\) začnou být členy menší než \(1\).}
{\(8\)}{\(7\)}{\(6\)}{\(9\)}{}{}}
\LANGsk{
\Question{
Určte, od ktorého člena geometrickej postupnosti \( (4096\, ,\ 1024 \, ,\ 256 \, , \ 64\, , \ \dots)\) začnú byť členy menšie ako \(1\).}
{\(8\)}{\(7\)}{\(6\)}{\(9\)}{}{}}
\LANGes{
\Question{
Averigua de cuál término de la progresión geométrica \( (4096\, ,\ 1024 \, ,\ 256 \, , \ 64\, , \ \dots)\) los términos empiezan a ser menores que \(1\)?}
{\(8\)}{\(7\)}{\(6\)}{\(9\)}{}{}}
\End

\Data\ID{2010004912}
\Updated{2022-08-25 10:08:56}
\Level{B}
\SubArea{Geometric Sequences}
\LANGen{
\Question{
At which term does the geometric sequence \( (10\, ;\ 12 \, ;\ 14.4 \, ; \ 17.28 \, ; \ \dots)\) exceed \(100\)? (Use a calculator.)}
{\(14\)}{\(13\)}{\(12\)}{\(15\)}{}{}}
\LANGpl{
\Question{
Od którego wyrazu ciągu geometrycznego \( (10\, ;\ 12 \, ;\ 14{,}4 \, ; \ 17{,}28 \, ; \ \dots)\) wartości wyrazów przekraczają \(100\)? (Użyj kalkulatora.)}
{\(14\)}{\(13\)}{\(12\)}{\(15\)}{}{}}
\LANGcs{
\Question{
Určete, od kolikátého členu geometrické posloupnosti \( (10\, ;\ 12 \, ;\ 14{,}4 \, ; \ 17{,}28 \, ; \ \dots)\) překročí hodnoty členů \(100\)? (Použijte kalkulačku.)}
{\(14\)}{\(13\)}{\(12\)}{\(15\)}{}{}}
\LANGsk{
\Question{
Určte, od ktorého člena geometrickej postupnosti \( (10\, ;\ 12 \, ;\ 14{,}4 \, ; \ 17{,}28 \, ; \ \dots)\) prekročia hodnoty členov \(100\). (Použite kalkulačku.)}
{\(14\)}{\(13\)}{\(12\)}{\(15\)}{}{}}
\LANGes{
\Question{
Averigua desde cuál término de la progresión geométrica \( (10\, ;\ 12 \, ;\ 14.4 \, ; \ 17.28 \, ; \ \dots)\) los términos empiezan a ser mayores que \(100\)? (Usa la calculadora.)}
{\(14\)}{\(13\)}{\(12\)}{\(15\)}{}{}}
\End

\Data\ID{2010004913}
\Updated{2022-08-25 10:08:56}
\Level{B}
\SubArea{Geometric Sequences}
\LANGen{
\Question{
At which term does the geometric sequence \( (8\, ;\ 10 \, ;\ 12.5 \, ; \ 15.625 \, ; \ \dots)\) exceed \(80\)? (Use a calculator.)}
{\(12\)}{\(11\)}{\(13\)}{\(14\)}{}{}}
\LANGpl{
\Question{
Od którego wyrazu ciągu geometrycznego \( (8\, ;\ 10 \, ;\ 12{,}5 \, ; \ 15{,}625 \, ; \ \dots)\) wartości wyrazów przekraczają \(80\)? (Użyj kalkulatora.)}
{\(12\)}{\(11\)}{\(13\)}{\(14\)}{}{}}
\LANGcs{
\Question{
Určete, od kolikátého členu geometrické posloupnosti \( (8\, ;\ 10 \, ;\ 12{,}5 \, ; \ 15{,}625 \, ; \ \dots)\) překročí hodnoty členů \(80\). (Použijte kalkulačku.)}
{\(12\)}{\(11\)}{\(13\)}{\(14\)}{}{}}
\LANGsk{
\Question{
Určte, od ktorého člena geometrickej postupnosti \( (8\, ;\ 10 \, ;\ 12{,}5 \, ; \ 15{,}625 \, ; \ \dots)\) prekročia hodnoty členov \(80\).  (Použite kalkulačku.)}
{\(12\)}{\(11\)}{\(13\)}{\(14\)}{}{}}
\LANGes{
\Question{
Averigua desde cuál término de la progresión geométrica \( (8\, ;\ 10 \, ;\ 12.5 \, ; \ 15.625 \, ; \ \dots)\) los términos empiezan a ser mayores que \(80\)? (Usa la calculadora.)}
{\(12\)}{\(11\)}{\(13\)}{\(14\)}{}{}}
\End

\Data\ID{2010004914}
\Updated{2022-08-25 10:08:56}
\Level{B}
\SubArea{Geometric Sequences}
\LANGen{
\Question{
Identify the real number \(x\) 
which converts the numbers \(a_{1} = 3^{x-6}\),
\(a_{2} = 1\) and
\(a_{3} = 3^{x}\) into
three consecutive terms of a geometric series.}
{\(x = 3\)}{\(x = 1\)}{\(x =\log 3\)}{\(x = 10\)}{\(x = 100\)}{}}
\LANGpl{
\Question{
Zidentyfikuj liczbę rzeczywistą \(x\), dla której 
liczby \(a_{1} = 3^{x-6}\),
\(a_{2} = 1\) i
\(a_{3} = 3^{x}\) są trzema kolejnymi wyrazami ciągu geometrycznego.}
{\(x = 3\)}{\(x = 1\)}{\(x =\log 3\)}{\(x = 10\)}{\(x = 100\)}{}}
\LANGcs{
\Question{
Vyberte reálné číslo \(x\) 
tak, aby čísla \(a_{1} = 3^{x-6}\),
\(a_{2} = 1\) a
\(a_{3} = 3^{x}\) 
tvořila tři po sobě jdoucí členy geometrické posloupnosti.}
{\(x = 3\)}{\(x = 1\)}{\(x =\log 3\)}{\(x = 10\)}{\(x = 100\)}{}}
\LANGsk{
\Question{
Vyberte reálne číslo \(x\) tak, aby čísla \(a_{1} = 3^{x-6}\),
\(a_{2} = 1\) a
\(a_{3} = 3^{x}\) tvorili tri po sebe idúce členy geometrickej postupnosti.}
{\(x = 3\)}{\(x = 1\)}{\(x =\log 3\)}{\(x = 10\)}{\(x = 100\)}{}}
\LANGes{
\Question{
Identifica el número real \(x\) 
que hace que los números \(a_{1} = 3^{x-6}\),
\(a_{2} = 1\) y
\(a_{3} = 3^{x}\) sean términos consecutivos de una progresión geométrica.}
{\(x = 3\)}{\(x = 1\)}{\(x =\log 3\)}{\(x = 10\)}{\(x = 100\)}{}}
\End

\Data\ID{2010004915}
\Updated{2022-08-25 10:08:56}
\Level{B}
\SubArea{Geometric Sequences}
\LANGen{
\Question{
Consider the geometric sequence with \(a_{2} = 50\) 
and \(a_{3} = 10\).
Find the sum of the first four terms.}
{\(312\)}{\(62.4\)}{\(66\)}{\(315\)}{\(312.4\)}{}}
\LANGpl{
\Question{
Rozważ ciąg geometryczny z \(a_{2} = 50\)
oraz \(a_{3} = 10\).
Znajdź sumę pierwszych czterech wyrazów.}
{\(312\)}{\(62{,}4\)}{\(66\)}{\(315\)}{\(312{,}4\)}{}}
\LANGcs{
\Question{
V geometrické posloupnosti je \(a_{2} = 50\) 
a \(a_{3} = 10\).
Určete součet prvních čtyř členů.}
{\(312\)}{\(62{,}4\)}{\(66\)}{\(315\)}{\(312{,}4\)}{}}
\LANGsk{
\Question{
V geometrickej postupnosti je  \(a_{2} = 50\) 
a \(a_{3} = 10\). Súčet prvých 4
 členov je:}
{\(312\)}{\(62{,}4\)}{\(66\)}{\(315\)}{\(312{,}4\)}{}}
\LANGes{
\Question{
Sea una progresión geométrica con \(a_{2} = 50\) 
y \(a_{3} = 10\).
Halla la suma de cuatro primeros términos.}
{\(312\)}{\(62.4\)}{\(66\)}{\(315\)}{\(312.4\)}{}}
\End

\Data\ID{2010005501}
\Updated{2022-08-25 10:08:56}
\Level{B}
\SubArea{Geometric Sequences}
\LANGen{
\Question{
The sum of the first and the second term of a geometric sequence is \( 2 \), the sum of the third and the fourth term is \( 18 \) and the common ratio is negative. Find the first term.}
{\( -1 \)}{\( 1\)}{\( 2 \)}{\( \frac12 \)}{\( -2 \)}{}}
\LANGpl{
\Question{
Suma pierwszego i drugiego wyrazu ciągu geometrycznego wynosi \( 2 \), suma trzeciego i czwartego wyrazu wynosi \( 18 \), a iloraz jest ujemny. Znajdź pierwszy wyraz ciągu.}
{\( -1 \)}{\( 1\)}{\( 2 \)}{\( \frac12 \)}{\( -2 \)}{}}
\LANGcs{
\Question{
Součet prvního a druhého členu geometrické posloupnosti je \( 2 \), součet třetího a čtvrtého členu je \( 18 \), kvocient je záporný.  Určete první člen posloupnosti.}
{\( -1 \)}{\( 1\)}{\( 2 \)}{\( \frac12 \)}{\( -2 \)}{}}
\LANGsk{
\Question{
Určte prvý člen geometrickej postupnosti, ak súčet prvého a druhého členu je \( 2 \), súčet tretieho a štvrtého je \( 18 \) a kvocient je záporný.}
{\( -1 \)}{\( 1\)}{\( 2 \)}{\( \frac12 \)}{\( -2 \)}{}}
\LANGes{
\Question{
La suma del primer y segundo término de una progresión geométrica es \( 2 \), la suma del tercer y cuarto término es \( 18 \) y la razón es negativa. Halla el primer término.}
{\( -1 \)}{\( 1\)}{\( 2 \)}{\( \frac12 \)}{\( -2 \)}{}}
\End

\Data\ID{2010005502}
\Updated{2022-08-25 10:08:56}
\Level{B}
\SubArea{Geometric Sequences}
\LANGen{
\Question{
The sum of the first three terms of a geometric sequence is \( \frac{13}9 \) and the common ratio is \( \frac13 \). Find the sum of all the terms from the \( 3 \)rd to the \( 5 \)th of the sequence.}
{\( \frac{13}{81} \)}{\( \frac{10}{81} \)}{\( \frac1{27} \)}{\( \frac{40}{81} \)}{\( \frac{121}{81} \)}{}}
\LANGpl{
\Question{
Suma pierwszych trzech wyrazów ciągu geometrycznego wynosi \( \frac{13}9 \), a iloraz to \( \frac13 \). Znajdż sumę od trzeciego do piątego wyrazu ciągu.}
{\( \frac{13}{81} \)}{\( \frac{10}{81} \)}{\( \frac1{27} \)}{\( \frac{40}{81} \)}{\( \frac{121}{81} \)}{}}
\LANGcs{
\Question{
Součet prvních tří členů geometrické posloupnosti je \( \frac{13}9 \) a kvocient je \( \frac13 \). Určete součet třetího až pátého člena posloupnosti.}
{\( \frac{13}{81} \)}{\( \frac{10}{81} \)}{\( \frac1{27} \)}{\( \frac{40}{81} \)}{\( \frac{121}{81} \)}{}}
\LANGsk{
\Question{
Určte súčet tretieho až piateho člena geometrickej postupnosti, ak súčet prvých troch členov je \( \frac{13}9 \)  a kvocient je \( \frac13 \).}
{\( \frac{13}{81} \)}{\( \frac{10}{81} \)}{\( \frac1{27} \)}{\( \frac{40}{81} \)}{\( \frac{121}{81} \)}{}}
\LANGes{
\Question{
La suma de tres primeros términos de una progresión geométrica es \( \frac{13}9 \) y la razón es \( \frac13 \). Halla la suma de todos los términos del  \( 3 \)er hasta el \( 5 \)to de la progresión.}
{\( \frac{13}{81} \)}{\( \frac{10}{81} \)}{\( \frac1{27} \)}{\( \frac{40}{81} \)}{\( \frac{121}{81} \)}{}}
\End

\Data\ID{2010005503}
\Updated{2022-08-25 10:08:56}
\Level{B}
\SubArea{Geometric Sequences}
\LANGen{
\Question{
Find the first term of the geometric sequence \( \{a_n\}_{n=1}^{\infty} \), if: \[\begin{aligned} a_5\cdot a_6&=-9, \\ a_6-a_5&=6. \end{aligned} \]}
{\( -3 \)}{\( 3 \)}{\( -1 \)}{\( 1 \)}{\( -9 \)}{}}
\LANGpl{
\Question{
Znajdź pierwszy wyraz ciągu geometrycznego \( \{a_n\}_{n=1}^{\infty} \), jeśli: \[\begin{aligned} a_5\cdot a_6&=-9, \\ a_6- a_5&=6. \end{aligned} \]}
{\( -3 \)}{\( 3 \)}{\( -1 \)}{\( 1 \)}{\( -9 \)}{}}
\LANGcs{
\Question{
Určete první člen geometrické posloupnosti \( \{a_n\}_{n=1}^{\infty} \), když: \[\begin{aligned} a_5\cdot a_6&=-9, \\ a_6-a_5&=6. \end{aligned} \]}
{\( -3 \)}{\( 3 \)}{\( -1 \)}{\( 1 \)}{\( -9 \)}{}}
\LANGsk{
\Question{
Určte prvý člen geometrickej postupnosti \( \{a_n\}_{n=1}^{\infty} \), ak: \[\begin{aligned} a_5\cdot a_6&=-9, \\ a_6-a_5&=6. \end{aligned} \]}
{\( -3 \)}{\( 3 \)}{\( -1 \)}{\( 1 \)}{\( -9 \)}{}}
\LANGes{
\Question{
Halla el primer término de la progresión geométrica \( \{a_n\}_{n=1}^{\infty} \), si: \[\begin{aligned} a_5\cdot a_6&=-9, \\ a_6-a_5&=6. \end{aligned} \]}
{\( -3 \)}{\( 3 \)}{\( -1 \)}{\( 1 \)}{\( -9 \)}{}}
\End

\Data\ID{2010005504}
\Updated{2022-08-25 10:08:56}
\Level{B}
\SubArea{Geometric Sequences}
\LANGen{
\Question{
Find the sum of the first five terms of the geometric sequence \( \{a_n\}_{n=1}^{\infty} \), if: \[ \begin{aligned} a_1+a_3&=-10, \\ a_1+a_2&=0. \end{aligned} \]}
{\( -5 \)}{\( 5 \)}{\( 0 \)}{\( -1 \)}{\( 1\)}{}}
\LANGpl{
\Question{
Znajdź sumę pierwszych pięciu wyrazów ciągu geometrycznego \( \{a_n\}_{n=1}^{\infty} \), jeśli: \[ \begin{aligned} a_1+a_3&=-10, \\ a_1+a_2&=0. \end{aligned} \]}
{\( -5 \)}{\( 5 \)}{\( 0 \)}{\( -1 \)}{\( 1\)}{}}
\LANGcs{
\Question{
Určete součet prvních pěti členů geometrické posloupnosti \( \{a_n\}_{n=1}^{\infty} \), když: \[ \begin{aligned} a_1+a_3&=-10, \\ a_1+a_2&=0. \end{aligned} \]}
{\( -5 \)}{\( 5 \)}{\( 0 \)}{\( -1 \)}{\( 1\)}{}}
\LANGsk{
\Question{
Určte súčet prvých piatich členov geometrickej postupnosti \( \{a_n\}_{n=1}^{\infty} \), ak: \[ \begin{aligned} a_1+a_3&=-10, \\ a_1+a_2&=0. \end{aligned} \]}
{\( -5 \)}{\( 5 \)}{\( 0 \)}{\( -1 \)}{\( 1\)}{}}
\LANGes{
\Question{
Halla la suma de los primeros cinco términos de la progresión geométrica \( \{a_n\}_{n=1}^{\infty} \), si: \[ \begin{aligned} a_1+a_3&=-10, \\ a_1+a_2&=0. \end{aligned} \]}
{\( -5 \)}{\( 5 \)}{\( 0 \)}{\( -1 \)}{\( 1\)}{}}
\End

\Data\ID{2010005505}
\Updated{2022-08-25 10:08:59}
\Level{B}
\SubArea{Geometric Sequences}
\LANGen{
\Question{
The sum of the first two terms of a geometric sequence is \( 54 \) and its first term equals \( 3 \). Which of the following statements about its common ratio \( q \) is not valid?}
{\( q \) is an even number.}{\( q > 10 \)}{\( q < 54 \)}{\( q \) is a prime number.}{\( q \) is a divisor of \( 51 \).}{}}
\LANGpl{
\Question{
Suma pierwszych dwóch wyrazów ciągu geometrycznego wynosi \( 54 \), a jego pierwszy wyraz równa się \( 3 \). Które z poniższych stwierdzeń dotyczących ilorazu \( q \) nie jest prawidłowe?}
{\( q \) jest liczbą parzystą.}{\( q > 10 \)}{\( q < 54 \)}{\( q \) jest liczbą pierwszą.}{\( q \) jest dzielnikem \( 51 \).}{}}
\LANGcs{
\Question{
Součet prvních dvou členů geometrické posloupnosti je \( 54 \) a první člen je \( 3 \). Které z následujících tvrzení neplatí pro kvocient \( q \)?}
{\( q \) je sudé číslo.}{\( q > 10 \)}{\( q < 54 \)}{\( q \) je prvočíslo.}{\( q \) je dělitel čísla \( 51 \).}{}}
\LANGsk{
\Question{
Súčet prvých dvoch členov geometrickej postupnosti je \( 54 \) a prvý člen je rovný \( 3 \). Ktoré z nasledujúcich tvrdení neplatí pre kvocient \( q \)?}
{\( q \) je párne číslo.}{\( q > 10 \)}{\( q < 54 \)}{\( q \) je prvočíslo.}{\( q \) je deliteľ \( 51 \).}{}}
\LANGes{
\Question{
La suma de primeros dos términos de una progresión geométrica es \( 54 \) y su primer término es igual a  \( 3 \). ¿Cuál de las declaraciones sobre la razón  \( q \) no es verdadera?}
{\( q \) es un número par.}{\( q > 10 \)}{\( q < 54 \)}{\( q \) es un número primo.}{\( q \) es divisor de \( 51 \).}{}}
\End

\Data\ID{2010005506}
\Updated{2022-08-25 10:08:59}
\Level{B}
\SubArea{Geometric Sequences}
\LANGen{
\Question{
The sum of the first \( n \) terms of a geometric sequence is \( 1 \), the common ratio is \( -3 \) and the first term equals \( \frac17 \). Find \( n \).}
{\( 3\)}{\( 4 \)}{\( 5 \)}{\( 6 \)}{\( 2 \)}{}}
\LANGpl{
\Question{
Suma pierwszych \( n \) wyrazów ciągu geometrycznego wynosi \( 1 \), iloraz wynosi \( -3 \), a pierwszy wyraz równa się \( \frac17 \). Znajdź \( n \).}
{\( 3\)}{\( 4 \)}{\( 5 \)}{\( 6 \)}{\( 2 \)}{}}
\LANGcs{
\Question{
Součet prvních \( n \) členů geometrické posloupnosti je \( 1 \), kvocient je \( -3 \) a první člen je \( \frac17 \). Určete \( n \).}
{\( 3\)}{\( 4 \)}{\( 5 \)}{\( 6 \)}{\( 2 \)}{}}
\LANGsk{
\Question{
Súčet prvých \( n \) členov geometrickej postupnosti je \( 1 \), kvocient je \( -3 \) a prvý člen je rovný \( \frac17 \). Určte \( n \).}
{\( 3\)}{\( 4 \)}{\( 5 \)}{\( 6 \)}{\( 2 \)}{}}
\LANGes{
\Question{
La suma de primeros \( n \) términos de una progresión geométrica es  \( 1 \), la razón es \( -3 \) y el primer término es igual a \( \frac17 \). Halla \( n \).}
{\( 3\)}{\( 4 \)}{\( 5 \)}{\( 6 \)}{\( 2 \)}{}}
\End

\Data\ID{9000068701}
\Updated{2022-01-13 02:01:06}
\Level{B}
\SubArea{Geometric Sequences}
\LANGen{
\Question{
Identify the real number \(x\) 
which converts the numbers \(a_{1} = 10^{2}\),
\(a_{2} = 10^{3}\) and
\(a_{3} = x\) into
three consecutive numbers of a geometric series.}
{\(x = 10\: 000\)}{\(x = 1\: 000\)}{\(x = 1\: 900\)}{\(x = 1\: 990\)}{\(x = 100\: 000\)}{}}
\LANGpl{
\Question{
Wyznacz liczbę rzeczywistą \(x\), dla której 
 \(a_{1} = 10^{2}\),
\(a_{2} = 10^{3}\) i
\(a_{3} = x\) 
są trzema kolejnymi wyrazami ciągu geometrycznego.}
{\(x = 10\: 000\)}{\(x = 1\: 000\)}{\(x = 1\: 900\)}{\(x = 1\: 990\)}{\(x = 100\: 000\)}{}}
\LANGcs{
\Question{
Vyberte reálné číslo \(x\) 
tak, aby čísla \(a_{1} = 10^{2}\),
\(a_{2} = 10^{3}\),
\(a_{3} = x\) 
tvořila tři po sobě jdoucí členy geometrické posloupnosti.}
{\(x = 10\: 000\)}{\(x = 1\: 000\)}{\(x = 1\: 900\)}{\(x = 1\: 990\)}{\(x = 100\: 000\)}{}}
\LANGsk{
\Question{
Vyberte reálne číslo \(x\) 
tak, aby čísla \(a_{1} = 10^{2}\),
\(a_{2} = 10^{3}\),
\(a_{3} = x\) 
tvorili tri po sebe idúce členy geometrickej postupnosti.}
{\(x = 10\: 000\)}{\(x = 1\: 000\)}{\(x = 1\: 900\)}{\(x = 1\: 990\)}{\(x = 100\: 000\)}{}}
\LANGes{
\Question{
Elige un número real \(x\) 
tal que los números \(a_{1} = 10^{2}\),
\(a_{2} = 10^{3}\),
\(a_{3} = x\) 
sean tres términos consecutivos de alguna progresión geométrica.}
{\(x = 10\: 000\)}{\(x = 1\: 000\)}{\(x = 1\: 900\)}{\(x = 1\: 990\)}{\(x = 100\: 000\)}{}}
\End

\Data\ID{9000068702}
\Updated{2022-01-13 02:01:46}
\Level{B}
\SubArea{Geometric Sequences}
\LANGen{
\Question{
Identify the real number \(x\) 
which converts the numbers \(a_{1} = -12\),
\(a_{2} = x\) and
\(a_{3} = -48\) into
three consecutive numbers of a geometric series.}
{\(x = -24\)}{\(x = 0\)}{\(x = 6\)}{\(x = 12\)}{\(x = -2\)}{}}
\LANGpl{
\Question{
Wyznacz liczbę rzeczywistą  \(x\), dla której \(a_{1} = -12\),
\(a_{2} = x\) i
\(a_{3} = -48\) 
są trzema kolejnymi wyrazami ciągu geometrycznego.}
{\(x = -24\)}{\(x = 0\)}{\(x = 6\)}{\(x = 12\)}{\(x = -2\)}{}}
\LANGcs{
\Question{
Vyberte reálné číslo \(x\) 
tak, aby čísla \(a_{1} = -12\),
\(a_{2} = x\),
\(a_{3} = -48\) 
tvořila tři po sobě jdoucí členy geometrické posloupnosti.}
{\(x = -24\)}{\(x = 0\)}{\(x = 6\)}{\(x = 12\)}{\(x = -2\)}{}}
\LANGsk{
\Question{
Vyberte reálne číslo \(x\) 
tak, aby čísla \(a_{1} = -12\),
\(a_{2} = x\),
\(a_{3} = -48\) 
tvorili tri po sebe idúce členy geometrickej postupnosti.}
{\(x = -24\)}{\(x = 0\)}{\(x = 6\)}{\(x = 12\)}{\(x = -2\)}{}}
\LANGes{
\Question{
Elige un número real \(x\) 
tal que los números \(a_{1} = -12\),
\(a_{2} = x\),
\(a_{3} = -48\) 
sean tres términos consecutivos de alguna progresión geométrica.}
{\(x = -24\)}{\(x = 0\)}{\(x = 6\)}{\(x = 12\)}{\(x = -2\)}{}}
\End

\Data\ID{9000068703}
\Updated{2022-01-13 02:01:22}
\Level{B}
\SubArea{Geometric Sequences}
\LANGen{
\Question{
Identify the real number \(x\) 
which converts the numbers \(a_{1} = -x\),
\(a_{2} = -5\) and
\(a_{3} = 5\) into
three consecutive numbers of a geometric series.}
{\(x = -5\)}{\(x = 0\)}{\(x = 5\)}{\(x = -10\)}{\(x = 10\)}{}}
\LANGpl{
\Question{
Wyznacz wartość liczby rzeczywistej  \(x\), dla której \(a_{1} = -x\),
\(a_{2} = -5\) i
\(a_{3} = 5\) 
są trzema kolejnymi wyrazami ciągu geometrycznego.}
{\(x = -5\)}{\(x = 0\)}{\(x = 5\)}{\(x = -10\)}{\(x = 10\)}{}}
\LANGcs{
\Question{
Vyberte reálné číslo \(x\) 
tak, aby čísla \(a_{1} = -x\),
\(a_{2} = -5\),
\(a_{3} = 5\) 
tvořila tři po sobě jdoucí členy geometrické posloupnosti.}
{\(x = -5\)}{\(x = 0\)}{\(x = 5\)}{\(x = -10\)}{\(x = 10\)}{}}
\LANGsk{
\Question{
Vyberte reálne číslo \(x\) 
tak, aby čísla \(a_{1} = -x\),
\(a_{2} = -5\),
\(a_{3} = 5\) 
tvorili tri po sebe idúce členy geometrickej postupnosti.}
{\(x = -5\)}{\(x = 0\)}{\(x = 5\)}{\(x = -10\)}{\(x = 10\)}{}}
\LANGes{
\Question{
Elige un número real \(x\) 
tal que los números \(a_{1} = -x\),
\(a_{2} = -5\),
\(a_{3} = 5\) 
sean tres términos consecutivos de alguna progresión geométrica.}
{\(x = -5\)}{\(x = 0\)}{\(x = 5\)}{\(x = -10\)}{\(x = 10\)}{}}
\End

\Data\ID{9000068704}
\Updated{2021-09-23 03:09:23}
\Level{B}
\SubArea{Geometric Sequences}
\LANGen{
\Question{
Identify the real number \(x\) 
which converts the numbers \(a_{1} = x\),
\(a_{2} = x + 5\) and
\(a_{3} = 4x\) into
three consecutive numbers of a geometric series.}
{\(x = 5\)}{\(x = 1\)}{\(x = 2\)}{\(x = 3\)}{\(x = 4\)}{}}
\LANGpl{
\Question{
Wyznacz wartość liczby rzeczywistej  \(x\), dla której \(a_{1} = x\),
\(a_{2} = x + 5\) i
\(a_{3} = 4x\) 
są trzema kolejnymi wyrazami ciągu geometrycznego.}
{\(x = 5\)}{\(x = 1\)}{\(x = 2\)}{\(x = 3\)}{\(x = 4\)}{}}
\LANGcs{
\Question{
Vyberte reálné číslo \(x\) 
tak, aby čísla \(a_{1} = x\),
\(a_{2} = x + 5\),
\(a_{3} = 4x\) 
tvořila tři po sobě jdoucí členy geometrické posloupnosti.}
{\(x = 5\)}{\(x = 1\)}{\(x = 2\)}{\(x = 3\)}{\(x = 4\)}{}}
\LANGsk{
\Question{
Vyberte reálne číslo \(x\) 
tak, aby čísla \(a_{1} = x\),
\(a_{2} = x + 5\),
\(a_{3} = 4x\) 
tvorili tri po sebe idúce členy geometrickej postupnosti.}
{\(x = 5\)}{\(x = 1\)}{\(x = 2\)}{\(x = 3\)}{\(x = 4\)}{}}
\LANGes{
\Question{
Elige el número real \(x\) 
tal que los números \(a_{1} = x\),
\(a_{2} = x + 5\),
\(a_{3} = 4x\) 
sean tres términos consecutivos de alguna progresión geométrica.}
{\(x = 5\)}{\(x = 1\)}{\(x = 2\)}{\(x = 3\)}{\(x = 4\)}{}}
\End

\Data\ID{9000068705}
\Updated{2025-01-26 12:01:01}
\Level{B}
\SubArea{Geometric Sequences}
\LANGen{
\Question{
Identify the real number \(x\) 
which converts the numbers \(a_{1} = x - 6\),
\(a_{2} = x\) and
\(a_{3} = -x\) into
three consecutive numbers of a geometric series.}
{\(x = 3\)}{\(x = 0\)}{\(x = 2\)}{\(x = 2.5\)}{\(x = -12\)}{}}
\LANGpl{
\Question{
Wyznacz wartość liczby rzeczywistej \(x\), dla której  \(a_{1} = x - 6\),
\(a_{2} = x\) i
\(a_{3} = -x\) 
są trzema kolejnymi wyrazami ciągu geometrycznego.}
{\(x = 3\)}{\(x = 0\)}{\(x = 2\)}{\(x = 2.5\)}{\(x = -12\)}{}}
\LANGcs{
\Question{
Vyberte reálné číslo \(x\) 
tak, aby čísla \(a_{1} = x - 6\),
\(a_{2} = x\),
\(a_{3} = -x\) 
tvořila tři po sobě jdoucí členy geometrické posloupnosti.}
{\(x = 3\)}{\(x = 0\)}{\(x = 2\)}{\(x = 2{,}5\)}{\(x = -12\)}{}}
\LANGsk{
\Question{
Vyberte reálne číslo \(x\) 
tak, aby čísla \(a_{1} = x - 6\),
\(a_{2} = x\),
\(a_{3} = -x\) 
tvorili tri po sebe idúce členy geometrickej postupnosti.}
{\(x = 3\)}{\(x = 0\)}{\(x = 2\)}{\(x = 2{,}5\)}{\(x = -12\)}{}}
\LANGes{
\Question{
Elige un número real \(x\) 
tal que los números \(a_{1} = x - 6\),
\(a_{2} = x\),
\(a_{3} = -x\) 
sean tres términos consecutivos de alguna progresión geométrica.}
{\(x = 3\)}{\(x = 0\)}{\(x = 2\)}{\(x = 2.5\)}{\(x = -12\)}{}}
\End

\Data\ID{9000068706}
\Updated{2025-01-26 12:01:01}
\Level{B}
\SubArea{Geometric Sequences}
\LANGen{
\Question{
Identify the real number \(x\) 
which converts the numbers \(a_{1} = x + 14\),
\(a_{2} = x + 2\) and
\(a_{3} = x - 4\) into
three consecutive numbers of a geometric series.}
{\(x = 10\)}{\(x = 5\)}{\(x = 2.5\)}{\(x = 2\)}{\(x = -5\)}{}}
\LANGpl{
\Question{
Wyznacz wartość liczby rzeczywistej  \(x\), dla której  \(a_{1} = x + 14\),
\(a_{2} = x + 2\) i
\(a_{3} = x - 4\) 
są trzema kolejnymi wyrazami ciągu geometrycznego.}
{\(x = 10\)}{\(x = 5\)}{\(x = 2.5\)}{\(x = 2\)}{\(x = -5\)}{}}
\LANGcs{
\Question{
Vyberte reálné číslo \(x\) 
tak, aby čísla \(a_{1} = x + 14\),
\(a_{2} = x + 2\),
\(a_{3} = x - 4\) 
tvořila tři po sobě jdoucí členy geometrické posloupnosti.}
{\(x = 10\)}{\(x = 5\)}{\(x = 2{,}5\)}{\(x = 2\)}{\(x = -5\)}{}}
\LANGsk{
\Question{
Vyberte reálne číslo \(x\) 
tak, aby čísla \(a_{1} = x + 14\),
\(a_{2} = x + 2\),
\(a_{3} = x - 4\) 
tvorili tri po sebe idúce členy geometrickej postupnosti.}
{\(x = 10\)}{\(x = 5\)}{\(x = 2{,}5\)}{\(x = 2\)}{\(x = -5\)}{}}
\LANGes{
\Question{
Elige el número real \(x\) 
tal que los números \(a_{1} = x + 14\),
\(a_{2} = x + 2\),
\(a_{3} = x - 4\) 
sean tres términos consecutivos de alguna progresión geométrica.}
{\(x = 10\)}{\(x = 5\)}{\(x = 2.5\)}{\(x = 2\)}{\(x = -5\)}{}}
\End

\Data\ID{9000068707}
\Updated{2021-09-23 03:09:38}
\Level{B}
\SubArea{Geometric Sequences}
\LANGen{
\Question{
Identify the real number \(x\) 
which converts the numbers \(a_{1} = x^{2} - 110\),
\(a_{2} = x^{2}\) and
\(a_{3} = x^{2} - 1\: 100\) into
three consecutive numbers of a geometric series.}
{\(x = -10\)}{\(x = -1\)}{\(x = -2\)}{\(x = -4\)}{\(x = -110\)}{}}
\LANGpl{
\Question{
Wyznacz wartość liczby rzeczywistej \(x\), dla której  \(a_{1} = x^{2} - 110\),
\(a_{2} = x^{2}\) i
\(a_{3} = x^{2} - 1\: 100\) 
są trzema kolejnymi wyrazami ciągu geometrycznego.}
{\(x = -10\)}{\(x = -1\)}{\(x = -2\)}{\(x = -4\)}{\(x = -110\)}{}}
\LANGcs{
\Question{
Vyberte reálné číslo \(x\) 
tak, aby čísla \(a_{1} = x^{2} - 110\),
\(a_{2} = x^{2}\),
\(a_{3} = x^{2} - 1\: 100\) 
tvořila tři po sobě jdoucí členy geometrické posloupnosti.}
{\(x = -10\)}{\(x = -1\)}{\(x = -2\)}{\(x = -4\)}{\(x = -110\)}{}}
\LANGsk{
\Question{
Vyberte reálne číslo \(x\) 
tak, aby čísla \(a_{1} = x^{2} - 110\),
\(a_{2} = x^{2}\),
\(a_{3} = x^{2} - 1\: 100\) 
tvorili tri po sebe idúce členy geometrickej postupnosti.}
{\(x = -10\)}{\(x = -1\)}{\(x = -2\)}{\(x = -4\)}{\(x = -110\)}{}}
\LANGes{
\Question{
Elige el número real \(x\) 
tal que los números \(a_{1} = x^{2} - 110\),
\(a_{2} = x^{2}\),
\(a_{3} = x^{2} - 1\: 100\) 
sean tres términos consecutivos de alguna progresión geométrica.}
{\(x = -10\)}{\(x = -1\)}{\(x = -2\)}{\(x = -4\)}{\(x = -110\)}{}}
\End

\Data\ID{9000068708}
\Updated{2022-04-23 05:04:14}
\Level{B}
\SubArea{Geometric Sequences}
\LANGen{
\Question{
Identify the real number \(x\) 
which converts the numbers \(a_{1} = 2^{x-4}\),
\(a_{2} = 1\) and
\(a_{3} = 2^{x}\) into
three consecutive terms of a geometric series.}
{\(x = 2\)}{\(x = 1\)}{\(x =\log 2\)}{\(x = 10\)}{\(x = 100\)}{}}
\LANGpl{
\Question{
Wyznacz wartość liczby rzeczywistej \(x\), dla której  \(a_{1} = 2^{x-4}\),
\(a_{2} = 1\)  i
\(a_{3} = 2^{x}\) 
są trzema kolejnymi wyrazami ciągu geometrycznego.}
{\(x = 2\)}{\(x = 1\)}{\(x =\log 2\)}{\(x = 10\)}{\(x = 100\)}{}}
\LANGcs{
\Question{
Vyberte reálné číslo \(x\) 
tak, aby čísla \(a_{1} = 2^{x-4}\),
\(a_{2} = 1\),
\(a_{3} = 2^{x}\) 
tvořila tři po sobě jdoucí členy geometrické posloupnosti.}
{\(x = 2\)}{\(x = 1\)}{\(x =\log 2\)}{\(x = 10\)}{\(x = 100\)}{}}
\LANGsk{
\Question{
Vyberte reálne číslo \(x\) 
tak, aby čísla \(a_{1} = 2^{x-4}\),
\(a_{2} = 1\),
\(a_{3} = 2^{x}\) 
tvorili tri po sebe idúce členy geometrickej postupnosti.}
{\(x = 2\)}{\(x = 1\)}{\(x =\log 2\)}{\(x = 10\)}{\(x = 100\)}{}}
\LANGes{
\Question{
Elige el número real \(x\) 
tal que los números \(a_{1} = 2^{x-4}\),
\(a_{2} = 1\),
\(a_{3} = 2^{x}\) 
sean tres términos consecutivos de alguna progresión geométrica.}
{\(x = 2\)}{\(x = 1\)}{\(x =\log 2\)}{\(x = 10\)}{\(x = 100\)}{}}
\End

\Data\ID{9000068709}
\Updated{2021-09-23 03:09:51}
\Level{B}
\SubArea{Geometric Sequences}
\LANGen{
\Question{
Identify the real number \(x\) 
which converts the numbers \(a_{1} =\log x\),
\(a_{2} = 2 +\log x\) and
\(a_{3} = 4\log x\) into
three consecutive numbers of a geometric series.}
{\(x = 100\)}{\(x = 1\)}{\(x =\log 2\)}{\(x = 2\)}{\(x = 10\)}{}}
\LANGpl{
\Question{
Wyznacz wartość liczby rzeczywistej \(x\), dla której  \(a_{1} =\log x\),
\(a_{2} = 2 +\log x\)  i 
\(a_{3} = 4\log x\) 
są trzema kolejnymi wyrazami ciągu geometrycznego.}
{\(x = 100\)}{\(x = 1\)}{\(x =\log 2\)}{\(x = 2\)}{\(x = 10\)}{}}
\LANGcs{
\Question{
Vyberte reálné číslo \(x\) 
tak, aby čísla \(a_{1} =\log x\),
\(a_{2} = 2 +\log x\),
\(a_{3} = 4\log x\) 
tvořila tři po sobě jdoucí členy geometrické posloupnosti.}
{\(x = 100\)}{\(x = 1\)}{\(x =\log 2\)}{\(x = 2\)}{\(x = 10\)}{}}
\LANGsk{
\Question{
Vyberte reálne číslo \(x\) 
tak, aby čísla \(a_{1} =\log x\),
\(a_{2} = 2 +\log x\),
\(a_{3} = 4\log x\) 
tvorili tri po sebe idúce členy geometrickej postupnosti.}
{\(x = 100\)}{\(x = 1\)}{\(x =\log 2\)}{\(x = 2\)}{\(x = 10\)}{}}
\LANGes{
\Question{
Elige el número real \(x\) 
tal que los números \(a_{1} =\log x\),
\(a_{2} = 2 +\log x\),
\(a_{3} = 4\log x\) 
sean tres términos consecutivos de alguna progresión geométrica.}
{\(x = 100\)}{\(x = 1\)}{\(x =\log 2\)}{\(x = 2\)}{\(x = 10\)}{}}
\End

\Data\ID{9000068710}
\Updated{2021-09-23 03:09:57}
\Level{B}
\SubArea{Geometric Sequences}
\LANGen{
\Question{
Identify the real number \(x\) 
which converts the numbers \(a_{1} = 10^{2x+2}\),
\(a_{2} = 10^{4x+1}\) and
\(a_{3} = 10^{12}\) into
three consecutive numbers of a geometric series.}
{\(x = 2\)}{\(x = 4\)}{\(x = 10^{2}\)}{\(x = \frac{1} 
{2}\)}{\(x = \frac{1} 
{100}\)}{}}
\LANGpl{
\Question{
Wyznacz wartość liczby rzeczywistej \(x\), dla której  \(a_{1} = 10^{2x+2}\),
\(a_{2} = 10^{4x+1}\) i
\(a_{3} = 10^{12}\) 
są trzema kolejnymi wyrazami ciągu geometrycznego.}
{\(x = 2\)}{\(x = 4\)}{\(x = 10^{2}\)}{\(x = \frac{1} 
{2}\)}{\(x = \frac{1} 
{100}\)}{}}
\LANGcs{
\Question{
Vyberte reálné číslo \(x\) 
tak, aby čísla \(a_{1} = 10^{2x+2}\),
\(a_{2} = 10^{4x+1}\),
\(a_{3} = 10^{12}\) 
tvořila tři po sobě jdoucí členy geometrické posloupnosti.}
{\(x = 2\)}{\(x = 4\)}{\(x = 10^{2}\)}{\(x = \frac{1} 
{2}\)}{\(x = \frac{1} 
{100}\)}{}}
\LANGsk{
\Question{
Vyberte reálne číslo \(x\) 
tak, aby čísla \(a_{1} = 10^{2x+2}\),
\(a_{2} = 10^{4x+1}\),
\(a_{3} = 10^{12}\) 
tvorili tri po sebe idúce členy geometrickej postupnosti.}
{\(x = 2\)}{\(x = 4\)}{\(x = 10^{2}\)}{\(x = \frac{1} 
{2}\)}{\(x = \frac{1} 
{100}\)}{}}
\LANGes{
\Question{
Elige el número real \(x\) 
tal que los números \(a_{1} = 10^{2x+2}\),
\(a_{2} = 10^{4x+1}\),
\(a_{3} = 10^{12}\) 
sean tres términos consecutivos de alguna progresión geométrica.}
{\(x = 2\)}{\(x = 4\)}{\(x = 10^{2}\)}{\(x = \frac{1} 
{2}\)}{\(x = \frac{1} 
{100}\)}{}}
\End

\Data\ID{9000070501}
\Updated{2025-01-26 12:01:01}
\Level{B}
\SubArea{Geometric Sequences}
\LANGen{
\Question{
Consider the geometric sequence with \(a_{2} = 50\) 
and \(a_{3} = 25\).
Find the sum of the first four terms.}
{\(187.5\)}{\(93.75\)}{\(250\)}{\(375\)}{\(500\)}{}}
\LANGpl{
\Question{
Rozważ ciąg geometryczny, w którym  \(a_{2} = 50\) 
i \(a_{3} = 25\).
Oblicz sumę czterech początkowych wyrazów tego ciągu.}
{\(187{,}5\)}{\(93{,}75\)}{\(250\)}{\(375\)}{\(500\)}{}}
\LANGcs{
\Question{
V geometrické posloupnosti je \(a_{2} = 50\),
\(a_{3} = 25\). Součet
prvních \(4\) 
členů je:}
{\(187{,}5\)}{\(93{,}75\)}{\(250\)}{\(375\)}{\(500\)}{}}
\LANGsk{
\Question{
V geometrickej postupnosti je \(a_{2} = 50\),
\(a_{3} = 25\). Súčet
prvých \(4\) 
členov je:}
{\(187{,}5\)}{\(93{,}75\)}{\(250\)}{\(375\)}{\(500\)}{}}
\LANGes{
\Question{
En una progresión geométrica se cumple que \(a_{2} = 50\),
\(a_{3} = 25\). La suma de los primeros
 \(4\) términos es:}
{\(187.5\)}{\(93.75\)}{\(250\)}{\(375\)}{\(500\)}{}}
\End

\Data\ID{9000070502}
\Updated{2022-01-16 08:01:43}
\Level{B}
\SubArea{Geometric Sequences}
\LANGen{
\Question{
Consider the geometric sequence with the first term
\(a_{1} = 243\) and quotient
\(q = \frac{1} 
{3}\). Find
\(n\) such that the
sum of the first \(n\) 
terms equals \(363\).}
{\(n = 5\)}{\(n = 2\)}{\(n = 3\)}{\(n = 4\)}{\(n = 6\)}{}}
\LANGpl{
\Question{
Rozważ ciąg geometryczny, którego pierwszym wyrazem  jest 
\(a_{1} = 243\), a iloraz
\(q = \frac{1} 
{3}\). Oblicz wartość
\(n\), taką, by suma 
początkowych  \(n\) 
wyrazów była równa \(363\).}
{\(n = 5\)}{\(n = 2\)}{\(n = 3\)}{\(n = 4\)}{\(n = 6\)}{}}
\LANGcs{
\Question{
V geometrické posloupnosti je \(q = \frac{1} 
{3}\),
\(a_{1} = 243\).
Vypočtěte, kolik členů je třeba sečíst, aby jejich součet byl roven
\(363\):}
{\(5\)}{\(2\)}{\(3\)}{\(4\)}{\(6\)}{}}
\LANGsk{
\Question{
V geometrickej postupnosti je \(q = \frac{1} 
{3}\),
\(a_{1} = 243\).
Vypočítajte, koľko členov je potrebné sčítať, aby ich súčet bol rovný
\(363\):}
{\(5\)}{\(2\)}{\(3\)}{\(4\)}{\(6\)}{}}
\LANGes{
\Question{
En una progresión geométrica sabemos que la razón es \(q = \frac{1} 
{3}\), y que 
\(a_{1} = 243\).
Calcula cuántos términos tenemos que sumar para que su suma sea igual
\(363\):}
{\(5\)}{\(2\)}{\(3\)}{\(4\)}{\(6\)}{}}
\End

\Data\ID{9000072801}
\Updated{2022-01-13 08:01:28}
\Level{B}
\SubArea{Geometric Sequences}
\LANGen{
\Question{
The following numbers form a geometric sequence. Find
\(x\).
\[ 
 2\, ,\ 4\, ,\ x
\]}
{\(8\)}{\(5\)}{\(6\)}{\(16\)}{}{}}
\LANGpl{
\Question{
Podane liczby tworzą ciąg geometryczny. Oblicz wartość 
\(x\).
\[ 
 2\, ,\ 4\, ,\ x
\]}
{\(8\)}{\(5\)}{\(6\)}{\(16\)}{}{}}
\LANGcs{
\Question{
Níže je uvedeno několik po sobě jdoucích členů geometrické posloupnosti. Určete \(x\).
\[ 
 2\, ,\ 4\, ,\ x
\]}
{\(8\)}{\(5\)}{\(6\)}{\(16\)}{}{}}
\LANGsk{
\Question{
Je daných niekoľko po sebe idúcich členov geometrickej postupnosti. Určte \(x\).
\[ 
 2\, ,\ 4\, ,\ x
\]}
{\(8\)}{\(5\)}{\(6\)}{\(16\)}{}{}}
\LANGes{
\Question{
Abajo tenemos varios términos consecutivos de una progresión geométrica. Averigua \(x\).
\[ 
 2\, ,\ 4\, ,\ x
\]}
{\(8\)}{\(5\)}{\(6\)}{\(16\)}{}{}}
\End

\Data\ID{9000072802}
\Updated{2025-01-26 12:01:01}
\Level{B}
\SubArea{Geometric Sequences}
\LANGen{
\Question{
The following numbers form a geometric sequence. Find
\(x\).
\[ 
 100\, ,\ a\, ,\ 1\, ,\ b\, ,\ x
\]}
{\(0.01\)}{\(- 100\)}{\(0.1\)}{\(- 10\)}{}{}}
\LANGpl{
\Question{
Podane liczby tworzą ciąg geometryczny. Oblicz wartość 
\(x\).
\[ 
 100\, ,\ a\, ,\ 1\, ,\ b\, ,\ x
\]}
{\(0.01\)}{\(- 100\)}{\(0.1\)}{\(- 10\)}{}{}}
\LANGcs{
\Question{
Níže je uvedeno několik po sobě jdoucích členů geometrické posloupnosti. Určete \(x\).
\[ 
 100\, ,\ a\, ,\ 1\, ,\ b\, ,\ x
\]}
{\(0{,}01\)}{\(- 100\)}{\(0{,}1\)}{\(- 10\)}{}{}}
\LANGsk{
\Question{
Je daných niekoľko po sebe idúcich členov geometrickej postupnosti. Určte \(x\).
\[ 
 100\, ,\ a\, ,\ 1\, ,\ b\, ,\ x
\]}
{\(0{,}01\)}{\(- 100\)}{\(0{,}1\)}{\(- 10\)}{}{}}
\LANGes{
\Question{
Abajo tenemos varios términos consecutivos de una progresión geométrica. Averigua \(x\).
\[ 
 100\, ,\ a\, ,\ 1\, ,\ b\, ,\ x
\]}
{\(0.01\)}{\(- 100\)}{\(0.1\)}{\(- 10\)}{}{}}
\End

\Data\ID{9000072803}
\Updated{2025-01-26 12:01:01}
\Level{B}
\SubArea{Geometric Sequences}
\LANGen{
\Question{
The following numbers form a geometric sequence. The last term satisfies
\(a > 0\). Find
\(x\).
\[ 
 1\, ,\ x\, ,\ 2\, ,\ a
\]}
{\(\sqrt{2}\)}{\(-\sqrt{2}\)}{\(1.5\)}{\(- 1.5\)}{}{}}
\LANGpl{
\Question{
Podane liczby tworzą ciąg geometryczny. Ostatni wyraz ciągu spełnia warunek 
\(a > 0\). Oblicz wartość 
\(x\).
\[ 
 1\, ,\ x\, ,\ 2\, ,\ a
\]}
{\(\sqrt{2}\)}{\(-\sqrt{2}\)}{\(1.5\)}{\(- 1.5\)}{}{}}
\LANGcs{
\Question{
Níže je uvedeno několik po sobě jdoucích členů geometrické posloupnosti. Určete \(x\), pokud platí \(a > 0\).
\[ 
 1\, ,\ x\, ,\ 2\, ,\ a
\]}
{\(\sqrt{2}\)}{\(-\sqrt{2}\)}{\(1{,}5\)}{\(- 1{,}5\)}{}{}}
\LANGsk{
\Question{
Je daných niekoľko po sebe idúcich členov geometrickej postupnosti. Určte \(x\), ak platí \(a > 0\). 
\[ 
 1\, ,\ x\, ,\ 2\, ,\ a
\]}
{\(\sqrt{2}\)}{\(-\sqrt{2}\)}{\(1{,}5\)}{\(- 1{,}5\)}{}{}}
\LANGes{
\Question{
Abajo tenemos varios términos consecutivos de progresión geométrica. Averigua \(x\), si sabemos que  \(a > 0\).
\[ 
 1\, ,\ x\, ,\ 2\, ,\ a
\]}
{\(\sqrt{2}\)}{\(-\sqrt{2}\)}{\(1.5\)}{\(- 1.5\)}{}{}}
\End

\Data\ID{9000072804}
\Updated{2022-01-13 08:01:05}
\Level{B}
\SubArea{Geometric Sequences}
\LANGen{
\Question{
The following numbers form a geometric sequence. Find
\(x\).
\[ 
 x\, ,\ a\, ,\ 3\, ,\ b\, ,\ 9
\]}
{\(1\)}{\(- 1\)}{\(- 3\)}{\(\sqrt{3}\)}{}{}}
\LANGpl{
\Question{
Następujące liczby tworzą ciąg geometryczny. Oblicz wartość
\(x\).
\[ 
 x\, ,\ a\, ,\ 3\, ,\ b\, ,\ 9
\]}
{\(1\)}{\(- 1\)}{\(- 3\)}{\(\sqrt{3}\)}{}{}}
\LANGcs{
\Question{
Níže je uvedeno několik po sobě jdoucích členů geometrické posloupnosti. Určete \(x\).
\[ 
 x\, ,\ a\, ,\ 3\, ,\ b\, ,\ 9
\]}
{\(1\)}{\(- 1\)}{\(- 3\)}{\(\sqrt{3}\)}{}{}}
\LANGsk{
\Question{
Je daných niekoľko po sebe idúcich členov geometrickej postupnosti. Určte \(x\). 
\[ 
 x\, ,\ a\, ,\ 3\, ,\ b\, ,\ 9
\]}
{\(1\)}{\(- 1\)}{\(- 3\)}{\(\sqrt{3}\)}{}{}}
\LANGes{
\Question{
Abajo tenemos varios términos consecutivos de una progresión geométrica. Averigua \(x\).
\[ 
 x\, ,\ a\, ,\ 3\, ,\ b\, ,\ 9
\]}
{\(1\)}{\(- 1\)}{\(- 3\)}{\(\sqrt{3}\)}{}{}}
\End

\Data\ID{9000072805}
\Updated{2022-01-13 08:01:05}
\Level{B}
\SubArea{Geometric Sequences}
\LANGen{
\Question{
The following numbers form a geometric sequence. The third term satisfies
\(a < 0\). Find
\(x\).
\[ 
 x\, ,\ 5\, ,\ a\, ,\ 25
\]}
{\(-\sqrt{5}\)}{\(\sqrt{5}\)}{\(- 5\)}{\(1\)}{}{}}
\LANGpl{
\Question{
Następujące liczby tworzą ciąg geometryczny. Trzeci wyraz spełnia warunek 
\(a < 0\). Oblicz wartość 
\(x\).
\[ 
 x\, ,\ 5\, ,\ a\, ,\ 25
\]}
{\(-\sqrt{5}\)}{\(\sqrt{5}\)}{\(- 5\)}{\(1\)}{}{}}
\LANGcs{
\Question{
Níže je uvedeno několik po sobě jdoucích členů geometrické posloupnosti. Určete \(x\), pokud platí \(a < 0\). 
\[ 
 x\, ,\ 5\, ,\ a\, ,\ 25
\]}
{\(-\sqrt{5}\)}{\(\sqrt{5}\)}{\(- 5\)}{\(1\)}{}{}}
\LANGsk{
\Question{
Je daných niekoľko po sebe idúcich členov geometrickej postupnosti. Určte \(x\), ak platí \(a < 0\). 
\[ 
 x\, ,\ 5\, ,\ a\, ,\ 25
\]}
{\(-\sqrt{5}\)}{\(\sqrt{5}\)}{\(- 5\)}{\(1\)}{}{}}
\LANGes{
\Question{
Abajo tenemos varios términos consecutivos de una progresión geométrica. Averigua \(x\) si sabemos que \(a < 0\). 
\[ 
 x\, ,\ 5\, ,\ a\, ,\ 25
\]}
{\(-\sqrt{5}\)}{\(\sqrt{5}\)}{\(- 5\)}{\(1\)}{}{}}
\End

\Data\ID{9000072806}
\Updated{2022-04-23 05:04:14}
\Level{B}
\SubArea{Geometric Sequences}
\LANGen{
\Question{
The following numbers form a geometric sequence. Find
\(x\).
\[ 
 2\, ,\ 1\, ,\ a\, ,\ x
\]}
{\(\frac{1}
{4}\)}{\(\frac{1}
{2}\)}{\(-\frac{1} 
{2}\)}{\(- 1\)}{}{}}
\LANGpl{
\Question{
Następujące liczby tworzą ciąg geometryczny. Oblicz wartość
\(x\).
\[ 
 2\, ,\ 1\, ,\ a\, ,\ x
\]}
{\(\frac{1}
{4}\)}{\(\frac{1}
{2}\)}{\(-\frac{1} 
{2}\)}{\(- 1\)}{}{}}
\LANGcs{
\Question{
Níže je uvedeno několik po sobě jdoucích členů geometrické posloupnosti. Určete \(x\).
\[ 
 2\, ,\ 1\, ,\ a\, ,\ x
\]}
{\(\frac{1}
{4}\)}{\(\frac{1}
{2}\)}{\(-\frac{1} 
{2}\)}{\(- 1\)}{}{}}
\LANGsk{
\Question{
Je daných niekoľko po sebe idúcich členov geometrickej postupnosti. Určte \(x\). 
\[ 
 2\, ,\ 1\, ,\ a\, ,\ x
\]}
{\(\frac{1}
{4}\)}{\(\frac{1}
{2}\)}{\(-\frac{1} 
{2}\)}{\(- 1\)}{}{}}
\LANGes{
\Question{
Abajo tenemos varios términos consecutivos de una progresión geométrica. Averigua \(x\).
\[ 
 2\, ,\ 1\, ,\ a\, ,\ x
\]}
{\(\frac{1}
{4}\)}{\(\frac{1}
{2}\)}{\(-\frac{1} 
{2}\)}{\(- 1\)}{}{}}
\End

\Data\ID{9000072807}
\Updated{2022-05-29 12:05:06}
\Level{B}
\SubArea{Geometric Sequences}
\LANGen{
\Question{
The following numbers form a geometric sequence. The third term satisfies
\(a < 0\). Find
\(x\).
\[ 
 x\, ,\ 1\, ,\ a\, ,\ \frac{1} 
{9}
\]}
{\(- 3\)}{\(9\)}{\(3\)}{\(-\frac{1} 
{3}\)}{}{}}
\LANGpl{
\Question{
Następujące liczby tworzą ciąg geometryczny. Trzeci wyraz spełnia warunek 
\(a < 0\). Oblicz wartość 
\(x\).
\[ 
 x\, ,\ 1\, ,\ a\, ,\ \frac{1} 
{9}
\]}
{\(- 3\)}{\(9\)}{\(3\)}{\(-\frac{1} 
{3}\)}{}{}}
\LANGcs{
\Question{
Níže je uvedeno několik po sobě jdoucích členů geometrické posloupnosti. Určete \(x\), pokud platí \(a < 0\).
\[ 
 x\, ,\ 1\, ,\ a\, ,\ \frac{1} 
{9}
\]}
{\(- 3\)}{\(9\)}{\(3\)}{\(-\frac{1} 
{3}\)}{}{}}
\LANGsk{
\Question{
Je daných niekoľko po sebe idúcich členov geometrickej postupnosti. Určte \(x\), ak platí \(a < 0\).
\[ 
 x\, ,\ 1\, ,\ a\, ,\ \frac{1} 
{9}
\]}
{\(- 3\)}{\(9\)}{\(3\)}{\(-\frac{1} 
{3}\)}{}{}}
\LANGes{
\Question{
Abajo tenemos varios términos consecutivos de una progresión geométrica. Averigua \(x\) si sabemos que \(a < 0\).
\[ 
 x\, ,\ 1\, ,\ a\, ,\ \frac{1} 
{9}
\]}
{\(- 3\)}{\(9\)}{\(3\)}{\(-\frac{1} 
{3}\)}{}{}}
\End

\Data\ID{9000072808}
\Updated{2021-09-27 05:09:46}
\Level{B}
\SubArea{Geometric Sequences}
\LANGen{
\Question{
The following numbers form a geometric sequence. Find
\(x\).
\[ 
 -2\, ,\ 4\, ,\ x
\]}
{\(- 8\)}{\(8\)}{\(6\)}{\(16\)}{}{}}
\LANGpl{
\Question{
Następujące liczby tworzą ciąg geometryczny. Oblicz wartość
\(x\).
\[ 
 -2\, ,\ 4\, ,\ x
\]}
{\(- 8\)}{\(8\)}{\(6\)}{\(16\)}{}{}}
\LANGcs{
\Question{
Je dán výčet několika po sobě jdoucích členů
geometrické posloupnosti. Doplňte správnou hodnotu pro člen
\(x\).
\[ 
 -2\, ,\ 4\, ,\ x
\]}
{\(- 8\)}{\(8\)}{\(6\)}{\(16\)}{}{}}
\LANGsk{
\Question{
Je daných niekoľko po sebe idúcich členov geometrickej postupnosti. Doplňte správnu hodnotu pre člen
\(x\).
\[ 
 -2\, ,\ 4\, ,\ x
\]}
{\(- 8\)}{\(8\)}{\(6\)}{\(16\)}{}{}}
\LANGes{
\Question{
Abajo tenemos varios términos consecutivos de una progresión geométrica. Averigua \(x\).
\[ 
 -2\, ,\ 4\, ,\ x
\]}
{\(- 8\)}{\(8\)}{\(6\)}{\(16\)}{}{}}
\End

\Data\ID{9000072809}
\Updated{2022-01-13 02:01:32}
\Level{B}
\SubArea{Geometric Sequences}
\LANGen{
\Question{
The following numbers form a geometric sequence. Find
\(x\).
\[ 
 1\, ,\ a\, ,\ x\, ,\ -1
\]}
{\(1\)}{\(-\frac{1} 
{3}\)}{\(0\)}{\(-\frac{2} 
{3}\)}{}{}}
\LANGpl{
\Question{
Następujące liczby tworzą ciąg geometryczny. Oblicz wartość
\(x\).
\[ 
 1\, ,\ a\, ,\ x\, ,\ -1
\]}
{\(1\)}{\(-\frac{1} 
{3}\)}{\(0\)}{\(-\frac{2} 
{3}\)}{}{}}
\LANGcs{
\Question{
Je dán výčet několika po sobě jdoucích členů geometrické posloupnosti.
Písmena \(a\) 
a \(x\) 
označují členy geometrické posloupnosti. Doplňte správnou hodnotu pro
člen \(x\).
\[ 
 1\, ,\ a\, ,\ x\, ,\ -1
\]}
{\(1\)}{\(-\frac{1} 
{3}\)}{\(0\)}{\(-\frac{2} 
{3}\)}{}{}}
\LANGsk{
\Question{
Je daný výpočet niekoľko po sebe idúcich členov geometrickej postupnosti.
Písmena \(a\) 
a \(x\) 
označujú členy geometrickej postupnosti. Doplňte správnu hodnotu pre
člen \(x\).
\[ 
 1\, ,\ a\, ,\ x\, ,\ -1
\]}
{\(1\)}{\(-\frac{1} 
{3}\)}{\(0\)}{\(-\frac{2} 
{3}\)}{}{}}
\LANGes{
\Question{
Tenemos varios términos consecutivos de una progresión geométrica. Las letras \(a\) y \(x\) 
representan términos de dicha progresión. Averigua el valor de \(x\).
\[ 
 1\, ,\ a\, ,\ x\, ,\ -1
\]}
{\(1\)}{\(-\frac{1} 
{3}\)}{\(0\)}{\(-\frac{2} 
{3}\)}{}{}}
\End

\Data\ID{9000072810}
\Updated{2021-09-27 05:09:56}
\Level{B}
\SubArea{Geometric Sequences}
\LANGen{
\Question{
The following numbers form a geometric sequence. Find
\(x\).
\[ 
 x\, ,\ 2\cdot 3\, ,\ 3\cdot 4
\]}
{\(1\cdot 3\)}{\(1\cdot 1\)}{\(1\cdot 2\)}{\(1\cdot 4\)}{}{}}
\LANGpl{
\Question{
Następujące liczby tworzą ciąg geometryczny. Oblicz wartość
\(x\).
\[ 
 x\, ,\ 2\cdot 3\, ,\ 3\cdot 4
\]}
{\(1\cdot 3\)}{\(1\cdot 1\)}{\(1\cdot 2\)}{\(1\cdot 4\)}{}{}}
\LANGcs{
\Question{
Je dán výčet tří po sobě jdoucích členů geometrické posloupnosti. Doplňte správnou
hodnotu pro člen \(x\).
\[ 
 x\, ,\ 2\cdot 3\, ,\ 3\cdot 4
\]}
{\(1\cdot 3\)}{\(1\cdot 1\)}{\(1\cdot 2\)}{\(1\cdot 4\)}{}{}}
\LANGsk{
\Question{
Je daný výpočet troch po sebe idúcich členov geometrickej postupnosti. Doplňte správnu
hodnotu pre člen \(x\).
\[ 
 x\, ,\ 2\cdot 3\, ,\ 3\cdot 4
\]}
{\(1\cdot 3\)}{\(1\cdot 1\)}{\(1\cdot 2\)}{\(1\cdot 4\)}{}{}}
\LANGes{
\Question{
Los siguientes números son tres términos consecutivos de una progresión geométrica. Averigua el valor de \(x\).
\[ 
 x\, ,\ 2\cdot 3\, ,\ 3\cdot 4
\]}
{\(1\cdot 3\)}{\(1\cdot 1\)}{\(1\cdot 2\)}{\(1\cdot 4\)}{}{}}
\End

\Data\ID{9000073001}
\Updated{2022-01-13 08:01:41}
\Level{B}
\SubArea{Geometric Sequences}
\LANGen{
\Question{
Consider a geometric sequence \((a_{n})_{n=1}^{\infty }\).
Let \(q\) be the quotient
and \(s_{n}\) be the sum
of the first \(n\) 
terms. Given \(a_{1} = 2\) 
and \(q = 2\),
find the sum of the first five terms of the sequence.}
{\(s_{5} = 62\)}{\(s_{5} = 18\)}{\(s_{5} = 32\)}{\(s_{5} = -59\)}{}{}}
\LANGpl{
\Question{
Rozważ ciąg geometryczny  \((a_{n})_{n=1}^{\infty }\). Załóżmy, że  \(q\) jest ilorazem,  
a \(s_{n}\) sumą 
 \(n\) początkowych wyrazów. Wiedząc, że  \(a_{1} = 2\) 
i \(q = 2\),
oblicz sumę pięciu początkowych wyrazów tego ciągu.}
{\(s_{5} = 62\)}{\(s_{5} = 18\)}{\(s_{5} = 32\)}{\(s_{5} = -59\)}{}{}}
\LANGcs{
\Question{
Určete součet prvních pěti členů geometrické posloupnosti, víte-li, že \(a_{1} = 2\), \(q = 2\). Přitom \(a_{n}\) značí \(n\)-tý člen posloupnosti, \(q\) kvocient a \(s_{n}\) součet prvních \(n\)-členů této posloupnosti.}
{\(s_{5} = 62\)}{\(s_{5} = 18\)}{\(s_{5} = 32\)}{\(s_{5} = -59\)}{}{}}
\LANGsk{
\Question{
\(s_{n}\) označuje súčet prvých
\(n\)-členov geometrickej
postupnosti, \(a_{n}\) označuje
\(n\)-tý člen geometrickej
postupnosti, \(q\) 
je kvocient geometrickej postupnosti. Určte súčet prvých piatich členov geometrickej
postupnosti, ak poznáte: \(a_{1} = 2\),
\(q = 2\).}
{\(s_{5} = 62\)}{\(s_{5} = 18\)}{\(s_{5} = 32\)}{\(s_{5} = -59\)}{}{}}
\LANGes{
\Question{
Calcula la suma de los primeros cinco términos de una progresión geométrica sabiendo : \(a_{1} = 2\), \(q = 2\). Además sabemos que \(a_{n}\) representa el término \(n\)-ésimo de la progresión, \(q\) es su razón y \(s_{n}\) es la suma de primeros \(n\) términos.}
{\(s_{5} = 62\)}{\(s_{5} = 18\)}{\(s_{5} = 32\)}{\(s_{5} = -59\)}{}{}}
\End

\Data\ID{9000073002}
\Updated{2022-04-23 05:04:57}
\Level{B}
\SubArea{Geometric Sequences}
\LANGen{
\Question{
Consider a geometric sequence \((a_{n})_{n=1}^{\infty }\).
Let \(q\) be the quotient
and \(s_{n}\) be the sum
of the first \(n\) 
terms. Given \(a_{6} = 5\) 
and \(q = 1\),
find the sum of the first five terms of the sequence.}
{\(s_{5} = 25\)}{\(s_{5} = 31\)}{\(s_{5} = 6\)}{\(s_{5} = 30\)}{}{}}
\LANGpl{
\Question{
Rozważmy ciąg geometryczny  \((a_{n})_{n=1}^{\infty }\).
Załóżmy, że \(q\) jest ilorazem, a \(s_{n}\) sumą 
 \(n\)  początkowych wyrazów. Wiedząc, że  \(a_{6} = 5\) 
i \(q = 1\),
oblicz sumę pięciu początkowych wyrazów tego ciągu.}
{\(s_{5} = 25\)}{\(s_{5} = 31\)}{\(s_{5} = 6\)}{\(s_{5} = 30\)}{}{}}
\LANGcs{
\Question{
\(s_{n}\) značí součet prvních
\(n\)-členů geometrické
posloupnosti, \(a_{n}\) značí
\(n\)-tý člen geometrické
posloupnosti, \(q\) 
je kvocient geometrické posloupnosti. Určete součet prvních pěti členů geometrické
posloupnosti, znáte-li: \(a_{6} = 5\),
\(q = 1\).}
{\(s_{5} = 25\)}{\(s_{5} = 31\)}{\(s_{5} = 6\)}{\(s_{5} = 30\)}{}{}}
\LANGsk{
\Question{
\(s_{n}\) označuje súčet prvých
\(n\)-členov geometrickej
postupnosti, \(a_{n}\) označuje
\(n\)-tý člen geometrickej
postupnosti, \(q\) 
je kvocient geometrickej postupnosti. Určte súčet prvých piatich členov geometrickej
postupnosti, ak poznáte: \(a_{6} = 5\),
\(q = 1\).}
{\(s_{5} = 25\)}{\(s_{5} = 31\)}{\(s_{5} = 6\)}{\(s_{5} = 30\)}{}{}}
\LANGes{
\Question{
\(s_{n}\) es la suma de los primeros
\(n\) términos de una progresión geométrica, \(a_{n}\) es el
\(n\)-ésimo término de la progresión y \(q\) es su razón. Calcula la suma de los primeros cinco términos de la progresión geométrica si \(a_{6} = 5\),
\(q = 1\).}
{\(s_{5} = 25\)}{\(s_{5} = 31\)}{\(s_{5} = 6\)}{\(s_{5} = 30\)}{}{}}
\End

\Data\ID{9000073003}
\Updated{2022-11-01 09:11:41}
\Level{B}
\SubArea{Geometric Sequences}
\LANGen{
\Question{
Consider a geometric sequence \((a_{n})_{n=1}^{\infty }\).
Let \(q\) be the quotient
and \(s_{n}\) be the sum
of the first \(n\) 
terms. Given \(a_{1} = 1\),
\(a_{3} = 4\) and
\(a_{2} > 0\), find
the sum of the first four terms of the sequence.}
{\(s_{4} = 15\)}{\(s_{4} = -5\)}{\(s_{4} = 14\)}{\(s_{4} = 8\)}{}{}}
\LANGpl{
\Question{
Rozważmy ciąg geometryczny  \((a_{n})_{n=1}^{\infty }\).
Załóżmy, że  \(q\) jest ilorazem, a \(s_{n}\)  sumą
\(n\) 
początkowych wyrazów. Wiedząc, że \(a_{1} = 1\),
\(a_{3} = 4\) i 
\(a_{2} > 0\), 
oblicz sumę czterech początkowych wyrazów tego ciągu.}
{\(s_{4} = 15\)}{\(s_{4} = -5\)}{\(s_{4} = 14\)}{\(s_{4} = 8\)}{}{}}
\LANGcs{
\Question{
Určete součet \(s_{4}\) prvních čtyř členů geometrické posloupnosti, platí-li: \(a_{1} = 1\), \(a_{3} = 4\), \(a_{2} > 0\).}
{\(s_{4} = 15\)}{\(s_{4} = -5\)}{\(s_{4} = 14\)}{\(s_{4} = 8\)}{}{}}
\LANGsk{
\Question{
Určte súčet \(s_{4}\) prvých štyroch členov geometrickej postupnosti, ak plati: \(a_{1} = 1\), \(a_{3} = 4\), \(a_{2} > 0\).}
{\(s_{4} = 15\)}{\(s_{4} = -5\)}{\(s_{4} = 14\)}{\(s_{4} = 8\)}{}{}}
\LANGes{
\Question{
Averigua la suma \(s_{4}\) de primeros cuatro términos de una progresión geométrica si sabemos que:  \(a_{1} = 1\), \(a_{3} = 4\), \(a_{2} > 0\).}
{\(s_{4} = 15\)}{\(s_{4} = -5\)}{\(s_{4} = 14\)}{\(s_{4} = 8\)}{}{}}
\End

\Data\ID{9000073004}
\Updated{2022-01-06 01:01:03}
\Level{B}
\SubArea{Geometric Sequences}
\LANGen{
\Question{
Consider a geometric sequence \((a_{n})_{n=1}^{\infty }\).
Let \(q\) be the quotient
and \(s_{n}\) be the sum
of the first \(n\) 
terms. Given \(a_{1} = 1\),
\(a_{3} = 4\) and
\(a_{2} < 0\), find
the sum of the first four terms of the sequence.}
{\(s_{4} = -5\)}{\(s_{4} = 15\)}{\(s_{4} = 14\)}{\(s_{4} = 8\)}{}{}}
\LANGpl{
\Question{
Rozważmy ciąg geometryczny \((a_{n})_{n=1}^{\infty }\).
Załóżmy, że \(q\) jest ilorazem, a 
\(s_{n}\)  sumą
 \(n\) początkowych wyrazów. Wiedząc, że \(a_{1} = 1\),
\(a_{3} = 4\) i 
\(a_{2} < 0\), oblicz sumę czterech początkowych wyrazów tego ciągu.}
{\(s_{4} = -5\)}{\(s_{4} = 15\)}{\(s_{4} = 14\)}{\(s_{4} = 8\)}{}{}}
\LANGcs{
\Question{
Určete součet \(s_{4}\) prvních čtyř členů geometrické posloupnosti, platí-li: \(a_{1} = 1\), \(a_{3} = 4\), \(a_{2} < 0\).}
{\(s_{4} = -5\)}{\(s_{4} = 15\)}{\(s_{4} = 14\)}{\(s_{4} = 8\)}{}{}}
\LANGsk{
\Question{
Určte súčet \(s_{4}\) prvých štyroch členov geometrickej postupnosti, ak plati: \(a_{1} = 1\), \(a_{3} = 4\), \(a_{2} < 0\).}
{\(s_{4} = -5\)}{\(s_{4} = 15\)}{\(s_{4} = 14\)}{\(s_{4} = 8\)}{}{}}
\LANGes{
\Question{
Averigua la suma \(s_{4}\) de los primeros cuatro términos de una progresión geométrica si se cumple que:  \(a_{1} = 1\), \(a_{3} = 4\), \(a_{2} < 0\).}
{\(s_{4} = -5\)}{\(s_{4} = 15\)}{\(s_{4} = 14\)}{\(s_{4} = 8\)}{}{}}
\End

\Data\ID{9000073005}
\Updated{2025-01-26 12:01:01}
\Level{B}
\SubArea{Geometric Sequences}
\LANGen{
\Question{
Consider a geometric sequence \((a_{n})_{n=1}^{\infty }\).
Let \(q\) be the quotient
and \(s_{n}\) be the sum
of the first \(n\) 
terms. Given \(a_{2} = 1\) 
and \(a_{3} = 10\),
find the sum of the first four terms of the sequence.}
{\(s_{4} = 111.1\)}{\(s_{4} = 99.9\)}{\(s_{4} = 111\)}{\(s_{4} = 100\)}{}{}}
\LANGpl{
\Question{
Rozważmy ciąg geometryczny  \((a_{n})_{n=1}^{\infty }\).
Załóżmy, że  \(q\) jest ilorazem, a 
\(s_{n}\) sumą
 \(n\)  początkowych 
wyrazów. Wiedząc, że \(a_{2} = 1\) 
i \(a_{3} = 10\),
oblicz sumę czterech początkowych wyrazów tego ciągu.}
{\(s_{4} = 111.1\)}{\(s_{4} = 99.9\)}{\(s_{4} = 111\)}{\(s_{4} = 100\)}{}{}}
\LANGcs{
\Question{
Určete součet \(s_{4}\) prvních čtyř členů geometrické posloupnosti, platí-li: \(a_{2} = 1\), \(a_{3} = 10\).}
{\(s_{4} = 111{,}1\)}{\(s_{4} = 99{,}9\)}{\(s_{4} = 111\)}{\(s_{4} = 100\)}{}{}}
\LANGsk{
\Question{
Určte súčet \(s_{4}\) prvých štyroch členov geometrickej postupnosti, ak plati: \(a_{2} = 1\), \(a_{3} = 10\).}
{\(s_{4} = 111{,}1\)}{\(s_{4} = 99{,}9\)}{\(s_{4} = 111\)}{\(s_{4} = 100\)}{}{}}
\LANGes{
\Question{
Averigua la suma \(s_{4}\) de los primeros cuatro términos de una progresión geométrica si sabemos que: \(a_{2} = 1\), \(a_{3} = 10\).}
{\(s_{4} = 111.1\)}{\(s_{4} = 99.9\)}{\(s_{4} = 111\)}{\(s_{4} = 100\)}{}{}}
\End

\Data\ID{9000073006}
\Updated{2022-01-13 08:01:33}
\Level{B}
\SubArea{Geometric Sequences}
\LANGen{
\Question{
Consider a geometric sequence \((a_{n})_{n=1}^{\infty }\).
Let \(q\) be the quotient
and \(s_{n}\) be the sum
of the first \(n\) 
terms. Given \(a_{1} = 1\) 
and \(a_{4} = -8\),
find the sum of the first five terms of the sequence.}
{\(s_{5} = 11\)}{\(s_{5} = 31\)}{\(s_{5} = 16\)}{\(s_{5} = -16\)}{}{}}
\LANGpl{
\Question{
Rozważmy ciąg geometryczny \((a_{n})_{n=1}^{\infty }\).
Załóżmy, że \(q\) jest ilorazem, a 
\(s_{n}\) sumą
\(n\) początkowych wyrazów.  Wiedząc, że \(a_{1} = 1\) 
i \(a_{4} = -8\),
oblicz sumę pięciu początkowych wyrazów tego ciągu.}
{\(s_{5} = 11\)}{\(s_{5} = 31\)}{\(s_{5} = 16\)}{\(s_{5} = -16\)}{}{}}
\LANGcs{
\Question{
Určete součet \(s_{5}\) prvních pěti členů geometrické posloupnosti, platí-li: \(a_{1} = 1\), \(a_{4} = -8\).}
{\(s_{5} = 11\)}{\(s_{5} = 31\)}{\(s_{5} = 16\)}{\(s_{5} = -16\)}{}{}}
\LANGsk{
\Question{
Určte súčet \(s_{5}\) prvých piatich členov geometrickej postupnosti, ak plati:  \(a_{1} = 1\), \(a_{4} = -8\).}
{\(s_{5} = 11\)}{\(s_{5} = 31\)}{\(s_{5} = 16\)}{\(s_{5} = -16\)}{}{}}
\LANGes{
\Question{
Averigua la suma \(s_{5}\) de los primeros cinco términos de una progresión geométrica si :  \(a_{1} = 1\), \(a_{4} = -8\).}
{\(s_{5} = 11\)}{\(s_{5} = 31\)}{\(s_{5} = 16\)}{\(s_{5} = -16\)}{}{}}
\End

\Data\ID{9000073007}
\Updated{2022-01-13 08:01:04}
\Level{B}
\SubArea{Geometric Sequences}
\LANGen{
\Question{
Consider a geometric sequence \((a_{n})_{n=1}^{\infty }\).
Let \(q\) be the quotient
and \(s_{n}\) be the sum
of the first \(n\) 
terms. Given \(a_{1} = -1\: 000\) 
and \(a_{2} = 100\),
find the sum of the first four terms of the sequence.}
{\(s_{4} = -909\)}{\(s_{4} = -900\)}{\(s_{4} = 911\)}{\(s_{4} = -911\)}{}{}}
\LANGpl{
\Question{
Rozważmy ciąg geometryczny  \((a_{n})_{n=1}^{\infty }\).
Załóżmy, że \(q\) jest ilorazem, a 
\(s_{n}\) sumą
 \(n\) początkowych wyrazów. Wiedząc, że \(a_{1} = -1\: 000\) 
i \(a_{2} = 100\),
oblicz sumę czterech początkowych wyrazów tego ciągu.}
{\(s_{4} = -909\)}{\(s_{4} = -900\)}{\(s_{4} = 911\)}{\(s_{4} = -911\)}{}{}}
\LANGcs{
\Question{
Určete součet \(s_{4}\) prvních čtyř členů geometrické posloupnosti, platí-li: \(a_{1} = -1\: 000\), \(a_{2} = 100\).}
{\(s_{4} = -909\)}{\(s_{4} = -900\)}{\(s_{4} = 911\)}{\(s_{4} = -911\)}{}{}}
\LANGsk{
\Question{
Určte súčet \(s_{4}\) prvých štyroch členov geometrickej postupnosti, ak plati: \(a_{1} = -1\: 000\), \(a_{2} = 100\).}
{\(s_{4} = -909\)}{\(s_{4} = -900\)}{\(s_{4} = 911\)}{\(s_{4} = -911\)}{}{}}
\LANGes{
\Question{
Averigua la suma \(s_{4}\) de los primeros cuatro términos de una progresión geométrica si : \(a_{1} = -1\: 000\), \(a_{2} = 100\).}
{\(s_{4} = -909\)}{\(s_{4} = -900\)}{\(s_{4} = 911\)}{\(s_{4} = -911\)}{}{}}
\End

\Data\ID{1003109201}
\Updated{2021-12-26 08:12:16}
\Level{C}
\SubArea{Geometric Sequences}
\LANGen{
\Question{
Karel made a resolution to go running. On the first day he ran one kilometre. Every day he extended the distance run by a quarter of the distance run on the previous day. How many whole kilometres did he run on the tenth day?}
{\( 7 \)}{\( 9 \)}{\( 8 \)}{\( 6 \)}{\( 11 \)}{}}
\LANGpl{
\Question{
Karol postanowił zacząć biegać. Pierwszego dnia przebiegł jeden kilometr. Każdego dnia wydłużał dystans o jedną czwartą w porównaniu do dnia poprzedniego. Ile kilometrów przebiegł dziesiątego dnia?}
{\( 7 \)}{\( 9 \)}{\( 8 \)}{\( 6 \)}{\( 11 \)}{}}
\LANGcs{
\Question{
Karel se rozhodl chodit běhat. První den uběhl jeden kilometr. Pak prodlužoval každý den délku trasy o čtvrtinu předchozí délky. Kolik celých kilometrů běžel desátý den?}
{\( 7 \)}{\( 9 \)}{\( 8 \)}{\( 6 \)}{\( 11 \)}{}}
\LANGsk{
\Question{
Karol sa rozhodol chodiť behať. Prvý deň odbehol jeden kilometer. Potom predlžoval každý deň dĺžku trasy o štvrtinu predchádzajúcej dĺžky. Koľko celých kilometrov bežal desiaty deň?}
{\( 7 \)}{\( 9 \)}{\( 8 \)}{\( 6 \)}{\( 11 \)}{}}
\LANGes{
\Question{
Carlos decidió empezar a correr. El primer día corrió un kilómetro. Luego extendió la longitud de su ruta cada día un cuarto de la longitud del día anterior. ¿Cuántos kilómetros corrió el décimo día?}
{\( 7 \)}{\( 9 \)}{\( 8 \)}{\( 6 \)}{\( 11 \)}{}}
\End

\Data\ID{1003109202}
\Updated{2022-05-16 08:05:18}
\Level{C}
\SubArea{Geometric Sequences}
\LANGen{
\Question{
A blouse is on sale and its price was reduced twice, each time by ten percent. The difference between its original and its final price is \( 133\,\mathrm{CZK} \). What was the blouse original price before the discounts?}
{\( 700\,\mathrm{CZK} \)}{\( 665\,\mathrm{CZK} \)}{\( 1\,330\,\mathrm{CZK} \)}{\( 1\,400\,\mathrm{CZK} \)}{\( 750\,\mathrm{CZK} \)}{}}
\LANGpl{
\Question{
Na wyprzedaży cena bluzki uległa dwukrotnej obniżce, za każdym razem o dziesięć procent. Różnica pomiędzy początkową i końcową ceną bluzki wynosi  \( 133\,\mathrm{PLN} \). Jaka był początkowa cena bluzki przed obniżkami?}
{\( 700\,\mathrm{PLN} \)}{\( 665\,\mathrm{PLN} \)}{\( 1\,330\,\mathrm{PLN} \)}{\( 1\,400\,\mathrm{PLN} \)}{\( 750\,\mathrm{PLN} \)}{}}
\LANGcs{
\Question{
Halenku zlevnili dvakrát o deset procent. Kolik stála původně, jestliže rozdíl původní a konečné ceny je \( 133\,\mathrm{CZK} \) ?}
{\( 700\,\mathrm{CZK} \)}{\( 665\,\mathrm{CZK} \)}{\( 1\,330\,\mathrm{CZK} \)}{\( 1\,400\,\mathrm{CZK} \)}{\( 750\,\mathrm{CZK} \)}{}}
\LANGsk{
\Question{
Košieľku zlacneli dvakrát o desať percent. Koľko stála pôvodne, ak rozdiel pôvodnej a konečnej ceny je \( 133\,\mathrm{CZK} \) ?}
{\( 700\,\mathrm{CZK} \)}{\( 665\,\mathrm{CZK} \)}{\( 1\,330\,\mathrm{CZK} \)}{\( 1\,400\,\mathrm{CZK} \)}{\( 750\,\mathrm{CZK} \)}{}}
\LANGes{
\Question{
Una camisa fue rebajada dos veces un \(10\%\). ¿Cuál era su precio inicial si la diferencia entre el precio inicial y el final es \( 133\,\mathrm{CZK} \) ?}
{\( 700\,\mathrm{CZK} \)}{\( 665\,\mathrm{CZK} \)}{\( 1\,330\,\mathrm{CZK} \)}{\( 1\,400\,\mathrm{CZK} \)}{\( 750\,\mathrm{CZK} \)}{}}
\End

\Data\ID{1003109203}
\Updated{2022-04-23 05:04:20}
\Level{C}
\SubArea{Geometric Sequences}
\LANGen{
\Question{
The half-life of Francium is \( 21 \) minutes. How long will it take before its weight decreases from \( 512\,\mathrm{g} \) to \( 1\,\mathrm{g} \)?}
{\( 3 \) hours \( 9 \) minutes}{\( 2 \) hours \( 48 \) minutes}{\( 3 \) hours \( 30 \) minutes}{\( 2 \) hours \( 27 \) minutes}{\( 3 \) hours \( 51 \) minutes}{}}
\LANGpl{
\Question{
Okres połowicznego rozpadu fransu wynosi \( 21 \) minut. Ile czasu minie zanim jego waga zmniejszy się z \( 512\,\mathrm{g} \) do \( 1\,\mathrm{g} \)?}
{\( 3 \) godziny \( 9 \) minuty}{\( 2 \) godziny \( 48 \) minuty}{\( 3 \) godziny \( 30 \) minuty}{\( 2 \) godziny \( 27 \) minuty}{\( 3 \) godziny \( 51 \) minuty}{}}
\LANGcs{
\Question{
Poločas rozpadu Francia je \( 21 \) minut. Za jak dlouho se jeho hmotnost sníží z \( 512\,\mathrm{g} \) na \( 1\,\mathrm{g} \)?}
{\( 3 \) hodiny \( 9 \) minut}{\( 2 \) hodiny \( 48 \) minut}{\( 3 \) hodiny \( 30 \) minut}{\( 2 \) hodiny \( 27 \) minut}{\( 3 \) hodiny \( 51 \) minut}{}}
\LANGsk{
\Question{
Polčas rozpadu Francia je \( 21 \) minút. Za ako dlho sa jeho hmotnosť zníži z \( 512\,\mathrm{g} \) na \( 1\,\mathrm{g} \)?}
{\( 3 \) hodiny \( 9 \) minút}{\( 2 \) hodiny \( 48 \) minút}{\( 3 \) hodiny \( 30 \) minút}{\( 2 \) hodiny \( 27 \) minút}{\( 3 \) hodiny \( 51 \) minút}{}}
\LANGes{
\Question{
La vida media del Francio es de \( 21 \) minutos. ¿Cuánto tarda en disminuir su masa de  \( 512\,\mathrm{g} \) a \( 1\,\mathrm{g} \)?}
{\( 3 \) horas \( 9 \) minutos}{\( 2 \) horas\( 48 \) minutos}{\( 3 \) horas \( 30 \) minutos}{\( 2 \) horas \( 27 \) minutos}{\( 3 \) horas \( 51 \) minutos}{}}
\End

\Data\ID{1003109204}
\Updated{2021-12-26 08:12:57}
\Level{C}
\SubArea{Geometric Sequences}
\LANGen{
\Question{
Candles are being lit at a commemoration in the square. Gradually, every \( 15 \) seconds, two candles are lit from one already burning candle. It also took \( 15 \) seconds to lit the first candle. After \( 4 \) minutes, all candles were lit. How many candles were finally burning in the square?}
{\( 65\,535 \)}{\( 32\,767 \)}{\( 32\,768 \)}{\( 131\,070 \)}{\( 131\,071 \)}{}}
\LANGpl{
\Question{
Podczas obchodów rocznicy na placu zapalane są świece. Stopniowo co \( 15 \) sekund, od jednej płonącej świecy zapalane są dwie kolejne.  \( 15 \) sekund zajęło również rozpalenie pierwszej świecy. Po \( 4 \) minutach, wszystkie świece zostały zapalone.  Ile świec płonęło na placu?}
{\( 65\,535 \)}{\( 32\,767 \)}{\( 32\,768 \)}{\( 131\,070 \)}{\( 131\,071 \)}{}}
\LANGcs{
\Question{
Na vzpomínkové slavnosti na náměstí se zapalují svíčky. Postupně se od již hořící zapálí další dvě během \( 15 \) sekund. Kolik nakonec hořelo na náměstí svíček, jestliže i zapálení první svíčky trvalo čtvrt minuty a celkem trvalo zapalování čtyři minuty?}
{\( 65\,535 \)}{\( 32\,767 \)}{\( 32\,768 \)}{\( 131\,070 \)}{\( 131\,071 \)}{}}
\LANGsk{
\Question{
Na spomienkovej slávnosti na námestí sa zapaľujú sviečky. Postupne sa od už horiacej zapália ďalšie dve behom \( 15 \) sekúnd. Koľko nakoniec horelo na námestí sviečok, ak i zapálenie prvej sviečky trvalo štvrť minúty a celkom trvalo zapaľovanie štyri minúty?}
{\( 65\,535 \)}{\( 32\,767 \)}{\( 32\,768 \)}{\( 131\,070 \)}{\( 131\,071 \)}{}}
\LANGes{
\Question{
En un evento conmemorativo están encendiendo velas. Cada vela encendida enciende dos velas nuevas cada \( 15 \) segundos. ¿Cuántas velas hay encendidas al final, si han tardado 15 segundos en encender la primera vela y en total tardaron cuatro minutos en encender las velas?}
{\( 65\,535 \)}{\( 32\,767 \)}{\( 32\,768 \)}{\( 131\,070 \)}{\( 131\,071 \)}{}}
\End

\Data\ID{1003109206}
\Updated{2021-12-26 08:12:05}
\Level{C}
\SubArea{Geometric Sequences}
\LANGen{
\Question{
In \( 2016 \), the population of the Czech Republic increased by \( 25\,000 \) to \( 10\,578\,820 \) inhabitants. How many inhabitants will there live in the Czech Republic at the end of the year \( 2026 \), if the percentage increase is the same every year?}
{\( 10\,832\,100 \)}{\( 10\,828\,820 \)}{\( 10\,846\,286 \)}{\( 10\,831\,495 \)}{\( 10\,603\,879 \)}{}}
\LANGpl{
\Question{
W \( 2016 \), populacja Czech wzrosła o  \( 25\,000 \) do \( 10\,578\,820 \) mieszkańców. Ilu mieszkańców będzie w Czechach pod koniec  \( 2026 \), jeśli procentowy wzrost jest jednakowy każdego roku?}
{\( 10\,832\,100 \)}{\( 10\,828\,820 \)}{\( 10\,846\,286 \)}{\( 10\,831\,495 \)}{\( 10\,603\,879 \)}{}}
\LANGcs{
\Question{
V roce \( 2016 \) se počet obyvatel České republiky zvýšil o \( 25\,000 \). Kolik bude obyvatel na konci roku \( 2026 \), bude-li procentuální nárůst každý rok stejný a na konci roku \( 2016 \) bylo v ČR \( 10\,578\,820 \) obyvatel?}
{\( 10\,832\,100 \)}{\( 10\,828\,820 \)}{\( 10\,846\,286 \)}{\( 10\,831\,495 \)}{\( 10\,603\,879 \)}{}}
\LANGsk{
\Question{
V roku \( 2016 \) sa počet obyvateľov Českej republiky zvýšil o \( 25\,000 \). Koľko bude obyvateľov na konci roku \( 2026 \), ak bude percentuálny nárast každý rok rovnaký a na konci roku \( 2016 \) bolo v ČR \( 10\,578\,820 \) obyvateľov?}
{\( 10\,832\,100 \)}{\( 10\,828\,820 \)}{\( 10\,846\,286 \)}{\( 10\,831\,495 \)}{\( 10\,603\,879 \)}{}}
\LANGes{
\Question{
En el año \( 2016 \) aumentó la población de la República Checa en \( 25\,000 \). ¿Cuántos habitantes tendrá al final del \( 2026 \), si el porcentaje de crecimiento se mantiene y si al final del año \( 2016 \) la República Checa tenía\( 10\,578\,820 \) habitantes?}
{\( 10\,832\,100 \)}{\( 10\,828\,820 \)}{\( 10\,846\,286 \)}{\( 10\,831\,495 \)}{\( 10\,603\,879 \)}{}}
\End

\Data\ID{1003109207}
\Updated{2021-12-26 08:12:52}
\Level{C}
\SubArea{Geometric Sequences}
\LANGen{
\Question{
In \( 2015 \), \( 8\,688 \) choir members sang in the largest gospel choir in Manila. Imagine, the conductor wrote an e-mail on January \( 1 \)st to three choir members. Each of them forwarded the e-mail to other three choir members on the next day and so on... On which day would all members receive the e-mail?}
{\( 8 \) January}{\( 15 \) January}{\( 2 \) February}{\( 8 \) February}{\( 12 \) January}{}}
\LANGpl{
\Question{
W \( 2015 \), \( 8\,688 \) członków chóru śpiewało w największym chórze gospelowym w Manili. Wyobraźmy sobie, że dyrygent napisał do trzech członków chóru wiadomość email  \( 1 \) stycznia. Każdy z nich następnego dnia przekazał tę wiadomość do trzech kolejnych członków chóru itd. Którego dnia wiadomość otrzymają wszyscy członkowie chóru?}
{\( 8 \) stycznia}{\( 15 \) stycznia}{\( 2 \) lutego}{\( 8 \) lutego}{\( 12 \) stycznia}{}}
\LANGcs{
\Question{
V největším gospelovém sboru v Manile v roce \( 2015 \) zpívalo \( 8\,688 \) účastníků. Kdyby dirigent napsal \( 1 \). ledna e-mail třem členům, každý z nich by ho další den přeposlal dalším třem členům, atd., který den by věděli zprávu všichni členové sboru?}
{\( 8 \). ledna}{\( 15 \). ledna}{\( 2 \). února}{\( 8 \). února}{\( 12 \). ledna}{}}
\LANGsk{
\Question{
V najväčšom gospelovom zbore v Manile v roku \( 2015 \) spievalo \( 8\,688 \) účastníkov. Keby dirigent napísal \( 1 \). januára e-mail trom členom, každý z nich by ho ďalší deň preposlal ďalším trom členom, atď., ktorý deň by vedeli správu všetci členovia zbora?}
{\( 8 \). januára}{\( 15 \). januára}{\( 2 \). februára}{\( 8 \). februára}{\( 12 \). januára}{}}
\LANGes{
\Question{
En el coro de gospel más grande de Manila en el año \( 2015 \) cantaron \( 8\,688 \) personas. Si el director escribe un mensaje el \( 1 \) de enero a tres personas y cada una de ellas lo manda a otras tres personas el día siguiente, etc. ¿qué día va a llegar el mensaje a todos los participantes?}
{\( 8 \) de enero}{\( 15 \) de enero}{\( 2 \) de febrero}{\( 8 \) de febrero}{\( 12 \) de enero}{}}
\End

\Data\ID{1003158501}
\Updated{2022-01-13 02:01:12}
\Level{C}
\SubArea{Geometric Sequences}
\LANGen{
\Question{
Let \( 3 \) numbers be \( 3 \) consecutive terms of a geometric sequence with the common ratio \( q=4 \). If we increase the second number by \( 9 \), we get \( 3 \) consecutive arithmetic sequence terms. Find the first number.}
{\( 2 \)}{\( 4 \)}{\( 8 \)}{\( 16 \)}{\( 32 \)}{}}
\LANGpl{
\Question{
Niech \( 3 \) liczby będą \( 3 \) kolejnymi wyrazami ciągu geometrycznego  o wspólnym współczynniku \( q=4 \). Jeśli zwiększymy drugą liczbę o \( 9 \), otrzymamy  \( 3 \) kolejne wyrazy ciągu arytmetycznego. Wyznacz pierwszą liczbę.}
{\( 2 \)}{\( 4 \)}{\( 8 \)}{\( 16 \)}{\( 32 \)}{}}
\LANGcs{
\Question{
\( 3 \) čísla tvoří \( 3 \) po sobě jdoucí členy geometrické posloupnosti s kvocientem \( q=4 \). Jestliže druhé číslo zvětšíme o \( 9 \), dostaneme \( 3 \) po sobě jdoucí členy aritmetické posloupnosti. Určete první číslo.}
{\( 2 \)}{\( 4 \)}{\( 8 \)}{\( 16 \)}{\( 32 \)}{}}
\LANGsk{
\Question{
\( 3 \) čísla tvoria \( 3 \) po sebe idúce členy geometrickej postupnosti s kvocientom \( q=4 \). Ak druhé číslo zväčšíme o \( 9 \), dostaneme \( 3 \) po sebe idúce členy aritmetickej postupnosti. Určte prvé číslo.}
{\( 2 \)}{\( 4 \)}{\( 8 \)}{\( 16 \)}{\( 32 \)}{}}
\LANGes{
\Question{
\( 3 \) números son \( 3 \) términos consecutivos de una progresión geométrica cuya razón es \( q=4 \). Si le sumamos \( 9 \) al segundo número, obtenemos  \( 3 \) términos consecutivos de una sucesión aritmética. Averigua el primer número.}
{\( 2 \)}{\( 4 \)}{\( 8 \)}{\( 16 \)}{\( 32 \)}{}}
\End

\Data\ID{1003158502}
\Updated{2022-01-13 02:01:14}
\Level{C}
\SubArea{Geometric Sequences}
\LANGen{
\Question{
Let there be two unknown positive numbers between numbers \( 12 \) and \( 54 \). The first three numbers of these four form \( 3 \) consecutive arithmetic sequence terms, and the last \( 3 \) numbers form three consecutive terms of a geometric sequence. Find the smaller of the two unknown numbers.}
{\( 24 \)}{\( 36 \)}{\( 15 \)}{\( 20 \)}{\( 32 \)}{}}
\LANGpl{
\Question{
Niech pomiędzy \( 12 \) i \( 54 \) będą dwie liczby dodatnie. Pierwsze trzy liczby tych czterech stanowią \( 3 \)  kolejne wyrazy ciągu arytmetycznego, a \( 3 \) ostatnie liczby tworzą trzy kolejne wyrazy ciągu geometrycznego. Wyznacz mniejszą z dwóch nieznanych liczb.}
{\( 24 \)}{\( 36 \)}{\( 15 \)}{\( 20 \)}{\( 32 \)}{}}
\LANGcs{
\Question{
Mezi čísly \( 12 \) a \( 54 \) leží dvě neznámá kladná čísla. V této čtveřici první tři čísla tvoří \( 3 \) po sobě jdoucí členy aritmetické posloupnosti a poslední tři čísla tvoří \( 3 \) po sobě jdoucí členy geometrické posloupnosti. Určete menší ze dvou neznámých čísel.}
{\( 24 \)}{\( 36 \)}{\( 15 \)}{\( 20 \)}{\( 32 \)}{}}
\LANGsk{
\Question{
Medzi číslami \( 12 \) a \( 54 \) ležia dve neznáme kladné čísla. V tejto štvorici prvé tri čísla tvoria \( 3 \) po sebe idúce členy aritmetickej postupnosti a posledné tri čísla tvoria \( 3 \) po sebe idúce členy geometrickej postupnosti. Určte menšie z dvoch neznámych čísel.}
{\( 24 \)}{\( 36 \)}{\( 15 \)}{\( 20 \)}{\( 32 \)}{}}
\LANGes{
\Question{
Entre los números \( 12 \) y \( 54 \) hay dos números positivos desconocidos.Los primeros tres números forman  \( 3 \) términos consecutivos de una sucesión aritmética y últimos tres números forman \( 3 \) términos consecutivos de una progresión geométrica. Averigua el menor de los dos números desconocidos.}
{\( 24 \)}{\( 36 \)}{\( 15 \)}{\( 20 \)}{\( 32 \)}{}}
\End

\Data\ID{1003158503}
\Updated{2022-01-13 03:01:42}
\Level{C}
\SubArea{Geometric Sequences}
\LANGen{
\Question{
Let there be four numbers, such that the first three numbers form consecutive arithmetic sequence terms with a difference of \( d=-6 \), and the last three numbers form consecutive terms of a geometric sequence with the common ratio \( q=\frac23 \). Find the fourth number.}
{\( 8 \)}{\( 18 \)}{\( 12 \)}{\( -24 \)}{\( -4 \)}{}}
\LANGpl{
\Question{
Załóżmy, że mamy cztery liczby, takie, że trzy pierwsze liczby stanowią kolejne wyrazy ciągu arytmetycznego o różnicy \( d=-6 \), a trzy ostatnie liczby stanowią kolejne wyrazy ciągu geometrycznego o wspólnym współczynniku  \( q=\frac23 \). Wyznacz czwartą liczbę.}
{\( 8 \)}{\( 18 \)}{\( 12 \)}{\( -24 \)}{\( -4 \)}{}}
\LANGcs{
\Question{
Ve čtveřici čísel tvoří první tři čísla tři po sobě jdoucí členy aritmetické posloupnosti s diferencí \( d=-6 \) a poslední tři čísla tvoří tři po sobě jdoucí členy geometrické posloupnosti s kvocientem \( q=\frac23 \). Určete čtvrté číslo.}
{\( 8 \)}{\( 18 \)}{\( 12 \)}{\( -24 \)}{\( -4 \)}{}}
\LANGsk{
\Question{
Vo štvorici čísel tvoria prvé tri čísla tri po sebe idúce členy aritmetickej postupnosti s diferenciou \( d=-6 \) a posledne tri čísla tvoria tri po sebe idúce členy geometrickej postupnosti s kvocientom \( q=\frac23 \). Určte štvrté číslo.}
{\( 8 \)}{\( 18 \)}{\( 12 \)}{\( -24 \)}{\( -4 \)}{}}
\LANGes{
\Question{
Tenemos cuatro números. Los tres primeros forman una sucesión aritmética con diferencia \( d=-6 \) y los últimos tres números son tres términos consecutivos de una progresión geométrica con razón \( q=\frac23 \). Averigua el cuarto número.}
{\( 8 \)}{\( 18 \)}{\( 12 \)}{\( -24 \)}{\( -4 \)}{}}
\End

\Data\ID{1003158504}
\Updated{2022-01-13 03:01:56}
\Level{C}
\SubArea{Geometric Sequences}
\LANGen{
\Question{
Let there be three numbers that are three consecutive arithmetic sequence terms with a difference of \( d=3 \). If the third number is decreased by \( \frac32 \), we get \( 3 \) consecutive terms of a geometric sequence. Find the third number (of an arithmetic sequence).}
{\( 0 \)}{\( 3 \)}{\( -3 \)}{\( \frac32 \)}{\( -\frac32 \)}{}}
\LANGpl{
\Question{
Załóżmy, ze mamy trzy liczby, które są trzema kolejnymi wyrazami ciągu arytmetycznego o różnicy \( d=3 \). Jeśli trzecią liczbę zwiększymy o  \( \frac32 \), otrzymamy \( 3 \) kolejne wyrazy ciągu geometrycznego. Wyznacz trzecia liczbę  (ciągu arytmetycznego).}
{\( 0 \)}{\( 3 \)}{\( -3 \)}{\( \frac32 \)}{\( -\frac32 \)}{}}
\LANGcs{
\Question{
Tři čísla tvoří tři po sobě jdoucí členy aritmetické posloupnosti s diferencí \( d=3 \). Jestliže třetí číslo zmenšíme o \( \frac32 \), dostaneme \( 3 \) po sobě jdoucí členy geometrické posloupnosti. Určete třetí číslo (aritmetické posloupnosti).}
{\( 0 \)}{\( 3 \)}{\( -3 \)}{\( \frac32 \)}{\( -\frac32 \)}{}}
\LANGsk{
\Question{
Tri čísla tvoria tri po sebe idúce členy aritmetickej postupnosti s diferenciou \( d=3 \). Ak tretie číslo zmenšíme o \( \frac32 \), dostaneme \( 3 \) po sebe idúce členy geometrickej postupnosti. Určte tretie číslo (aritmetickej postupnosti).}
{\( 0 \)}{\( 3 \)}{\( -3 \)}{\( \frac32 \)}{\( -\frac32 \)}{}}
\LANGes{
\Question{
Tres números forman tres términos consecutivos de una sucesión aritmética con diferencia \( d=3 \). Si el tercer número lo disminuimos en \( \frac32 \), obtenemos \( 3 \) términos consecutivos de una progresión geométrica. Averigua el tercer número (de la sucesión aritmética).}
{\( 0 \)}{\( 3 \)}{\( -3 \)}{\( \frac32 \)}{\( -\frac32 \)}{}}
\End

\Data\ID{1003158505}
\Updated{2022-01-13 03:01:12}
\Level{C}
\SubArea{Geometric Sequences}
\LANGen{
\Question{
Let there be three numbers that are \( 3 \) consecutive arithmetic sequence terms. Their sum is \( 9 \). If the first number is divided by \( -3 \), we get \( 3 \) consecutive terms of a geometric sequence. Find the largest number of the given three.}
{\( 9 \)}{\( 3 \)}{\( 12 \)}{\( 6 \)}{\( 4 \)}{}}
\LANGpl{
\Question{
Załóżmy, ze mamy trzy liczby, które są \( 3 \) są trzema kolejnymi wyrazami ciągu arytmetycznego.  Ich suma wynosi \( 9 \). Jeśli pierwszą liczbę podzielimy przez \( -3 \), otrzymamy \( 3 \) kolejne wyrazy ciągu geometrycznego. Wyznacz największą liczbę z trzech podanych.}
{\( 9 \)}{\( 3 \)}{\( 12 \)}{\( 6 \)}{\( 4 \)}{}}
\LANGcs{
\Question{
Tři čísla tvoří \( 3 \) po sobě jdoucí členy aritmetické posloupnosti a jejich součet je \( 9 \). Jestliže první číslo vydělíme \( -3 \), dostaneme \( 3 \) po sobě jdoucí členy geometrické posloupnosti. Určete největší číslo z dané trojice.}
{\( 9 \)}{\( 3 \)}{\( 12 \)}{\( 6 \)}{\( 4 \)}{}}
\LANGsk{
\Question{
Tri čísla tvoria \( 3 \) po sebe idúce členy aritmetickej postupnosti a ich súčet je \( 9 \). Ak prvé číslo vydelíme \( -3 \), dostaneme \( 3 \) po sebe idúce členy geometrickej postupnosti. Určte najväčšie číslo z danej trojice.}
{\( 9 \)}{\( 3 \)}{\( 12 \)}{\( 6 \)}{\( 4 \)}{}}
\LANGes{
\Question{
Tres números son  \( 3 \) términos consecutivos de una progresión aritmética y su suma es \( 9 \). Si el primer número lo dividimos por \( -3 \), obtenemos \( 3 \) términos consecutivos de una progresión geométrica. Averigua el mayor número de los tres.}
{\( 9 \)}{\( 3 \)}{\( 12 \)}{\( 6 \)}{\( 4 \)}{}}
\End

\Data\ID{1003158506}
\Updated{2022-01-13 08:01:57}
\Level{C}
\SubArea{Geometric Sequences}
\LANGen{
\Question{
Let there be the first nine terms of an arithmetic sequence with the first term \( a_1=1 \) and the difference \( d=1 \). Suppose we are creating ordered triples of \( 3 \) different numbers out of the given \( 9 \) numbers so that they form \( 3 \) consecutive terms of a geometric sequence. How many of such triples can be created?}
{\( 8 \)}{\( 6 \)}{\( 4 \)}{\( 3 \)}{\( 9 \)}{}}
\LANGpl{
\Question{
Załóżmy, że mamy dziewięć pierwszych wyrazów ciągu arytmetycznego z pierwszym wyrazem równym \( a_1=1 \) i różnicą \( d=1 \). Załóżmy, ze tworzymy uporządkowane trójki  \( 3 \) różnych liczb z  \( 9 \) podanych tak, że tworzą  \( 3 \) kolejne wyrazy ciągu geometrycznego. Ile takich trójek możemy utworzyć?}
{\( 8 \)}{\( 6 \)}{\( 4 \)}{\( 3 \)}{\( 9 \)}{}}
\LANGcs{
\Question{
Z prvních devíti členů aritmetické posloupnosti s prvním členem \( a_1=1 \) a diferencí \( d=1 \) vybíráme uspořádanou trojici různých čísel tak, aby tvořila \( 3 \) po sobě jdoucí členy geometrické posloupnosti. Kolik takových trojic lze vytvořit?}
{\( 8 \)}{\( 6 \)}{\( 4 \)}{\( 3 \)}{\( 9 \)}{}}
\LANGsk{
\Question{
Z prvých deviatich členov aritmetickej postupnosti s prvým členom \( a_1=1 \) a diferenciou \( d=1 \) vyberáme usporiadanú trojicu rôznych čísel tak, aby tvorila \( 3 \) po sebe idúce členy geometrickej postupnosti. Koľko takých trojíc sa dá vytvoriť?}
{\( 8 \)}{\( 6 \)}{\( 4 \)}{\( 3 \)}{\( 9 \)}{}}
\LANGes{
\Question{
De los primeros nueve términos de una sucesión aritmética cuyo primer término es \( a_1=1 \) y la diferencia \( d=1 \) vamos a elegir tres números para que sean \( 3 \) términos consecutivos de una progresión geométrica. ¿Cuántos grupos de tres números podemos hacer?}
{\( 8 \)}{\( 6 \)}{\( 4 \)}{\( 3 \)}{\( 9 \)}{}}
\End

\Data\ID{1003158507}
\Updated{2025-01-26 12:01:01}
\Level{C}
\SubArea{Geometric Sequences}
\LANGen{
\Question{
Let there be a row of five yellow cubes lying side by side. The first cube has an edge length of \( 100\,\mathrm{cm} \) and each next cube has an edge length by \( 10\,\mathrm{cm} \) shorter than the previous one. Let there be a second row of five blue cubes lying side by side. The first blue cube has an edge length of \( 100\,\mathrm{cm} \) and each next cube has an edge length by \( 10\% \) smaller than the previous one. What is the difference between the length of these two rows?}
{\( 9.51\,\mathrm{cm} \)}{\( 34.51\,\mathrm{cm} \)}{\( 0\,\mathrm{cm} \)}{\( 20\,\mathrm{cm} \)}{\( 20.51\,\mathrm{cm} \)}{}}
\LANGpl{
\Question{
Załóżmy, że mamy rząd pięciu żółtych sześcianów leżących ściana przy ścianie. Pierwszy sześcian ma bok o długości  \( 100\,\mathrm{cm} \), a każdy następny sześcian ma bok o  \( 10\,\mathrm{cm} \) krótszy od poprzedniego. Załóżmy, że mamy drugi rząd pięciu niebieskich sześcianów leżących obok siebie. Długość boku pierwszego sześcianu wynosi \( 100\,\mathrm{cm} \), a każdy następny sześcian ma nok o  \( 10\% \) krótszy niż poprzedni. Jaka jest różnica pomiędzy długością tych dwóch rzędów?}
{\( 9{,}51\,\mathrm{cm} \)}{\( 34{,}51\,\mathrm{cm} \)}{\( 0\,\mathrm{cm} \)}{\( 20\,\mathrm{cm} \)}{\( 20{,}51\,\mathrm{cm} \)}{}}
\LANGcs{
\Question{
Jaký je rozdíl mezi délkou řady pěti žlutých krychlí ležících těsně vedle sebe, z nichž první má hranu délky \( 100\,\mathrm{cm} \) a každá další o \( 10\,\mathrm{cm} \) menší než předchozí, a délkou řady modrých krychlí ležících těsně vedle sebe, z nichž první má hranu délky \( 100\,\mathrm{cm} \) a každá další o \( 10\% \) menší než předchozí?}
{\( 9{,}51\,\mathrm{cm} \)}{\( 34{,}51\,\mathrm{cm} \)}{\( 0\,\mathrm{cm} \)}{\( 20\,\mathrm{cm} \)}{\( 20{,}51\,\mathrm{cm} \)}{}}
\LANGsk{
\Question{
Aký je rozdiel medzi dĺžkou radu piatich žltých kociek ležiacich tesne vedľa seba, z nich prvá má hranu dĺžky \( 100\,\mathrm{cm} \) a každá ďalšia o \( 10\,\mathrm{cm} \) menšia než predchádzajúca, a dĺžkou radu modrých kociek ležiacich tesne vedľa seba, z nich prvá má hranu dĺžky \( 100\,\mathrm{cm} \) a každá ďalšia o \( 10\% \) menšia než predchádzajúca?}
{\( 9{,}51\,\mathrm{cm} \)}{\( 34{,}51\,\mathrm{cm} \)}{\( 0\,\mathrm{cm} \)}{\( 20\,\mathrm{cm} \)}{\( 20{,}51\,\mathrm{cm} \)}{}}
\LANGes{
\Question{
Tenemos dos filas de cubos. La primera fila está formada por cubos amarillos. El primer cubo tiene un lado de \( 100\,\mathrm{cm} \) y cada cubo siguiente tiene un lado \( 10\,\mathrm{cm} \) menor que el anterior. La otra fila es de cubos azules, el primer cubo tiene un lado de \( 100\,\mathrm{cm} \) y cada cubo siguiente tiene un lado \( 10\% \) menor que el anterior. ¿Qué diferencia hay entre las filas?}
{\( 9.51\,\mathrm{cm} \)}{\( 34.51\,\mathrm{cm} \)}{\( 0\,\mathrm{cm} \)}{\( 20\,\mathrm{cm} \)}{\( 20.51\,\mathrm{cm} \)}{}}
\End

\Data\ID{1003170601}
\Updated{2022-04-23 05:04:20}
\Level{C}
\SubArea{Geometric Sequences}
\LANGen{
\Question{
How many numbers do we need to insert between the numbers \( 6 \) and \( 1\,458 \) so that the inserted numbers with the given two numbers are consecutive terms of a geometric sequence? The sum of all numbers inserted must be \( 720 \).}
{\( 4 \)}{\( 6 \)}{\( 3 \)}{\( 5 \)}{\( 7 \)}{}}
\LANGpl{
\Question{
Ile liczb musimy wstawić pomiędzy liczby \( 6 \) i \( 1\,458 \), tak by wstawione liczby z dwoma podanymi liczbami tworzyły kolejne wyrazy ciągu geometrycznego? Suma wszystkich wstawionych liczb musi wynosić \( 720 \).}
{\( 4 \)}{\( 6 \)}{\( 3 \)}{\( 5 \)}{\( 7 \)}{}}
\LANGcs{
\Question{
Kolik čísel musíme vložit mezi čísla \( 6 \) a \( 1\,458 \), aby spolu s danými čísly tvořila po sobě jdoucí členy geometrické posloupnosti a součet vložených čísel byl \( 720 \)?}
{\( 4 \)}{\( 6 \)}{\( 3 \)}{\( 5 \)}{\( 7 \)}{}}
\LANGsk{
\Question{
Koľko čísel musíme vložiť medzi čísla \( 6 \) a \( 1\,458 \), aby spolu s danými číslami tvorili po sebe idúce členy geometrickej postupnosti a súčet vložených čísel bol \( 720 \)?}
{\( 4 \)}{\( 6 \)}{\( 3 \)}{\( 5 \)}{\( 7 \)}{}}
\LANGes{
\Question{
¿Cuántos números tenemos que intercalar entre \( 6 \) y \( 1\,458 \) para que junto con los números dados formen términos consecutivos de una progresión geométrica y para que la suma de los números intercalados sea \( 720 \)?}
{\( 4 \)}{\( 6 \)}{\( 3 \)}{\( 5 \)}{\( 7 \)}{}}
\End

\Data\ID{1003170602}
\Updated{2025-01-26 12:01:01}
\Level{C}
\SubArea{Geometric Sequences}
\LANGen{
\Question{
The lengths of a cuboid edges form \( 3 \) consecutive terms of a geometric sequence. The volume of the cuboid is \( 140\,608\,\mathrm{cm}^3 \), the sum of its shortest and its longest edge is \( 221\,\mathrm{cm} \). Find the length of its shortest edge.}
{\( 13\,\mathrm{cm} \)}{\( 52\,\mathrm{cm} \)}{\( 4\,\mathrm{cm} \)}{\( 208\,\mathrm{cm} \)}{\( 0.25\,\mathrm{cm} \)}{}}
\LANGpl{
\Question{
Długości boków prostopadłościanu stanowią \( 3 \) kolejne wyrazy ciągu geometrycznego. Objętość prostopadłościanu wynosi  \( 140\,608\,\mathrm{cm}^3 \), suma jego  krótszych i dłuższych boków wynosi  \( 221\,\mathrm{cm} \). Wyznacz długość jego krótszego boku.}
{\( 13\,\mathrm{cm} \)}{\( 52\,\mathrm{cm} \)}{\( 4\,\mathrm{cm} \)}{\( 208\,\mathrm{cm} \)}{\( 0{,}25\,\mathrm{cm} \)}{}}
\LANGcs{
\Question{
Velikosti hran kvádru tvoří tři po sobě jdoucí členy geometrické posloupnosti. Objem kvádru je \( 140\,608\,\mathrm{cm}^3 \), součet nejkratší a nejdelší strany je \( 221\,\mathrm{cm} \). Určete velikost nejmenší strany.}
{\( 13\,\mathrm{cm} \)}{\( 52\,\mathrm{cm} \)}{\( 4\,\mathrm{cm} \)}{\( 208\,\mathrm{cm} \)}{\( 0{,}25\,\mathrm{cm} \)}{}}
\LANGsk{
\Question{
Veľkosti hrán kvádra tvoria tri po sebe idúce členy geometrickej postupnosti. Objem kvádra je \( 140\,608\,\mathrm{cm}^3 \), súčet najkratšej a najdlhšej strany je \( 221\,\mathrm{cm} \). Určte veľkosť najmenšej strany.}
{\( 13\,\mathrm{cm} \)}{\( 52\,\mathrm{cm} \)}{\( 4\,\mathrm{cm} \)}{\( 208\,\mathrm{cm} \)}{\( 0{,}25\,\mathrm{cm} \)}{}}
\LANGes{
\Question{
Las dimensiones de un prisma rectángular forman tres términos consecutivos de una progresión geométrica. El volumen del prisma es \( 140\,608\,\mathrm{cm}^3 \) y la suma del lado más corto con el lado más largo es \( 221\,\mathrm{cm} \). Averigua el tamaño del lado más corto.}
{\( 13\,\mathrm{cm} \)}{\( 52\,\mathrm{cm} \)}{\( 4\,\mathrm{cm} \)}{\( 208\,\mathrm{cm} \)}{\( 0.25\,\mathrm{cm} \)}{}}
\End

\Data\ID{1003170603}
\Updated{2022-01-16 08:01:52}
\Level{C}
\SubArea{Geometric Sequences}
\LANGen{
\Question{
Between the roots of the equation \( 9x^2+130x-75=0 \) insert two numbers so that the roots and the new numbers form \( 4 \) consecutive terms of a geometric sequence.
What is the smaller of the two inserted numbers?}
{\( -\frac53 \)}{\( \frac59 \)}{\( -\frac59 \)}{\( \frac53 \)}{\( -5 \)}{}}
\LANGpl{
\Question{
Pomiędzy pierwiastkami równania \( 9x^2+130x-75=0 \) wstaw dwie liczby takie, że pierwiastki i te dwie liczby tworzą \( 4 \) kolejne wyrazy ciągu geometrycznego.
Jaka jest wartość mniejszej z dwóch wstawionych liczb?}
{\( -\frac53 \)}{\( \frac59 \)}{\( -\frac59 \)}{\( \frac53 \)}{\( -5 \)}{}}
\LANGcs{
\Question{
Mezi kořeny rovnice \( 9x^2+130x-75=0 \) vložte dvě čísla tak, aby spolu s kořeny rovnice tvořila čtyři po sobě jdoucí členy geometrické posloupnosti. Menší z nich je rovno:}
{\( -\frac53 \)}{\( \frac59 \)}{\( -\frac59 \)}{\( \frac53 \)}{\( -5 \)}{}}
\LANGsk{
\Question{
Medzi korene rovnice \( 9x^2+130x-75=0 \) vložte dva čísla tak, aby spolu s koreňmi rovnice tvorili štyri po sebe idúce členy geometrickej postupnosti. Menší z nich je rovný:}
{\( -\frac53 \)}{\( \frac59 \)}{\( -\frac59 \)}{\( \frac53 \)}{\( -5 \)}{}}
\LANGes{
\Question{
Entre las raíces de la ecuación \( 9x^2+130x-75=0 \) interpola dos números tales que con las raíces formen cuatro términos consecutivos de una progresión geométrica. ¿Cuál es el menor de los números?}
{\( -\frac53 \)}{\( \frac59 \)}{\( -\frac59 \)}{\( \frac53 \)}{\( -5 \)}{}}
\End

\Data\ID{1003170604}
\Updated{2022-01-16 08:01:13}
\Level{C}
\SubArea{Geometric Sequences}
\LANGen{
\Question{
Find the sum of the first ten natural powers of the number \( 2 \).}
{\( 2\,046 \)}{\( 2\,048 \)}{\( 2\,047 \)}{\( 2\,000 \)}{\( 1\,536 \)}{}}
\LANGpl{
\Question{
Wyznacz sumę pierwszych dziesięciu naturalnych potęg liczby \( 2 \).}
{\( 2\,046 \)}{\( 2\,048 \)}{\( 2\,047 \)}{\( 2\,000 \)}{\( 1\,536 \)}{}}
\LANGcs{
\Question{
Určete součet prvních deseti přirozených mocnin čísla \( 2 \).}
{\( 2\,046 \)}{\( 2\,048 \)}{\( 2\,047 \)}{\( 2\,000 \)}{\( 1\,536 \)}{}}
\LANGsk{
\Question{
Určte súčet prvých desiatich prirodzených mocnín čísla \( 2 \).}
{\( 2\,046 \)}{\( 2\,048 \)}{\( 2\,047 \)}{\( 2\,000 \)}{\( 1\,536 \)}{}}
\LANGes{
\Question{
Averigua la suma de las primeras diez potencias del número \( 2 \).}
{\( 2\,046 \)}{\( 2\,048 \)}{\( 2\,047 \)}{\( 2\,000 \)}{\( 1\,536 \)}{}}
\End

\Data\ID{1003170605}
\Updated{2022-01-16 08:01:45}
\Level{C}
\SubArea{Geometric Sequences}
\LANGen{
\Question{
Suppose we insert number \( 10\,530 \) between \( 2 \) unknown numbers such that they form \( 3 \) consecutive terms of a geometric sequence and their sum is \( 31\,707 \). What is the smaller of the two unknown numbers?}
{\( 9\,477 \)}{\( 10\,500 \)}{\( 9\,832 \)}{\( 10\,034 \)}{\( 5\,265 \)}{}}
\LANGpl{
\Question{
Załóżmy, że wstawimy liczbę \( 10\,530 \) pomiędzy \( 2 \) nieznane liczby, tak, że stanowią \( 3 \) kolejne wyrazy ciągu geometrycznego, a ich suma wynosi  \( 31\,707 \). Jaka jest wartość mniejszej z dwóch nieznanych liczb?}
{\( 9\,477 \)}{\( 10\,500 \)}{\( 9\,832 \)}{\( 10\,034 \)}{\( 5\,265 \)}{}}
\LANGcs{
\Question{
Vložíme-li mezi dvě neznámá čísla číslo \( 10\,530 \), bude spolu s nimi tvořit tři po sobě jdoucí členy geometrické posloupnosti a jejich součet bude \( 31\,707 \). Určete menší ze dvou neznámých čísel.}
{\( 9\,477 \)}{\( 10\,500 \)}{\( 9\,832 \)}{\( 10\,034 \)}{\( 5\,265 \)}{}}
\LANGsk{
\Question{
Ak vložíme medzi dva neznáme čísla číslo \( 10\,530 \), bude spolu s nimi tvoriť tri po sebe idúce členy geometrickej postupnosti a ich súčet bude \( 31\,707 \). Určte menšie z dvoch neznámych čísel.}
{\( 9\,477 \)}{\( 10\,500 \)}{\( 9\,832 \)}{\( 10\,034 \)}{\( 5\,265 \)}{}}
\LANGes{
\Question{
Si interpolamos \( 10\,530 \) entre dos números desconocidos, los tres números van a formar tres términos consecutivos de una progresión geométrica y su suma será  \( 31\,707 \). Calcula el menor número.}
{\( 9\,477 \)}{\( 10\,500 \)}{\( 9\,832 \)}{\( 10\,034 \)}{\( 5\,265 \)}{}}
\End

\Data\ID{1003170606}
\Updated{2022-01-16 08:01:52}
\Level{C}
\SubArea{Geometric Sequences}
\LANGen{
\Question{
Which of the following numbers cannot be a term of the geometric sequence given by its second term \( a_2=360 \) and the common ratio \( q = \frac35 \)?}
{\( 250 \)}{\( 216 \)}{\( 600 \)}{\( \frac{1\,944}{25} \)}{\( \frac{648}5 \)}{}}
\LANGpl{
\Question{
Która z podanych liczb nie może być wyrazem ciągu geometrycznego, którego drugi wyraz równy jest \( a_2=360 \), a wspólny współczynnik wynosi \( q = \frac35 \)?}
{\( 250 \)}{\( 216 \)}{\( 600 \)}{\( \frac{1\,944}{25} \)}{\( \frac{648}5 \)}{}}
\LANGcs{
\Question{
Které z následujících čísel nemůže být členem geometrické posloupnosti dané druhým členem \( a_2=360 \) a kvocientem \( q = \frac35 \)?}
{\( 250 \)}{\( 216 \)}{\( 600 \)}{\( \frac{1\,944}{25} \)}{\( \frac{648}5 \)}{}}
\LANGsk{
\Question{
Ktoré z nasledujúcich čísel nemôže byť členom geometrickej postupnosti danej druhým členom \( a_2=360 \) a kvocientom \( q = \frac35 \)?}
{\( 250 \)}{\( 216 \)}{\( 600 \)}{\( \frac{1\,944}{25} \)}{\( \frac{648}5 \)}{}}
\LANGes{
\Question{
¿Cuál de los números no puede ser  término de una progresión geométrica si sabemos que su segundo término es \( a_2=360 \) y la razón \( q = \frac35 \)?}
{\( 250 \)}{\( 216 \)}{\( 600 \)}{\( \frac{1\,944}{25} \)}{\( \frac{648}5 \)}{}}
\End

\Data\ID{1103170607}
\Updated{2022-01-16 08:01:46}
\Level{C}
\SubArea{Geometric Sequences}
\IMG{1103170607.pdf}
\LANGen{
\Question{
Let there be an equilateral triangle with a side length of \( 16\,\mathrm{cm} \). Connecting the centres of its sides with line segments we form another equilateral triangle. If we add three more triangles using the same procedure, what is the sum of their perimeters?}
{\( 93\,\mathrm{cm} \)}{\( 72\,\mathrm{cm} \)}{\( 144\,\mathrm{cm} \)}{\( 31\,\mathrm{cm} \)}{\( 90\,\mathrm{cm} \)}{}}
\LANGpl{
\Question{
Dany jest trójkąt równoboczny o boku długości \( 16\,\mathrm{cm} \). Łącząc środki boków  odcinkami tworzymy kolejny trójkąt równoboczny. Jeśli dodamy, trzy kolejne trójkąty w ten sam sposób jaka jest suma ich obwód?}
{\( 93\,\mathrm{cm} \)}{\( 72\,\mathrm{cm} \)}{\( 144\,\mathrm{cm} \)}{\( 31\,\mathrm{cm} \)}{\( 90\,\mathrm{cm} \)}{}}
\LANGcs{
\Question{
Je dán rovnostranný trojúhelník se stranou délky \( 16\,\mathrm{cm} \). Spojnice středů jeho stran tvoří opět rovnostranný trojúhelník. Takto postupně vepíšeme ještě další tři trojúhelníky. Jaký je součet jejich obvodů?}
{\( 93\,\mathrm{cm} \)}{\( 72\,\mathrm{cm} \)}{\( 144\,\mathrm{cm} \)}{\( 31\,\mathrm{cm} \)}{\( 90\,\mathrm{cm} \)}{}}
\LANGsk{
\Question{
Je daný rovnostranný trojuholník so stranou dĺžky \( 16\,\mathrm{cm} \). Spojnice stredov jeho strán tvoria opäť rovnostranný trojuholník. Takto postupne vpíšeme ešte ďalšie tri trojuholníky. Aký je súčet ich obvodov?}
{\( 93\,\mathrm{cm} \)}{\( 72\,\mathrm{cm} \)}{\( 144\,\mathrm{cm} \)}{\( 31\,\mathrm{cm} \)}{\( 90\,\mathrm{cm} \)}{}}
\LANGes{
\Question{
Sea un triángulo equilátero cuyo lado mide \( 16\,\mathrm{cm} \). Conectando los centros de sus lados hacemos otro triángulo equilátero. De la misma manera vamos a construir tres triángulos más. ¿Cuál es la suma de su perímetros?}
{\( 93\,\mathrm{cm} \)}{\( 72\,\mathrm{cm} \)}{\( 144\,\mathrm{cm} \)}{\( 31\,\mathrm{cm} \)}{\( 90\,\mathrm{cm} \)}{}}
\End

\Data\ID{1103170608}
\Updated{2022-01-16 08:01:54}
\Level{C}
\SubArea{Geometric Sequences}
\IMG{1103170608.pdf}
\LANGen{
\Question{
Let there be an equilateral triangle with a side length of \( 16\,\mathrm{cm} \). Connecting the centres of its sides with line segments we form another equilateral triangle. If we add two more triangles using the same procedure, what is the sum of their areas?}
{\( 85\sqrt3\,\mathrm{cm}^2 \)}{\( 128\sqrt3\,\mathrm{cm}^2 \)}{\( \frac{341}4\sqrt3\,\mathrm{cm}^2 \)}{\( 90\,\mathrm{cm}^2 \)}{\( 148\sqrt3\,\mathrm{cm}^2 \)}{}}
\LANGpl{
\Question{
Dany jest trójkąt równoboczny o boku długości \( 16\,\mathrm{cm} \). Łącząc środki boków za pomocą odcinków tworzymy kolejne trójkąty równoboczne. Jeśli dodamy dwa kolejne trójkąty za pomocą tego samego sposobu jaka będzie ich suma ich powierzchni?}
{\( 85\sqrt3\,\mathrm{cm}^2 \)}{\( 128\sqrt3\,\mathrm{cm}^2 \)}{\( \frac{341}4\sqrt3\,\mathrm{cm}^2 \)}{\( 90\,\mathrm{cm}^2 \)}{\( 148\sqrt3\,\mathrm{cm}^2 \)}{}}
\LANGcs{
\Question{
Je dán rovnostranný trojúhelník se stranou délky \( 16\,\mathrm{cm} \). Spojnice středů jeho stran tvoří opět rovnostranný trojúhelník. Takto postupně vepíšeme ještě další dva trojúhelníky. Jaký je součet jejich obsahů?}
{\( 85\sqrt3\,\mathrm{cm}^2 \)}{\( 128\sqrt3\,\mathrm{cm}^2 \)}{\( \frac{341}4\sqrt3\,\mathrm{cm}^2 \)}{\( 90\,\mathrm{cm}^2 \)}{\( 148\sqrt3\,\mathrm{cm}^2 \)}{}}
\LANGsk{
\Question{
Je daný rovnostranný trojuholník so stranou dĺžky \( 16\,\mathrm{cm} \). Spojnice stredov jeho strán tvorí opäť rovnostranný trojuholník. Takto postupne vpíšeme ešte ďalšie dva trojuholníky. Aký je súčet ich obsahov?}
{\( 85\sqrt3\,\mathrm{cm}^2 \)}{\( 128\sqrt3\,\mathrm{cm}^2 \)}{\( \frac{341}4\sqrt3\,\mathrm{cm}^2 \)}{\( 90\,\mathrm{cm}^2 \)}{\( 148\sqrt3\,\mathrm{cm}^2 \)}{}}
\LANGes{
\Question{
Sea un triángulo equilátero cuyo lado mide \( 16\,\mathrm{cm} \). Conectando los centros de sus lados hacemos otro triángulo equilátero. De la misma manera formamos dos triángulos más. ¿Cuál es la suma de su superficies?}
{\( 85\sqrt3\,\mathrm{cm}^2 \)}{\( 128\sqrt3\,\mathrm{cm}^2 \)}{\( \frac{341}4\sqrt3\,\mathrm{cm}^2 \)}{\( 90\,\mathrm{cm}^2 \)}{\( 148\sqrt3\,\mathrm{cm}^2 \)}{}}
\End

\Data\ID{2010005507}
\Updated{2022-08-25 10:08:59}
\Level{C}
\SubArea{Geometric Sequences}
\LANGen{
\Question{
A motorbike depreciates (loses value) at the rate of \(12\, \%\) 
per year. After how many full years will the value of the motorbike decrease to less than one third of its initial value?}
{\(9\)}{\(6\)}{\(7\)}{\(8\)}{\(10\)}{}}
\LANGpl{
\Question{
Motocykl amortyzuje się (traci wartość) w tempie \(12\, \%\)
na rok. Po ilu pełnych latach wartość motocykla spadnie do mniej niż jednej trzeciej jego wartości początkowej?}
{\(9\)}{\(6\)}{\(7\)}{\(8\)}{\(10\)}{}}
\LANGcs{
\Question{
Motorka ztrácí na hodnotě \(12\, \%\) ročně. Za kolik (celých) let klesne cena motorky na méně než jednu třetinu své původní hodnoty?}
{\(9\)}{\(6\)}{\(7\)}{\(8\)}{\(10\)}{}}
\LANGsk{
\Question{
Za koľko rokov klesne hodnota automobilu na menej než tretinu pôvodnej hodnoty, ak ročne stráca automobil \(12\, \%\)  svojej aktuálnej hodnoty?}
{\(9\)}{\(6\)}{\(7\)}{\(8\)}{\(10\)}{}}
\LANGes{
\Question{
Una moto deprecia (pierde valor) por \(12\, \%\) cada año. ¿Cuántos años tardará en tener un tercio del valor inicial?}
{\(9\)}{\(6\)}{\(7\)}{\(8\)}{\(10\)}{}}
\End

\Data\ID{2010005508}
\Updated{2022-08-25 10:08:59}
\Level{C}
\SubArea{Geometric Sequences}
\LANGen{
\Question{
The half-life of Copper-60 is approximately \( 24 \) minutes. How long will it take before its weight decreases from \( 1\,024\,\mathrm{g} \) to \( 8\,\mathrm{g} \)?}
{\( 2 \) hours \( 48 \) minutes}{\( 3 \) hours \( 9 \) minutes}{\( 3 \) hours \( 30 \) minutes}{\( 2 \) hours \( 27 \) minutes}{\( 3 \) hours \( 51 \) minutes}{}}
\LANGpl{
\Question{
Okres półtrwania Copper-60 wynosi około \( 24 \) minut. Ile czasu potrzeba, by jego waga spadła z \( 1\,024\,\mathrm{g} \) do \( 8\,\mathrm{g} \)?}
{\( 2 \) godziny \( 48 \) minuty}{\( 3 \) godziny \( 9 \) minuty}{\( 3 \) godziny \( 30 \) minuty}{\( 2 \)  godziny \( 27 \) minuty}{\( 3 \) godziny \( 51 \) minuty}{}}
\LANGcs{
\Question{
Poločas rozpadu izotopu mědi $ ^{\textrm{60}}\textrm{Cu}$ je přibližně \( 24 \) minut. Za jak dlouho klesne hmotnost z \( 1\,024\,\mathrm{g} \) na \( 8\,\mathrm{g} \)?}
{\( 2 \) hodiny \( 48 \) minut}{\( 3 \) hodiny \( 9 \) minut}{\( 3 \) hodiny \( 30 \) minut}{\( 2 \) hodiny \( 27 \) minut}{\( 3 \) hodiny \( 51 \) minut}{}}
\LANGsk{
\Question{
Polčas rozpadu Medi-60 je približne \( 24 \)  minút. Za ako dlho sa jeho hmotnosť zníži z \( 1\,024\,\mathrm{g} \) na \( 8\,\mathrm{g} \)?}
{\( 2 \) hodiny \( 48 \) minút}{\( 3 \) hodiny \( 9 \) minút}{\( 3 \) hodiny \( 30 \) minút}{\( 2 \) hodiny \( 27 \) minút}{\( 3 \) hodiny \( 51 \) minút}{}}
\LANGes{
\Question{
La vida media de Cobre-60 es aproximadamente \( 24 \) minutos. ¿Cuánto tardará en disminuir su masa de  \( 1\,024\,\mathrm{g} \) a \( 8\,\mathrm{g} \)?}
{\( 2 \) horas \( 48 \) minutos}{\( 3 \) horas \( 9 \) minutos}{\( 3 \) horas \( 30 \) minutos}{\( 2 \) horas \( 27 \) minutos}{\( 3 \) horas \( 51 \) minutos}{}}
\End

\Data\ID{2010005509}
\Updated{2022-08-25 10:08:59}
\Level{C}
\SubArea{Geometric Sequences}
\LANGen{
\Question{
How many numbers do we need to insert between the numbers \( 5 \) and \( 640 \) so that the inserted numbers with the given two numbers are consecutive terms of a geometric sequence? The sum of all numbers inserted must be \( 630 \).}
{\( 6 \)}{\( 4 \)}{\( 3 \)}{\( 5 \)}{\( 7 \)}{}}
\LANGpl{
\Question{
Ile liczb musimy wstawić między liczby \( 5 \) i \( 640 \), aby wstawione liczby o podanych dwóch liczbach były kolejnymi wyrazami ciągu geometrycznego? Suma wszystkich wstawionych liczb musi wynosić \( 630 \).}
{\( 6 \)}{\( 4 \)}{\( 3 \)}{\( 5 \)}{\( 7 \)}{}}
\LANGcs{
\Question{
Kolik čísel musíme vložit mezi čísla \( 5 \) a \( 640 \), aby spolu s danými čísly tvořila po sobě jdoucí členy geometrické posloupnosti? Součet vložených čísel musí být \( 630 \).}
{\( 6 \)}{\( 4 \)}{\( 3 \)}{\( 5 \)}{\( 7 \)}{}}
\LANGsk{
\Question{
Koľko čísel musíme vložiť medzi čísla \( 5 \) a \( 640 \), aby spolu s danými číslami tvorili po sebe idúce členy geometrickej postupnosti a súčet vložených čísel bol \( 630 \)?}
{\( 6 \)}{\( 4 \)}{\( 3 \)}{\( 5 \)}{\( 7 \)}{}}
\LANGes{
\Question{
¿Cuántos números tenemos que insertar entre \( 5 \) y \( 640 \) para que los números insertados junto con los números dados formen términos consecutivos de alguna progresión geométrica.? La suma de todos los números insertados tiene que ser \( 630 \).}
{\( 6 \)}{\( 4 \)}{\( 3 \)}{\( 5 \)}{\( 7 \)}{}}
\End

\Data\ID{2110014005}
\Updated{2023-02-23 11:02:23}
\Level{C}
\SubArea{Geometric Sequences}
\LANGen{
\Question{
Between the roots of the quadratic equation \(4x^2-35x+54=0\) insert two numbers, so that the numbers with the roots together form four consecutive terms of a geometric sequence. This part of the geometric sequence is shown in one of the graphs. Choose the graph.}
{}{}{}{}{}{}}
\LANGpl{
\Question{
Między pierwiastki równania kwadratowego \(4x^2-35x+54=0\) wstawiamy dwie liczby tak, aby liczby razem z pierwiastkami tworzyły cztery kolejne wyrazy ciągu geometrycznego. Ta część ciągu geometrycznego jest pokazana na jednym z wykresów. Wybierz wykres.}
{}{}{}{}{}{}}
\LANGcs{
\Question{
Mezi dva kořeny kvadratické rovnice \(4x^2-35x+54=0\) vložte dvě čísla tak, aby všechna čtyři čísla společně tvořila po sobě jdoucí členy geometrické posloupnosti. Vyberte graf, který tuto část posloupnosti zobrazuje.}
{}{}{}{}{}{}}
\LANGsk{
\Question{
Medzi korene kvadratickej rovnice \(4x^2-35x+54=0\) vložte dve čísla tak, aby všetky spolu tvorili štyri po sebe idúce členy geometrickej postupnosti. Na jednom z obrázkov je daná časť grafu tejto postupnosti. Vyberte, ktorý to je.}
{}{}{}{}{}{}}
\LANGes{
\Question{
Entre las raíces de la ecuación cuadrática \(4x^2-35x+54=0\) inserta dos números tales que dichos números junto con las raíces formen cuatro términos consecutivos de una progresión geométrica. Esta parte de la progresión geométrica se muestra en una de las gráficas. Elige la gráfica.}
{}{}{}{}{}{}}

\IMGanswer{2010014001-figure8_0.pdf}
\IMGanswer{2010014001-figure9_0.pdf}
\IMGanswer{2010014001-figure10_0.pdf}
\IMGanswer{2010014001-figure11_0.pdf}\End

\Data\ID{2110014006}
\Updated{2023-02-26 10:02:46}
\Level{C}
\SubArea{Geometric Sequences}
\LANGen{
\Question{
Between the roots of the quadratic equation \(9x^2-35x+24=0\) insert two numbers, so that the numbers with the roots together form four consecutive terms of a geometric sequence. This part of the geometric sequence is shown in one of the graphs. Choose the graph.}
{}{}{}{}{}{}}
\LANGpl{
\Question{
Między pierwiastki równania kwadratowego \(9x^2-35x+24=0\) wstawiamy dwie liczby tak, aby liczby razem z pierwiastkami tworzyły cztery kolejne wyrazy ciągu geometrycznego. Ta część ciągu geometrycznego jest pokazana na jednym z wykresów. Wybierz wykres.}
{}{}{}{}{}{}}
\LANGcs{
\Question{
Mezi dva kořeny kvadratické rovnice \(9x^2-35x+24=0\) vložte dvě čísla tak, aby všechna čtyři čísla společně tvořila po sobě jdoucí členy geometrické posloupnosti. Vyberte graf, který tuto část posloupnosti zobrazuje.}
{}{}{}{}{}{}}
\LANGsk{
\Question{
Medzi korene kvadratickej rovnice \(9x^2-35x+24=0\) vložte dve čísla tak, aby všetky spolu tvorili štyri po sebe idúce členy geometrickej postupnosti. Na jednom z obrázkov je daná časť grafu tejto postupnosti. Vyberte, ktorý to je.}
{}{}{}{}{}{}}
\LANGes{
\Question{
Entre las raíces de la ecuación cuadrática \(9x^2-35x+24=0\) inserta dos números tales que dichos números junto con las raíces formen cuatro términos consecutivos de una progresión geométrica. Esta parte de la progresión geométrica se muestra en una de las gráficas. Elige la gráfica.}
{}{}{}{}{}{}}

\IMGanswer{2010014001-figure12_0.pdf}
\IMGanswer{2010014001-figure13_0.pdf}
\IMGanswer{2010014001-figure14_0.pdf}
\IMGanswer{2010014001-figure15.pdf}\End

\Data\ID{9000070503}
\Updated{2025-01-26 12:01:01}
\Level{C}
\SubArea{Geometric Sequences}
\LANGen{
\Question{
Three numbers form a geometric sequence. The sum of these numbers is
\(39\) and the
product is \(1\: 000\).
Find the smallest of these numbers.}
{\(4\)}{\(2\)}{\(2.5\)}{\(10\)}{\(25\)}{}}
\LANGpl{
\Question{
Trzy liczby tworzą ciąg geometryczny. Suma tych wyrazów wynosi 
\(39\), a iloczyn jest równy \(1\: 000\).
Jaka jest najmniejsza z tych liczb?}
{\(4\)}{\(2\)}{\(2.5\)}{\(10\)}{\(25\)}{}}
\LANGcs{
\Question{
Tři čísla, která tvoří geometrickou posloupnost, mají součet
\(39\) a
součin \(1\: 000\).
Nejmenší z těchto čísel je:}
{\(4\)}{\(2\)}{\(2{,}5\)}{\(10\)}{\(25\)}{}}
\LANGsk{
\Question{
Tri čísla, ktoré tvoria geometrickú postupnosť, majú súčet
\(39\) a
súčin \(1\: 000\).
Najmenšie z týchto čísel je:}
{\(4\)}{\(2\)}{\(2{,}5\)}{\(10\)}{\(25\)}{}}
\LANGes{
\Question{
La suma de los tres números que forman una progresión geométrica es
\(39\) y su producto es \(1\: 000\).
El menor de los números es:}
{\(4\)}{\(2\)}{\(2.5\)}{\(10\)}{\(25\)}{}}
\End

\Data\ID{9000070504}
\Updated{2022-01-16 08:01:34}
\Level{C}
\SubArea{Geometric Sequences}
\LANGen{
\Question{
The geometric sequence is formed by consecutive powers of the number
\(3\). The eighth
term is \(a_{8} = 3^{10}\).
Find the sum of the first five terms.}
{\(3\: 267\)}{\(1\: 089\)}{\(2\: 178\)}{\(4\: 356\)}{\(6\: 543\)}{}}
\LANGpl{
\Question{
Ciąg geometryczny tworzą kolejne potęgi liczby 
\(3\). Ósmy wyraz ciągu wynosi
 \(a_{8} = 3^{10}\).
Oblicz sumę pięciu początkowych wyrazów tego ciągu.}
{\(3\: 267\)}{\(1\: 089\)}{\(2\: 178\)}{\(4\: 356\)}{\(6\: 543\)}{}}
\LANGcs{
\Question{
V posloupnosti, která je tvořena po sobě jdoucími mocninami čísla
\(3\), platí
\(a_{8} = 3^{10}\). Součet
prvních \(5\) 
členů je:}
{\(3\: 267\)}{\(1\: 089\)}{\(2\: 178\)}{\(4\: 356\)}{\(6\: 543\)}{}}
\LANGsk{
\Question{
V postupnosti, ktorá je tvorená po sebe idúcimi mocninami čísla
\(3\), platí
\(a_{8} = 3^{10}\). Súčet
prvých \(5\) 
členov je:}
{\(3\: 267\)}{\(1\: 089\)}{\(2\: 178\)}{\(4\: 356\)}{\(6\: 543\)}{}}
\LANGes{
\Question{
En una progresión geométrica formada por potencias de 
\(3\) consecutivas, sabemos que
\(a_{8} = 3^{10}\). La suma de sus primeros  \(5\) términos es:}
{\(3\: 267\)}{\(1\: 089\)}{\(2\: 178\)}{\(4\: 356\)}{\(6\: 543\)}{}}
\End

\Data\ID{9000070505}
\Updated{2025-01-26 12:01:01}
\Level{C}
\SubArea{Geometric Sequences}
\LANGen{
\Question{
The dimensions of the box form a geometric sequence. The volume of the box is
\(27\, \mathrm{cm}^{3}\) and the length of
the shortest side is \(2\, \mathrm{cm}\).
Find the surface area of the box.}
{\(57\, \mathrm{cm}^{2}\)}{\(28.5\, \mathrm{cm}^{2}\)}{\(27\, \mathrm{cm}^{2}\)}{\(35\, \mathrm{cm}^{2}\)}{\(45\, \mathrm{cm}^{2}\)}{}}
\LANGpl{
\Question{
Wymiary pudełka tworzą ciąg geometryczny. Objętość pudełka wynosi
\(27\, \mathrm{cm}^{3}\),  a długość najkrótszego 
boku jest równa  \(2\, \mathrm{cm}\).
Oblicz pole powierzchni tego pudełka.}
{\(57\, \mathrm{cm}^{2}\)}{\(28.5\, \mathrm{cm}^{2}\)}{\(27\, \mathrm{cm}^{2}\)}{\(35\, \mathrm{cm}^{2}\)}{\(45\, \mathrm{cm}^{2}\)}{}}
\LANGcs{
\Question{
Délky hran kvádru tvoří geometrickou posloupnost. Objem kvádru je
\(27\, \mathrm{cm}^{3}\). Jeho nejkratší
hrana měří \(2\, \mathrm{cm}\).
Jeho povrch je:}
{\(57\, \mathrm{cm}^{2}\)}{\(28{,}5\, \mathrm{cm}^{2}\)}{\(27\, \mathrm{cm}^{2}\)}{\(35\, \mathrm{cm}^{2}\)}{\(45\, \mathrm{cm}^{2}\)}{}}
\LANGsk{
\Question{
Dĺžky hrán kvádra tvoria geometrickú postupnosť. Objem kvádra je
\(27\, \mathrm{cm}^{3}\). Jeho najkratšia
hrana meria \(2\, \mathrm{cm}\).
Jeho povrch je:}
{\(57\, \mathrm{cm}^{2}\)}{\(28{,}5\, \mathrm{cm}^{2}\)}{\(27\, \mathrm{cm}^{2}\)}{\(35\, \mathrm{cm}^{2}\)}{\(45\, \mathrm{cm}^{2}\)}{}}
\LANGes{
\Question{
Las dimensiones de un prisma rectangular forman una progresión geométrica. Su volumen es
\(27\, \mathrm{cm}^{3}\) y su lado más corto mide \(2\, \mathrm{cm}\).
Su superficie es:}
{\(57\, \mathrm{cm}^{2}\)}{\(28.5\, \mathrm{cm}^{2}\)}{\(27\, \mathrm{cm}^{2}\)}{\(35\, \mathrm{cm}^{2}\)}{\(45\, \mathrm{cm}^{2}\)}{}}
\End

\Data\ID{9000070506}
\Updated{2025-01-26 12:01:01}
\Level{C}
\SubArea{Geometric Sequences}
\LANGen{
\Question{
The intensity of the light is reduced by
\(8\, \%\) when the
light shines through the glass panel. How much of the light intensity remains if the light shines
through \(6\) 
glass panels?}
{\(60.6\, \%\)}{\(39.4\, \%\)}{\(52\, \%\)}{\(48\, \%\)}{\(3.14\, \%\)}{}}
\LANGpl{
\Question{
Intensywność światła zmniejsza się o 
\(8\, \%\), gdy świeci ono przez szklany panel. Ile intensywności światła zostaje, gdy świeci ono przez
 \(6\) szklanych paneli.}
{\(60{,}6\, \%\)}{\(39{,}4\, \%\)}{\(52\, \%\)}{\(48\, \%\)}{\(3{,}14\, \%\)}{}}
\LANGcs{
\Question{
Při průchodu skleněnou deskou ztrácí světlo
\(8\, \%\) své
intenzity. Kolik procent původní intenzity světla zůstane po průchodu
\(6\) 
takovými deskami?}
{\(60{,}6\, \%\)}{\(39{,}4\, \%\)}{\(52\, \%\)}{\(48\, \%\)}{\(3{,}14\, \%\)}{}}
\LANGsk{
\Question{
Pri prechode sklenenou doskou stráca svetlo
\(8\, \%\) svoju
intenzitu. Koľko percent pôvodnej intenzity svetla zostane po prechode
\(6\) 
takými doskami?}
{\(60{,}6\, \%\)}{\(39{,}4\, \%\)}{\(52\, \%\)}{\(48\, \%\)}{\(3{,}14\, \%\)}{}}
\LANGes{
\Question{
La luz pierde \(8\, \%\) de su intensidad pasando por una lámina de vidrio. ¿Qué porcentaje de su intensidad inicial tendrá la luz después de pasar por 
\(6\) láminas?}
{\(60.6\, \%\)}{\(39.4\, \%\)}{\(52\, \%\)}{\(48\, \%\)}{\(3.14\, \%\)}{}}
\End

\Data\ID{9000070507}
\Updated{2022-04-23 05:04:20}
\Level{C}
\SubArea{Geometric Sequences}
\LANGen{
\Question{
A car depreciates (loses value) at the rate of \(15\, \%\) 
per year. After how many full years will the value of the car decrease to less than one quarter of its initial value?}
{\(9\)}{\(6\)}{\(7\)}{\(8\)}{\(10\)}{}}
\LANGpl{
\Question{
Wartość samochodu zmniejsza się o \(15\, \%\) 
rocznie. (Cena samochodu po roku jest mniejsza od aktualnej ceny o 
\(15\, \%\).) Po ilu latach samochód osiągnie mniej niż jedną czwartą swojej wartości?}
{\(9\)}{\(6\)}{\(7\)}{\(8\)}{\(10\)}{}}
\LANGcs{
\Question{
Za kolik let klesne hodnota automobilu na méně než čtvrtinu
původní hodnoty, jestliže ročně ztrácí automobil
\(15\, \%\) své
aktuální hodnoty?}
{\(9\)}{\(6\)}{\(7\)}{\(8\)}{\(10\)}{}}
\LANGsk{
\Question{
Za koľko rokov klesne hodnota automobilu na menej než štvrtinu
pôvodnej hodnoty, ak ročne stráca automobil
\(15\, \%\) svojej
aktuálnej hodnoty?}
{\(9\)}{\(6\)}{\(7\)}{\(8\)}{\(10\)}{}}
\LANGes{
\Question{
Un coche pierde su valor un \(15\, \%\)  cada año. ¿En cuántos años su valor será menor que un cuarto de su valor inicial?}
{\(9\)}{\(6\)}{\(7\)}{\(8\)}{\(10\)}{}}
\End

\Data\ID{9000070508}
\Updated{2022-01-13 08:01:30}
\Level{C}
\SubArea{Geometric Sequences}
\LANGen{
\Question{
The sum of the first three terms of a geometric sequence is
\(27\). The sum of the
next three terms is \(512\).
Find the quotient.}
{\(\frac{8}
{3}\)}{\(2\)}{\(3\)}{\(\frac{3}
{2}\)}{\(\frac{4}
{3}\)}{}}
\LANGpl{
\Question{
Suma trzech początkowych wyrazów ciągu geometrycznego wynosi 
\(27\). Suma trzech kolejnych wyrazów jest równa \(512\).
Oblicz iloraz.}
{\(\frac{8}
{3}\)}{\(2\)}{\(3\)}{\(\frac{3}
{2}\)}{\(\frac{4}
{3}\)}{}}
\LANGcs{
\Question{
Součet prvních \(3\) členů
geometrické posloupnosti je \(27\).
Součet následujících tří členů je
\(512\).
Kvocient této posloupnosti je roven:}
{\(\frac{8}
{3}\)}{\(2\)}{\(3\)}{\(\frac{3}
{2}\)}{\(\frac{4}
{3}\)}{}}
\LANGsk{
\Question{
Súčet prvých \(3\) členov
geometrickej postupnosti je \(27\).
Súčet následujúcich troch členov je
\(512\).
Kvocient tejto postupnosti je rovný:}
{\(\frac{8}
{3}\)}{\(2\)}{\(3\)}{\(\frac{3}
{2}\)}{\(\frac{4}
{3}\)}{}}
\LANGes{
\Question{
La suma de los primeros \(3\) términos de una progresión geométrica es  \(27\).
La suma de los siguientes tres números es
\(512\).
La razón de esta progresión es:}
{\(\frac{8}
{3}\)}{\(2\)}{\(3\)}{\(\frac{3}
{2}\)}{\(\frac{4}
{3}\)}{}}
\End
